\newMonth{Februarius}
\setrunningtitles{Februarius}{February}

% ---- martyrology/mart02/mart0201.htm
\needspace{10\baselineskip}
\begin{paracol}{2}
\selectlanguage{latin}
\begin{center}{\color{gregoriocolor} Kaléndis Februárii. 
 Luna\dots\ }\end{center}
\switchcolumn
\selectlanguage{english}
\begin{center}{\color{gregoriocolor} The 1\textsuperscript{st} Day of 
 February. The\dots\ Day of the Moon.}\end{center}
\end{paracol}

\noindent\begin{tabularx}{\linewidth}{*{19}{>{\centering\arraybackslash}X}}
 \textcolor{gregoriocolor}{a} & \textcolor{gregoriocolor}{b} & \textcolor{gregoriocolor}{c} & \textcolor{gregoriocolor}{d} & \textcolor{gregoriocolor}{e} & \textcolor{gregoriocolor}{f} & \textcolor{gregoriocolor}{g} & \textcolor{gregoriocolor}{h} & \textcolor{gregoriocolor}{i} & \textcolor{gregoriocolor}{k} & \textcolor{gregoriocolor}{l} & \textcolor{gregoriocolor}{m} & \textcolor{gregoriocolor}{n} & \textcolor{gregoriocolor}{p} & \textcolor{gregoriocolor}{q} & \textcolor{gregoriocolor}{r} & \textcolor{gregoriocolor}{s} & \textcolor{gregoriocolor}{t} & \textcolor{gregoriocolor}{u} \\
 3 & 4 & 5 & 6 & 7 & 8 & 9 & 10 & 11 & 12 & 13 & 14 & 15 & 16 & 17 & 18 & 19 & 20 & 21 \\
\end{tabularx}
\vspace{0.5\baselineskip}
\noindent\begin{tabularx}{\linewidth}{*{12}{>{\centering\arraybackslash}X}}
 \textcolor{gregoriocolor}{A} & \textcolor{gregoriocolor}{B} & \textcolor{gregoriocolor}{C} & \textcolor{gregoriocolor}{D} & \textcolor{gregoriocolor}{E} & F & \textcolor{gregoriocolor}{F} & \textcolor{gregoriocolor}{G} & \textcolor{gregoriocolor}{H} & \textcolor{gregoriocolor}{M} & \textcolor{gregoriocolor}{N} & \textcolor{gregoriocolor}{P} \\
 22 & 23 & 24 & 25 & 26 & 27 & 27 & 28 & 29 & 30 & 1 & 2 \\
\end{tabularx}

\begin{paracol}{2}
\selectlanguage{latin}
\lettrine[lines=2]{S}{ancti} Ignátii, Epíscopi Antiochéni et Mártyris, qui glorióse martyrium 
 consummávit tertiodécimo Kaléndas Januárii.
\switchcolumn
\selectlanguage{english}
\lettrine[lines=2]{S}{t.} Ignatius, bishop of Antioch and martyr, who gloriously suffered 
 martyrdom on the 20th of December.
\switchcolumn*
\selectlanguage{latin}
Smyrnæ sancti Piónii, Presbyteri et Mártyris; 
 qui, post apologías pro fide Christiána conscríptas, post squalórem cárceris, 
 ubi multos fratrum ad martyrii tolerántiam suis exhortatiónibus roborávit, 
 tandem, cruciátibus multis vexátus, clavis confíxus et ardénti rogo 
 superpósitus, beátum pro Christo finem sortítus est. Cum ipso autem at 
 álii quíndecim passi sunt.
\switchcolumn
\selectlanguage{english}
At Smyrna, St. Pionius, priest and martyr, who, after writing apologies for 
 the Catholic faith, and after suffering imprisonment in a loathsome dungeon, 
 where by his exhortations he encouraged many of his brethren even to 
 martyrdom, and after enduring excruciating pains from being pierced with 
 nails and laid on a hot fire, ended happily his life for Christ. With 
 him suffered fifteen others.
\switchcolumn*
\selectlanguage{latin}
Ravénnæ sancti Sevéri Epíscopi, qui, ob 
 præclára mérita, signo colúmbæ fuit eléctus.
\switchcolumn
\selectlanguage{english}
At Ravenna, the holy bishop Severus, whose great virtues deserved that he 
 should be raised to the episcopate, which action was confirmed with the sign 
 of a dove.
\switchcolumn*
\selectlanguage{latin}
In civitáte Tricastína, in Gállia, sancti Pauli Epíscopi, cujus vita 
 virtútibus cláruit, et mors pretiósa miráculis commendátur.
\switchcolumn
\selectlanguage{english}
At Trois-Châteaux in France, St. Paul, bishop, whose life was eminent for 
 virtues, and whose death was made precious by miracles.
\switchcolumn*
\selectlanguage{latin}
Apud Kildáriam, in Hibérnia, sanctæ Brígidæ 
 Vírginis, quæ, cum lignum altáris tetigísset in testimónium virginitátis suæ, 
 lignum ipsum statim víride factum est.
\switchcolumn
\selectlanguage{english}
At Kildare in Ireland, St. Bridget, virgin. Once, when she touched the 
 wood of an altar, it immediately sprouted into life, in testimony of her 
 virginity.
\switchcolumn*
\selectlanguage{latin}
Apud Castrum Florentínum, in Etrúria, beátæ 
 Viridiánæ, Vírginis reclúsæ, ex Ordine Vallis Umbrósæ.
\switchcolumn
\selectlanguage{english}
At Castel-Fiorentino in Tuscany, the blessed virgin Veridiana, a recluse of 
 the Order Vallombrosa.
\switchcolumn*
\selectlanguage{latin}
\end{paracol}

% \continueMonth{Februarius}

% ---- martyrology/mart02/mart0202.htm
\needspace{10\baselineskip}
\begin{paracol}{2}
\selectlanguage{latin}
\begin{center}{\color{gregoriocolor} Quarto Nonas Februárii. 
 Luna\dots\ }\end{center}
\switchcolumn
\selectlanguage{english}
\begin{center}{\color{gregoriocolor} The 2\textsuperscript{nd} Day of 
 February. The\dots\ Day of the Moon.}\end{center}
\end{paracol}

\noindent\begin{tabularx}{\linewidth}{*{19}{>{\centering\arraybackslash}X}}
 \textcolor{gregoriocolor}{a} & \textcolor{gregoriocolor}{b} & \textcolor{gregoriocolor}{c} & \textcolor{gregoriocolor}{d} & \textcolor{gregoriocolor}{e} & \textcolor{gregoriocolor}{f} & \textcolor{gregoriocolor}{g} & \textcolor{gregoriocolor}{h} & \textcolor{gregoriocolor}{i} & \textcolor{gregoriocolor}{k} & \textcolor{gregoriocolor}{l} & \textcolor{gregoriocolor}{m} & \textcolor{gregoriocolor}{n} & \textcolor{gregoriocolor}{p} & \textcolor{gregoriocolor}{q} & \textcolor{gregoriocolor}{r} & \textcolor{gregoriocolor}{s} & \textcolor{gregoriocolor}{t} & \textcolor{gregoriocolor}{u} \\
 4 & 5 & 6 & 7 & 8 & 9 & 10 & 11 & 12 & 13 & 14 & 15 & 16 & 17 & 18 & 19 & 20 & 21 & 22 \\
\end{tabularx}
\vspace{0.5\baselineskip}
\noindent\begin{tabularx}{\linewidth}{*{12}{>{\centering\arraybackslash}X}}
 \textcolor{gregoriocolor}{A} & \textcolor{gregoriocolor}{B} & \textcolor{gregoriocolor}{C} & \textcolor{gregoriocolor}{D} & \textcolor{gregoriocolor}{E} & F & \textcolor{gregoriocolor}{F} & \textcolor{gregoriocolor}{G} & \textcolor{gregoriocolor}{H} & \textcolor{gregoriocolor}{M} & \textcolor{gregoriocolor}{N} & \textcolor{gregoriocolor}{P} \\
 23 & 24 & 25 & 26 & 27 & 28 & 28 & 29 & 30 & 1 & 2 & 3 \\
\end{tabularx}

\begin{paracol}{2}
\selectlanguage{latin}
\lettrine[lines=2]{P}{urificátio} beátæ Maríæ Vírginis, quæ a Græcis 
 Hypapánte Dómini appellátur.
\switchcolumn
\selectlanguage{english}
\lettrine[lines=2]{T}{he} Purification of the Blessed Virgin Mary, called by the Greeks the 
 Hypapante (meeting) of the Lord.
\switchcolumn*
\selectlanguage{latin}
Cæsaréæ, in Palæstína, sancti Cornélii 
 Centuriónis, quem beátus Petrus Apóstolus baptizávit, et apud præfátam urbem 
 Episcopáli sublimávit honóre.
\switchcolumn
\selectlanguage{english}
At Caesarea in Palestine, St. Cornelius, a centurion, whom the blessed 
 apostle Peter baptized, and raised to the episcopal dignity in that city.
\switchcolumn*
\selectlanguage{latin}
Romæ, via Salária, pássio sancti Aproniáni 
 Commentariénsis, qui, adhuc Gentílis, cum sanctum Sisínium e cárcere 
 edúceret ut Laodício Præfécto præsentáret, vocémque de cælo factam audíret: 
 « Veníte, benedícti Patris mei, percípite regnum, quod vobis parátum est a 
 constitutióne mundi », credens baptizátur, et póstea, in confessióne Dómini, 
 vitæ finem senténtia capitáli accépit.
\switchcolumn
\selectlanguage{english}
At Rome, on the Salarian Way, the passion of St. Apronian, a notary. 
 While he was yet a heathen, and was leading St. Sisinius out of prison to 
 present him before the governor Laodicius, he head a voice from heaven 
 saying: Come, ye blessed of my Father, possess the kingdom which I have 
 prepared for you from the beginning of the world.`` At once he 
 believed, was baptized, and after confessing our Lord, received sentence of 
 death.
\switchcolumn*
\selectlanguage{latin}
Item Romæ sanctórum Mártyrum Fortunáti, 
 Feliciáni, Firmi et Cándidi.
\switchcolumn
\selectlanguage{english}
Also at Rome, the holy martyrs Fortunatus, Felician, Firmus and Candidus.
\switchcolumn*
\selectlanguage{latin}
Aureliánis, in Gállia, sancti Flósculi Epíscopi.
\switchcolumn
\selectlanguage{english}
At Orleans in France, the holy bishop Flosculus.
\switchcolumn*
\selectlanguage{latin}
Cantuáriæ, in Anglia, natális sancti Lauréntii 
 Epíscopi, qui, post sanctum Augustínum, eam Ecclésiam gubernávit, et Regem 
 ipsum ad fidem convértit.
\switchcolumn
\selectlanguage{english}
At Canterbury in England, the birthday of St. Lawrence, bishop, who 
 succeeded St. Augustine in the government of that church, and converted the 
 king himself to the faith.
\switchcolumn*
\selectlanguage{latin}
Prati, in Etrúria, sanctæ Catharínæ de Rícciis, 
 Vírginis Florentínæ, ex Ordine Prædicatórum, ob cæléstium donórum copiam 
 insígnis; quam Benedíctus Décimus quartus, Póntifex Máximus, sanctárum 
 Vírginum fastis adscrípsit.
\switchcolumn
\selectlanguage{english}
At Prati in Tuscany, St. Catherine de Ricci, a virgin of Florence, member of 
 the Order of Preachers, famous for a plenitude of heavenly gifts. Pope 
 Benedict XIV placed her name on the roll of holy virgins.
\switchcolumn*
\selectlanguage{latin}
Burdígalæ, in Gállia, sanctæ Joánnæ de 
 Léstonnac, Víduæ, Institúti Filiárum beátæ Vírginis Maríæ Fundatrícis, 
 caritátis stúdio ac puellárum instituendárum cura insígnis; quam Pius Papa 
 Duodécimus Sanctárum número accénsuit.
\switchcolumn
\selectlanguage{english}
At Bordeaux in France, St. Joan de Lestonnac, widow, foundress of the 
 Daughters of the blessed Virgin Mary, renowned for the practice of charity 
 and the care of her girl pupils, and whom Pope Pius XII raised to the number 
 of the saints.
\switchcolumn*
\selectlanguage{latin}
\end{paracol}


% ---- martyrology/mart02/mart0203.htm
\needspace{10\baselineskip}
\begin{paracol}{2}
\selectlanguage{latin}
\begin{center}{\color{gregoriocolor} Tértio Nonas Februárii. 
 Luna\dots\ }\end{center}
\switchcolumn
\selectlanguage{english}
\begin{center}{\color{gregoriocolor} The 3\textsuperscript{rd} Day of 
 February. The\dots\ Day of the Moon.}\end{center}
\end{paracol}

\noindent\begin{tabularx}{\linewidth}{*{19}{>{\centering\arraybackslash}X}}
 \textcolor{gregoriocolor}{a} & \textcolor{gregoriocolor}{b} & \textcolor{gregoriocolor}{c} & \textcolor{gregoriocolor}{d} & \textcolor{gregoriocolor}{e} & \textcolor{gregoriocolor}{f} & \textcolor{gregoriocolor}{g} & \textcolor{gregoriocolor}{h} & \textcolor{gregoriocolor}{i} & \textcolor{gregoriocolor}{k} & \textcolor{gregoriocolor}{l} & \textcolor{gregoriocolor}{m} & \textcolor{gregoriocolor}{n} & \textcolor{gregoriocolor}{p} & \textcolor{gregoriocolor}{q} & \textcolor{gregoriocolor}{r} & \textcolor{gregoriocolor}{s} & \textcolor{gregoriocolor}{t} & \textcolor{gregoriocolor}{u} \\
 5 & 6 & 7 & 8 & 9 & 10 & 11 & 12 & 13 & 14 & 15 & 16 & 17 & 18 & 19 & 20 & 21 & 22 & 23 \\
\end{tabularx}
\vspace{0.5\baselineskip}
\noindent\begin{tabularx}{\linewidth}{*{12}{>{\centering\arraybackslash}X}}
 \textcolor{gregoriocolor}{A} & \textcolor{gregoriocolor}{B} & \textcolor{gregoriocolor}{C} & \textcolor{gregoriocolor}{D} & \textcolor{gregoriocolor}{E} & F & \textcolor{gregoriocolor}{F} & \textcolor{gregoriocolor}{G} & \textcolor{gregoriocolor}{H} & \textcolor{gregoriocolor}{M} & \textcolor{gregoriocolor}{N} & \textcolor{gregoriocolor}{P} \\
 24 & 25 & 26 & 27 & 28 & 29 & 29 & 30 & 1 & 2 & 3 & 4 \\
\end{tabularx}

\begin{paracol}{2}
\selectlanguage{latin}
\lettrine[lines=2]{S}{ebáste,} in Arménia, pássio sancti Blásii, Epíscopi et Mártyris; qui, 
 multórum patrátor miraculórum, sub Agricoláo Præside, 
 post diútinam cæsiónem, atque in ligno suspensiónem, ubi férreis pectínibus 
 carnes ejus dirúptæ sunt, post tetérrimum cárcerem et in lacum demersiónem, 
 unde salvus exívit, tandem, jubénte eódem Júdice, una cum duóbus púeris, 
 cápite truncátur. Ante ipsum vero septem mulíeres, quæ guttas 
 sánguinis ex ejúsdem Mártyris córpore defluéntes, dum torquerétur, 
 colligébant, proptérea, deprehénsæ quod essent Christiánæ, omnes, post dira 
 torménta, gládio percússæ sunt.
\switchcolumn
\selectlanguage{english}
\lettrine[lines=2]{A}{t} Sebaste in Armenia, in the time of the governor Agricolaus, the passion 
 of St. Blase, bishop and martyr, who, after working many miracles, was 
 scourged for a long time, suspended from a tree where his flesh was 
 lacerated with iron combs. He was then imprisoned in a dark dungeon, 
 thrown into a lake from which he came out safe, and finally, by order of the 
 judge, he and two boys were beheaded. Before him, seven women who were 
 gathering the drops of his blood during his torture, were recognized as 
 Christians, and after undergoing severe torments, were put to death by the 
 sword.
\switchcolumn*
\selectlanguage{latin}
In Africa sancti Celeríni Diáconi, qui, decem et novem dies custódia 
 cárceris septus, in nervo et ferro variísque pœnis 
 gloriósus fuit Christi Conféssor; et, dum inexpugnábili firmitáte certáminis 
 sui vicit adversárium, vincéndi céteris viam fecit.
\switchcolumn
\selectlanguage{english}
In Africa, St. Celerinus, deacon, who was kept nineteen days in prison 
 burdened with fetters, and who gloriously confessed Christ in the midst of 
 afflictions. By overcoming the enemy with invincible constancy, he 
 shewed to others the road to victory.
\switchcolumn*
\selectlanguage{latin}
Ibídem sanctórum trium Mártyrum, ipsíus Celeríni Diáconi consanguineórum, 
 scílicet Laurentíni pátrui, Ignátii avúnculi, et Celerínæ 
 áviæ, qui ántea martyrio coronáti fúerant; de quorum ómnium gloriósis 
 láudibus exstat beáti Cypriáni epístola.
\switchcolumn
\selectlanguage{english}
In the same place, three holy martyrs who were relatives of the same deacon 
 Celerinus; his father's brother Laurentinus, his mother's brother Ignatius 
 and his grandmother Celerina. They were crowned with martyrdom 
 earlier, and were praised highly in an epistle by blessed Cyprian.
\switchcolumn*
\selectlanguage{latin}
Item in Africa sanctórum Mártyrum Felícis, Symphrónii, Hippólyti et Sociórum.
\switchcolumn
\selectlanguage{english}
Likewise in Africa, the holy martyrs Felix, Symphronius, Hippolytus, and 
 their companions.
\switchcolumn*
\selectlanguage{latin}
In óppido Vapíngo, in Gállia, sanctórum Tigídis 
 et Remédii Episcopórum.
\switchcolumn
\selectlanguage{english}
In the town of Gap in France, the holy bishops Tigides and Remedius.
\switchcolumn*
\selectlanguage{latin}
Lugdúni, in Gállia, sanctórum Lupicíni et Felícis,
 ítidem Episcopórum.
\switchcolumn
\selectlanguage{english}
At Lyons in France, Saints Lupicinus and Felix, also bishops.
\switchcolumn*
\selectlanguage{latin}
Bremæ sancti Anschárii, Hamburgénsis ac póstea 
 Breménsis simul Epíscopi, qui Suécos et Danos ad Christi fidem convértit, et 
 a Gregório Papa Quarto Legátus Apostólicus totíus Septentriónis fuit 
 institútus.
\switchcolumn
\selectlanguage{english}
At Bremen, St. Ansgar, bishop of Hamburg and later of Bremen, who converted 
 the Swedes and the Danes to the faith of Christ. He was appointed 
 Apostolic Delegate of all the North by Pope Gregory IV.
\switchcolumn*
\selectlanguage{latin}
\end{paracol}


% ---- martyrology/mart02/mart0204.htm
\needspace{10\baselineskip}
\begin{paracol}{2}
\selectlanguage{latin}
\begin{center}{\color{gregoriocolor} Prídie Nonas Februárii. 
 Luna\dots\ }\end{center}
\switchcolumn
\selectlanguage{english}
\begin{center}{\color{gregoriocolor} The 4\textsuperscript{th} Day of 
 February. The\dots\ Day of the Moon.}\end{center}
\end{paracol}

\noindent\begin{tabularx}{\linewidth}{*{19}{>{\centering\arraybackslash}X}}
 \textcolor{gregoriocolor}{a} & \textcolor{gregoriocolor}{b} & \textcolor{gregoriocolor}{c} & \textcolor{gregoriocolor}{d} & \textcolor{gregoriocolor}{e} & \textcolor{gregoriocolor}{f} & \textcolor{gregoriocolor}{g} & \textcolor{gregoriocolor}{h} & \textcolor{gregoriocolor}{i} & \textcolor{gregoriocolor}{k} & \textcolor{gregoriocolor}{l} & \textcolor{gregoriocolor}{m} & \textcolor{gregoriocolor}{n} & \textcolor{gregoriocolor}{p} & \textcolor{gregoriocolor}{q} & \textcolor{gregoriocolor}{r} & \textcolor{gregoriocolor}{s} & \textcolor{gregoriocolor}{t} & \textcolor{gregoriocolor}{u} \\
 6 & 7 & 8 & 9 & 10 & 11 & 12 & 13 & 14 & 15 & 16 & 17 & 18 & 19 & 20 & 21 & 22 & 23 & 24 \\
\end{tabularx}
\vspace{0.5\baselineskip}
\noindent\begin{tabularx}{\linewidth}{*{12}{>{\centering\arraybackslash}X}}
 \textcolor{gregoriocolor}{A} & \textcolor{gregoriocolor}{B} & \textcolor{gregoriocolor}{C} & \textcolor{gregoriocolor}{D} & \textcolor{gregoriocolor}{E} & F & \textcolor{gregoriocolor}{F} & \textcolor{gregoriocolor}{G} & \textcolor{gregoriocolor}{H} & \textcolor{gregoriocolor}{M} & \textcolor{gregoriocolor}{N} & \textcolor{gregoriocolor}{P} \\
 25 & 26 & 27 & 28 & 29 & 1 & 30 & 1 & 2 & 3 & 4 & 5 \\
\end{tabularx}

\begin{paracol}{2}
\selectlanguage{latin}
\lettrine[lines=2]{S}{ancti} Andréæ Corsíni, ex Ordine Carmelitárum, 
 Epíscopi Fæsuláni et Confessóris; cujus dies natális ágitur octávo Idus 
 Januárii.
\switchcolumn
\selectlanguage{english}
\lettrine[lines=2]{S}{t.} Andrew Corsini, Carmelite bishop of Fiesole, confessor, whose birthday 
 is the 6th of January.
\switchcolumn*
\selectlanguage{latin}
Romæ sancti Eutychii Mártyris, qui illústre 
 martyrium consummávit, ac sepúltus est in cœmetério Callísti; ejúsque 
 sepúlcrum póstea sanctus Dámasus Papa vérsibus exornávit.
\switchcolumn
\selectlanguage{english}
At Rome, St. Eutychius, who endured a glorious martyrdom and was buried in 
 the cemetery of Callistus. Pope St. Damasus wrote an epitaph in verse 
 for his tomb.
\switchcolumn*
\selectlanguage{latin}
Thumi, in Ægypto, pássio beáti Philéæ, ejúsdem 
 civitátis Epíscopi, et Philorómi, Tribúni mílitum; qui, in persecutióne 
 Diocletiáni, cum a cognátis et amícis suadéri non possunt ut sibi párcerent, 
 ambo, datis cervícibus, palmas a Dómino meruérunt. Cum ipsis innúmera 
 étiam multitúdo fidélium ex eádem urbe, pastóris sui vestígia sequens, 
 martyrio coronáta est.
\switchcolumn
\selectlanguage{english}
At Thumis in Egypt, in the persecution of Diocletian, the passion of blessed 
 Philaeus, bishop of that city, and of Philoromus, military tribune, who 
 rejected the exhortations of their relatives and friends to save themselves, 
 offered themselves to death, and so merited immortal palms from God. 
 With them was crowned with martyrdom a numberless multitude of the faithful 
 of the same place, who followed the example of their pastor.
\switchcolumn*
\selectlanguage{latin}
Foro Semprónii sanctórum Mártyrum Aquilíni, Gémini, Gelásii, Magni et Donáti.
\switchcolumn
\selectlanguage{english}
At Fossombrone, the holy martyrs Aquilinus, Geminus, Gelasius, Magnus, and 
 Donatus.
\switchcolumn*
\selectlanguage{latin}
In regno Maravénsi apud Indos Orientáles, sancti Joánnis de Britto, 
 Sacerdótis e Societáte Jesu, qui cum multos infidéles ad fidem convertísset, 
 glorióso martyrio coronátus est.
\switchcolumn
\selectlanguage{english}
In Marava Kingdom in India, St. John de Britto, priest of the Society of 
 Jesus, who having converted many infidels to the faith, was gloriously 
 crowned with martyrdom.
\switchcolumn*
\selectlanguage{latin}
Trecis, in Gállia, sancti Aventíni, Presbyteri et Confessóris.
\switchcolumn
\selectlanguage{english}
At Troyes in France, St. Aventin, priest and confessor.
\switchcolumn*
\selectlanguage{latin}
Pelúsii, in Ægypto, sancti Isidóri, Presbyteri 
 et Mónachi, méritis et doctrína conspícui.
\switchcolumn
\selectlanguage{english}
At Pelusium in Egypt, St. Isidore, a monk renowned for merit and learning.
\switchcolumn*
\selectlanguage{latin}
Sempringhámiæ, in Anglia, sancti Gilbérti, 
 Presbyteri et Confessóris; qui Ordinis Sempringhamiénsis fuit Institútor.
\switchcolumn
\selectlanguage{english}
At Sempringham in England, St. Gilbert, priest and confessor, who founded a 
 religious order at Sempringham.
\switchcolumn*
\selectlanguage{latin}
In oppido Amatrícis, in Aprútio, deposítio sancti Joséphi a Leoníssa, 
 Sacerdótis ex Ordine Minórum Capuccinórum et Confessóris; quem, ob fídei prædicatiónem 
 a Mahumetánis dira perpéssum, labóribus apostólicis et miráculis clarum, 
 Benedíctus Décimus quartus, Póntifex Máximus, in Sanctórum cánonem rétulit.
\switchcolumn
\selectlanguage{english}
In the town of Amatrice, in the diocese of Rieti, the death of St. Joseph of 
 Leonissa, a Capuchin priest who suffered greatly from the Mohammedans. 
 As he was celebrated for his apostolic labours and miracles, he was placed 
 on the list of holy confessors by the Sovereign Pontiff, Benedict XIV.
\switchcolumn*
\selectlanguage{latin}
Bremæ Commemorátio sancti Rembérti, qui, sancti 
 Anschárii discípulus, in ipsíus locum, hac die, óbitum magístri sui próxime 
 subsequénti, olim Hamburgénsis simul ac Breménsis Epíscopus eléctus est.
\switchcolumn
\selectlanguage{english}
At Bremen, the commemoration of St. Rembert, who was a disciple of St. 
 Ansgar, and on this day took his place as bishop of Hamburg and Bremen, the 
 day after the death of his master.
\switchcolumn*
\selectlanguage{latin}
Bitúrcis, in Aquitánia, sanctæ Joánnæ de Valois, 
 Gálliæ Regínæ, Ordinis sanctíssimæ Annuntiatiónis beátæ Maríæ Vírginis 
 Fundatrícis, pietáte et singulári Crucis participatióne illústris, a Pio 
 Papa Duodécimo Sanctárum fastis adscríptæ.
\switchcolumn
\selectlanguage{english}
At Bourges in Aquitaine, St. Jane de Valois, Queen of France, foundress of 
 the Order of Sisters of the Annunciation of the blessed Virgin Mary, 
 renowned for her piety and singular devotion to the Cross, whom Pope Pius 
 XII added to the catalogue of saints.
\switchcolumn*
\selectlanguage{latin}
\end{paracol}


% ---- martyrology/mart02/mart0205.htm
\needspace{10\baselineskip}
\begin{paracol}{2}
\selectlanguage{latin}
\begin{center}{\color{gregoriocolor} Nonis Februárii. 
 Luna\dots\ }\end{center}
\switchcolumn
\selectlanguage{english}
\begin{center}{\color{gregoriocolor} The 5\textsuperscript{th} Day of 
 February. The\dots\ Day of the Moon.}\end{center}
\end{paracol}

\noindent\begin{tabularx}{\linewidth}{*{19}{>{\centering\arraybackslash}X}}
 \textcolor{gregoriocolor}{a} & \textcolor{gregoriocolor}{b} & \textcolor{gregoriocolor}{c} & \textcolor{gregoriocolor}{d} & \textcolor{gregoriocolor}{e} & \textcolor{gregoriocolor}{f} & \textcolor{gregoriocolor}{g} & \textcolor{gregoriocolor}{h} & \textcolor{gregoriocolor}{i} & \textcolor{gregoriocolor}{k} & \textcolor{gregoriocolor}{l} & \textcolor{gregoriocolor}{m} & \textcolor{gregoriocolor}{n} & \textcolor{gregoriocolor}{p} & \textcolor{gregoriocolor}{q} & \textcolor{gregoriocolor}{r} & \textcolor{gregoriocolor}{s} & \textcolor{gregoriocolor}{t} & \textcolor{gregoriocolor}{u} \\
 7 & 8 & 9 & 10 & 11 & 12 & 13 & 14 & 15 & 16 & 17 & 18 & 19 & 20 & 21 & 22 & 23 & 24 & 25 \\
\end{tabularx}
\vspace{0.5\baselineskip}
\noindent\begin{tabularx}{\linewidth}{*{12}{>{\centering\arraybackslash}X}}
 \textcolor{gregoriocolor}{A} & \textcolor{gregoriocolor}{B} & \textcolor{gregoriocolor}{C} & \textcolor{gregoriocolor}{D} & \textcolor{gregoriocolor}{E} & F & \textcolor{gregoriocolor}{F} & \textcolor{gregoriocolor}{G} & \textcolor{gregoriocolor}{H} & \textcolor{gregoriocolor}{M} & \textcolor{gregoriocolor}{N} & \textcolor{gregoriocolor}{P} \\
 26 & 27 & 28 & 29 & 1 & 2 & 1 & 2 & 3 & 4 & 5 & 6 \\
\end{tabularx}

\begin{paracol}{2}
\selectlanguage{latin}
\lettrine[lines=2]{C}{átanæ,} in Sicília, natális sanctæ Agathæ, 
 Vírginis et Mártyris; quæ, tempóribus Décii Imperatóris, sub Quinctiáno 
 Júdice, post álapas et cárcerem, post equúleum et torsiónes, post mamillárum 
 abscissiónem, post volutatiónem in téstulis et carbónibus, tandem in cárcere, 
 Deum precans, consummáta est.
\switchcolumn
\selectlanguage{english}
\lettrine[lines=2]{A}{t} Catana in Sicily, in the time of Emperor Decius and the judge Quinctian, 
 the birthday of St. Agatha, virgin and martyr. After being buffeted, 
 imprisoned, tortured, racked, dragged over pieces of earthenware and burning 
 coals, and having her breasts cut away, she completed her sacrifice in 
 prison while engaged in prayer.
\switchcolumn*
\selectlanguage{latin}
Nangasáchii, in Japónia, pássio vigínti sex Mártyrum, e quibus tres 
 Sacerdótes atque unus Cléricus et duo Láici ad Ordinem Minórum, tres, et in 
 eis unus quidem Cléricus, ad Societátem Jesu, ac septémdecim ad tértium 
 sancti Francísci Ordinem revocántur. Hi omnes pro cathólica fide, in 
 crucem acti et lanceárum íctibus perfóssi, inter divínas laudes ejusdémque 
 fídei prædicatiónem, glorióse occubuérunt; et a 
 Pio Nono, Pontífice Máximo, Sanctórum fastis adscrípti sunt.
\switchcolumn
\selectlanguage{english}
At Nagasaki in Japan, the passion of twenty-six martyrs. Three 
 priests, one cleric, and two lay brothers were members of the Order of 
 Friars Minor; one cleric was of the Society of Jesus, and seventeen belonged 
 to the Third Order of St. Francis. All of them, placed upon crosses 
 for the Catholic faith, and pierced with lances, gloriously died in praising 
 God and preaching that same faith. Their names were added to the roll 
 of saints by Pope Pius IX.
\switchcolumn*
\selectlanguage{latin}
In Ponto commemorátio plurimórum sanctórum Mártyrum, in persecutióne 
 Maximiáni; quorum álii plumbo liquénti perfúsi, álii acútis arundínibus in 
 únguibus cruciáti, ac multis horréndis vexáti torméntis, iisdémque sæpius 
 iterátis, palmas a Dómino et corónas illústri passióne meruérunt.
\switchcolumn
\selectlanguage{english}
In Pontus, during the persecution of Maximian, the commemoration of many 
 holy martyrs, some of whom had molten lead poured on them, others had sharp 
 reeds thrust under their nails, and were often horribly tormented in many 
 other ways. Thus, by their glorious suffering, they deserved to 
 receive at the hands of God palms of victory and their crowns.
\switchcolumn*
\selectlanguage{latin}
Alexandríæ sancti Isidóri, mílitis et Mártyris; 
 qui, in persecutióne Décii, a Numeriáno, exércitus Duce, ob Christi fidem, 
 cápite cæsus est.
\switchcolumn
\selectlanguage{english}
At Alexandria, during the persecution of Decius, St. Isidore, martyr, who 
 was beheaded for the faith of Christ by Numerian, general of the army.
\switchcolumn*
\selectlanguage{latin}
Viénnæ beáti Avíti, Epíscopi et Confessóris; 
 cujus fide, indústria atque admirábili doctrína ab Ariánæ hæresis 
 infestatióne sunt Gálliæ defénsæ.
\switchcolumn
\selectlanguage{english}
At Vienne, blessed Avitus, bishop and confessor, whose faith, labours, and 
 admirable learning protected France against the ravages of the Arian heresy.
\switchcolumn*
\selectlanguage{latin}
Sabióne, in Rhǽtia secúnda, sancti Ingenuíni 
 Epíscopi, cujus vita miráculis éxstitit gloriósa. Sacrum vero ipsíus 
 corpus Brixinónem póstea translátum fuit, ibíque honorífice asservátum.
\switchcolumn
\selectlanguage{english}
At Sabion in the Tyrol, St. Genuinus, bishop, whose illustrious life 
 abounded in miracles. His revered body was afterwards taken to Brixen 
 where a shrine was erected in his honour.
\switchcolumn*
\selectlanguage{latin}
Brixinóne sancti Albuíni Epíscopi, qui eam in civitátem e Sabióne Cáthedram 
 Episcopálem tránstulit, et ibídem, virtútum signis
 émicans, migrávit ad Dóminum.
\switchcolumn
\selectlanguage{english}
At Brixen, St. Albinus, bishop, who moved the Episcopal See from Sabion to 
 that city, and there, eminent by virtue of his miracles, passed to the Lord.
\switchcolumn*
\selectlanguage{latin}
\end{paracol}


% ---- martyrology/mart02/mart0206.htm
\needspace{10\baselineskip}
\begin{paracol}{2}
\selectlanguage{latin}
\begin{center}{\color{gregoriocolor} Octávo Idus Februárii. 
 Luna\dots\ }\end{center}
\switchcolumn
\selectlanguage{english}
\begin{center}{\color{gregoriocolor} The 6\textsuperscript{th} Day of 
 February. The\dots\ Day of the Moon.}\end{center}
\end{paracol}

\noindent\begin{tabularx}{\linewidth}{*{19}{>{\centering\arraybackslash}X}}
 \textcolor{gregoriocolor}{a} & \textcolor{gregoriocolor}{b} & \textcolor{gregoriocolor}{c} & \textcolor{gregoriocolor}{d} & \textcolor{gregoriocolor}{e} & \textcolor{gregoriocolor}{f} & \textcolor{gregoriocolor}{g} & \textcolor{gregoriocolor}{h} & \textcolor{gregoriocolor}{i} & \textcolor{gregoriocolor}{k} & \textcolor{gregoriocolor}{l} & \textcolor{gregoriocolor}{m} & \textcolor{gregoriocolor}{n} & \textcolor{gregoriocolor}{p} & \textcolor{gregoriocolor}{q} & \textcolor{gregoriocolor}{r} & \textcolor{gregoriocolor}{s} & \textcolor{gregoriocolor}{t} & \textcolor{gregoriocolor}{u} \\
 8 & 9 & 10 & 11 & 12 & 13 & 14 & 15 & 16 & 17 & 18 & 19 & 20 & 21 & 22 & 23 & 24 & 25 & 26 \\
\end{tabularx}
\vspace{0.5\baselineskip}
\noindent\begin{tabularx}{\linewidth}{*{12}{>{\centering\arraybackslash}X}}
 \textcolor{gregoriocolor}{A} & \textcolor{gregoriocolor}{B} & \textcolor{gregoriocolor}{C} & \textcolor{gregoriocolor}{D} & \textcolor{gregoriocolor}{E} & F & \textcolor{gregoriocolor}{F} & \textcolor{gregoriocolor}{G} & \textcolor{gregoriocolor}{H} & \textcolor{gregoriocolor}{M} & \textcolor{gregoriocolor}{N} & \textcolor{gregoriocolor}{P} \\
 27 & 28 & 29 & 1 & 2 & 3 & 2 & 3 & 4 & 5 & 6 & 7 \\
\end{tabularx}

\begin{paracol}{2}
\selectlanguage{latin}
\lettrine[lines=2]{S}{ancti} Titi, Epíscopi Creténsium et Confessóris, cujus dies natális occúrrit 
 prídie Nonas Januárii.
\switchcolumn
\selectlanguage{english}
\lettrine[lines=2]{S}{t.} Titus, confessor and bishop of Crete, whose birthday is on the fourth of 
 January.
\switchcolumn*
\selectlanguage{latin}
Cæsaréæ, in Cappadócia, natális sanctæ Dorothéæ, 
 Vírginis et Mártyris; quæ, sub Saprício, illíus Provínciæ Præside, primum 
 per equúlei extensiónem vexáta, dehinc palmis diutíssime cæsa, capitáli 
 senténtia ad últimum puníta est. In ejus confessióne Theóphilus quidam 
 scholásticus ad Christi fidem convérsus, et mox equúleo acérimme tortus, 
 novíssime gládio cæsus est.
\switchcolumn
\selectlanguage{english}
At Caesarea in Cappadocia, the birthday of St. Dorothy, virgin and martyr, 
 who was stretched on the rack, then scourged for a long time with the boughs 
 of a palm tree, and finally condemned to capital punishment by Sapricius, 
 governor of the province. Her noble confession of Christ converted a 
 lawyer named Theophilus, who also was tortured in a barbarous manner, and 
 finally put to death by the sword.
\switchcolumn*
\selectlanguage{latin}
Eméssæ, in Phœnícia, sancti Silváni Epíscopi, 
 qui cum eídem Ecclésiæ annis quadragínta præfuísset, tandem, sub Maximiáno 
 Imperatóre, una cum duóbus áliis, objéctus feris membratímque discérptus, 
 martyrii palmam accépit.
\switchcolumn
\selectlanguage{english}
At Emessa in Phoenicia, in the time of Emperor Maximian, St. Silvanus, 
 bishop, who, after having governed that church for forty years, was 
 delivered to the beasts with two other Christians, and having his limbs all 
 mangled, received the crown of martyrdom.
\switchcolumn*
\selectlanguage{latin}
Eódem die sanctórum Mártyrum Saturníni, Theóphili et Revocátæ.
\switchcolumn
\selectlanguage{english}
The same day, the holy martyrs Caturninus, Theophilus, and Revocata.
\switchcolumn*
\selectlanguage{latin}
Arvérnis, in Gállia, sancti Antholiáni Mártyris.
\switchcolumn
\selectlanguage{english}
In Auvergne in France, St. Atholian, martyr.
\switchcolumn*
\selectlanguage{latin}
Atrébati, in Gálliis, sancti Vedásti, ejúsdem 
 civitátis Epíscopi, cujus vita et mors plúrimis miráculis éxstitit gloriósa.
\switchcolumn
\selectlanguage{english}
At Arras in France, St. Vedast, bishop of that city. The glory of his 
 life and death is attested by many miracles.
\switchcolumn*
\selectlanguage{latin}
Elnóne, in Gállia, sancti Amándi, Epíscopi Trajecténsis, qui miráculis, cum 
 vivus tum mórtuus, glorióse refúlsit; cujus nómine póstmodum insignítum est 
 óppidum, in quo ille monastérium exstrúxerat et mortálem vitam absólverat.
\switchcolumn
\selectlanguage{english}
At Elnon in France, St. Amand, bishop of Maestricht, who was renowned for 
 his miracles during his life and in death. In the town which was named 
 after him, he lived and died in a monastery that he had built.
\switchcolumn*
\selectlanguage{latin}
Bonóniæ sancti Guaríni Cardinális et Epíscopi 
 Prænestíni, vitæ sanctitáte conspícui.
\switchcolumn
\selectlanguage{english}
At Bologna, St. Guarinus, bishop of Palestrina and cardinal, conspicuous for 
 his holiness of life.
\switchcolumn*
\selectlanguage{latin}
\end{paracol}


% ---- martyrology/mart02/mart0207.htm
\needspace{10\baselineskip}
\begin{paracol}{2}
\selectlanguage{latin}
\begin{center}{\color{gregoriocolor} Séptimo Idus Februárii. 
 Luna\dots\ }\end{center}
\switchcolumn
\selectlanguage{english}
\begin{center}{\color{gregoriocolor} The 7\textsuperscript{th} Day of 
 February. The\dots\ Day of the Moon.}\end{center}
\end{paracol}

\noindent\begin{tabularx}{\linewidth}{*{19}{>{\centering\arraybackslash}X}}
 \textcolor{gregoriocolor}{a} & \textcolor{gregoriocolor}{b} & \textcolor{gregoriocolor}{c} & \textcolor{gregoriocolor}{d} & \textcolor{gregoriocolor}{e} & \textcolor{gregoriocolor}{f} & \textcolor{gregoriocolor}{g} & \textcolor{gregoriocolor}{h} & \textcolor{gregoriocolor}{i} & \textcolor{gregoriocolor}{k} & \textcolor{gregoriocolor}{l} & \textcolor{gregoriocolor}{m} & \textcolor{gregoriocolor}{n} & \textcolor{gregoriocolor}{p} & \textcolor{gregoriocolor}{q} & \textcolor{gregoriocolor}{r} & \textcolor{gregoriocolor}{s} & \textcolor{gregoriocolor}{t} & \textcolor{gregoriocolor}{u} \\
 9 & 10 & 11 & 12 & 13 & 14 & 15 & 16 & 17 & 18 & 19 & 20 & 21 & 22 & 23 & 24 & 25 & 26 & 27 \\
\end{tabularx}
\vspace{0.5\baselineskip}
\noindent\begin{tabularx}{\linewidth}{*{12}{>{\centering\arraybackslash}X}}
 \textcolor{gregoriocolor}{A} & \textcolor{gregoriocolor}{B} & \textcolor{gregoriocolor}{C} & \textcolor{gregoriocolor}{D} & \textcolor{gregoriocolor}{E} & F & \textcolor{gregoriocolor}{F} & \textcolor{gregoriocolor}{G} & \textcolor{gregoriocolor}{H} & \textcolor{gregoriocolor}{M} & \textcolor{gregoriocolor}{N} & \textcolor{gregoriocolor}{P} \\
 28 & 29 & 1 & 2 & 3 & 4 & 3 & 4 & 5 & 6 & 7 & 8 \\
\end{tabularx}

\begin{paracol}{2}
\selectlanguage{latin}
\lettrine[lines=2]{S}{ancti} Romuáldi Abbátis, Monachórum Camaldulénsium Patris, cujus dies 
 natális tertiodécimo Kaléndas Júlii recensétur, sed festívitas hac die, ob 
 Translatiónem córporis ejus, potíssimum celebrátur.
\switchcolumn
\selectlanguage{english}
\lettrine[lines=2]{S}{t.} Romuald, founder of the Camaldolese monks, whose birthday is the 19th of 
 June, but celebrated today because of the transference of his body.
\switchcolumn*
\selectlanguage{latin}
Augústæ, cui nunc Londíni nomen, in Británnia, natális beáti Auguli Epíscopi, qui, ætátis cursu per martyrium expléto, 
 ætérna præmia suscípere méruit.
\switchcolumn
\selectlanguage{english}
At London, England, the birthday of blessed Augulus, bishop, who ended the 
 course of his life by martyrdom, and deserved to receive an eternal 
 recompense.
\switchcolumn*
\selectlanguage{latin}
In Phrygia sancti Adáuci Mártyris, qui, ex Itálico génere clarus, et omni 
 fere dignitátum gradu ab Imperatóribus insignítus, tandem, cum adhuc Quæstóris 
 offício fungerétur, martyrii coróna pro fídei defensióne dignátus est.
\switchcolumn
\selectlanguage{english}
In Phrygia, St. Adaucus, martyr, an Italian of noble birth, who was honoured 
 by the emperors with almost every dignity. While he was still 
 discharging the office of quæstor, he was judged worthy of the crown of 
 martyrdom for his defence of the faith.
\switchcolumn*
\selectlanguage{latin}
Ibídem plurimórum sanctórum Mártyrum, urbis uníus cívium, quorum dux erat 
 idem Adáucus; qui, cum omnes Christiáni essent, et constánter in fídei 
 confessióne persísterent, a Galério Maximiáno Imperatóre sunt igne consúmpti.
\switchcolumn
\selectlanguage{english}
Also, many holy martyrs, citizens of this same city of which Adaucus was 
 mayor. As they were all Christians, and persisted in the confession of 
 the faith, they were burned to death by Emperor Galerius Maximian.
\switchcolumn*
\selectlanguage{latin}
Heracléæ, in Ponto, sancti Theodóri, ductóris 
 mílitum, qui, Licínio imperánte, post multa torménta, truncátus cápite, 
 victor migrávit in cælum.
\switchcolumn
\selectlanguage{english}
At Heraclea, in the reign of Licinius, St. Theodore, a military officer, who 
 was beheaded after undergoing many torments, and went victoriously to 
 heaven.
\switchcolumn*
\selectlanguage{latin}
In Ægypto sancti Móysis, Epíscopi venerábilis, 
 qui primum in erémo vitam solitáriam duxit; deínde, peténte Regína 
 Saracenórum Máuvia, Epíscopus factus, gentem illam ferocíssimam magna ex 
 parte ad fidem convértit, et gloriósus méritis quiévit in pace.
\switchcolumn
\selectlanguage{english}
In Egypt, St. Moses, a venerable bishop, who first led a solitary life in 
 the desert, and afterwards, at the request of Mauvia, queen of the Saracens, 
 converted to the faith the greater part of that barbarous people. 
 Being made a bishop, and rich in merits, he peacefully went to his reward.
\switchcolumn*
\selectlanguage{latin}
Lucæ, in Túscia, deposítio sancti Richárdi, 
 Regis Anglórum, qui pater éxstitit sancti Willebáldi, Eistetténsis Epíscopi, 
 ac sanctæ Walbúrgæ Vírginis.
\switchcolumn
\selectlanguage{english}
At Lucca in Tuscany, the death of St. Richard, king of England. He was 
 the father of St. Willebald, bishop of Eichstadt, and of St. Walburga, 
 virgin.
\switchcolumn*
\selectlanguage{latin}
Bonóniæ sanctæ Juliánæ Víduæ.
\switchcolumn
\selectlanguage{english}
At Bologna, St. Juliana, widow.
\switchcolumn*
\selectlanguage{latin}
\end{paracol}


% ---- martyrology/mart02/mart0208.htm
\needspace{10\baselineskip}
\begin{paracol}{2}
\selectlanguage{latin}
\begin{center}{\color{gregoriocolor} Sexto Idus Februárii. 
 Luna\dots\ }\end{center}
\switchcolumn
\selectlanguage{english}
\begin{center}{\color{gregoriocolor} The 8\textsuperscript{th} Day of 
 February. The\dots\ Day of the Moon.}\end{center}
\end{paracol}

\noindent\begin{tabularx}{\linewidth}{*{19}{>{\centering\arraybackslash}X}}
 \textcolor{gregoriocolor}{a} & \textcolor{gregoriocolor}{b} & \textcolor{gregoriocolor}{c} & \textcolor{gregoriocolor}{d} & \textcolor{gregoriocolor}{e} & \textcolor{gregoriocolor}{f} & \textcolor{gregoriocolor}{g} & \textcolor{gregoriocolor}{h} & \textcolor{gregoriocolor}{i} & \textcolor{gregoriocolor}{k} & \textcolor{gregoriocolor}{l} & \textcolor{gregoriocolor}{m} & \textcolor{gregoriocolor}{n} & \textcolor{gregoriocolor}{p} & \textcolor{gregoriocolor}{q} & \textcolor{gregoriocolor}{r} & \textcolor{gregoriocolor}{s} & \textcolor{gregoriocolor}{t} & \textcolor{gregoriocolor}{u} \\
 10 & 11 & 12 & 13 & 14 & 15 & 16 & 17 & 18 & 19 & 20 & 21 & 22 & 23 & 24 & 25 & 26 & 27 & 28 \\
\end{tabularx}
\vspace{0.5\baselineskip}
\noindent\begin{tabularx}{\linewidth}{*{12}{>{\centering\arraybackslash}X}}
 \textcolor{gregoriocolor}{A} & \textcolor{gregoriocolor}{B} & \textcolor{gregoriocolor}{C} & \textcolor{gregoriocolor}{D} & \textcolor{gregoriocolor}{E} & F & \textcolor{gregoriocolor}{F} & \textcolor{gregoriocolor}{G} & \textcolor{gregoriocolor}{H} & \textcolor{gregoriocolor}{M} & \textcolor{gregoriocolor}{N} & \textcolor{gregoriocolor}{P} \\
 29 & 1 & 2 & 3 & 4 & 5 & 4 & 5 & 6 & 7 & 8 & 9 \\
\end{tabularx}

\begin{paracol}{2}
\selectlanguage{latin}
\lettrine[lines=2]{S}{ancti} Joánnis de Matha, Presbyteri et Confessóris, qui Ordinis sanctíssimæ 
 Trinitátis redemptiónis captivórum fuit Institútor, et sextodécimo Kaléndas 
 Januárii obdormívit in Dómino.
\switchcolumn
\selectlanguage{english}
\lettrine[lines=2]{S}{t.} John of Matha, priest and confessor, founder of the Order of the Most 
 Holy Trinity for the redemption of captives, who went to repose in the Lord 
 on the 17th of December.
\switchcolumn*
\selectlanguage{latin}
Somáschæ, in território Bergoménsi, natális 
 sancti Hierónymi Æmiliáni Confessóris, qui Congregatiónis Somáschæ Fundátor 
 éxstitit; atque, plúribus in vita et post mortem miráculis illústris, a 
 Cleménte Décimo tértio, Pontífice Máximo, Sanctórum fastis adscríptus est, 
 et a Pio Papa Undécimo universális orphanórum ac derelíctæ juventútis 
 Patrónus apud Deum eléctus et declarátus. Ejus tamen festívitas 
 tertiodécimo Kaléndas Augústi recólitur.
\switchcolumn
\selectlanguage{english}
At Somascha, in the district of Bergamo, the birthday of St. Jerome Emilian, 
 confessor, who was the founder of the Congregation of Somascha. 
 Illustrious both during his life and after death for many miracles, he was 
 inscribed in the roll of the saints by Pope Clement XIII. Pope Pius XI 
 chose and declared him to be the heavenly patron of orphans and abandoned 
 children. His feast is celebrated on the 20th of July.
\switchcolumn*
\selectlanguage{latin}
Romæ sanctórum Mártyrum Pauli, Lúcii et Cyríaci.
\switchcolumn
\selectlanguage{english}
At Rome, the holy martyrs Paul, Lucius, and Cyriacus.
\switchcolumn*
\selectlanguage{latin}
In Arménia minóre pássio sanctórum Mártyrum Dionysii,
 Æmiliáni et Sebastiáni.
\switchcolumn
\selectlanguage{english}
In Lesser Armenia, the birthday of the holy martyrs Denis, Aemilian, and 
 Sebastian.
\switchcolumn*
\selectlanguage{latin}
Constantinópoli natális sanctórum Mártyrum Monachórum monastérii Dii, 
 qui, ob defensiónem fídei cathólicæ, cum 
 tulíssent lítteras sancti Felícis Papæ Tértii advérsus Acácium, diríssime 
 cæsi sunt.
\switchcolumn
\selectlanguage{english}
At Constantinople, the birthday of the holy martyrs, monks of the monastery 
 of Dius. While bringing the letter of Pope St. Felix against Acacius, 
 they were barbarously killed for their defence of the Catholic faith.
\switchcolumn*
\selectlanguage{latin}
In Pérside commemorátio sanctórum Mártyrum, qui, sub Rege Persárum Cábade, 
 ob Christiánam fidem, divérsis suplíciis necáti sunt.
\switchcolumn
\selectlanguage{english}
In Persia, in the time of King Cabades, the commemoration of the holy 
 martyrs, who were put to death by various kinds of torments on account of 
 their Christian faith.
\switchcolumn*
\selectlanguage{latin}
Alexandríæ pássio sanctæ Coínthæ Mártyris, quam 
 Pagáni, sub Décio Imperatóre, corréptam et ad idóla perdúctam, hæc adoráre 
 cogébant; quod cum illa éxsecrans recusáret, ipsíus pedes vínculis 
 innexuérunt, eámque, trahéntes sic vinctam per civitátis platéas, horréndo 
 supplício discerpsérunt.
\switchcolumn
\selectlanguage{english}
At Alexandria, under Emperor Decius, the martyr St. Cointha, whom the pagans 
 seized, led to the idols, and urged to adore them. As she refused with 
 horror, they put her feet in chains, and dragged her through the streets of 
 the city, mangling her body in a most barbarous manner.
\switchcolumn*
\selectlanguage{latin}
Papíæ sancti Juvéntii Epíscopi, qui strénue in 
 Evangélio laborávit.
\switchcolumn
\selectlanguage{english}
At Pavia, St. Juventius, bishop, who laboured with zeal in preaching the 
 Gospel.
\switchcolumn*
\selectlanguage{latin}
Medioláni deposítio sancti Honoráti, Epíscopi 
 et Confessóris.
\switchcolumn
\selectlanguage{english}
At Milan, the death of St. Honoratus, bishop and confessor.
\switchcolumn*
\selectlanguage{latin}
Virodúni, in Gállia, sancti Pauli Epíscopi, miraculórum dono illústris.
\switchcolumn
\selectlanguage{english}
At Verdun in France, St. Paul, a bishop renowned for his miracles.
\switchcolumn*
\selectlanguage{latin}
Apud Murétum, in agro Lemovicénsi, natális sancti Stéphani Abbátis, qui 
 Grandimonténsis Ordinis Institútor fuit, ac virtútibus et miráculis cláruit.
\switchcolumn
\selectlanguage{english}
At Muret, near Limoges, the birthday of the abbot St. Stephen, founder of 
 the order of Grandmont, celebrated for his virtues and miracles.
\switchcolumn*
\selectlanguage{latin}
In monastério Vallis Umbrósæ Beáti Petri, 
 Cardinális et Epíscopi Albanénsis, ex Ordine Vallis Umbrósæ, cognoménto 
 Ignei, quia per ignem illæsus transívit.
\switchcolumn
\selectlanguage{english}
In the monastery of Vallombrosa, blessed Peter, cardinal and bishop of 
 Albano, a member of the Congregation of Vallombrosa of the Order of St. 
 Benedict. He was surnamed Igneus because he passed through fire 
 unharmed.
\switchcolumn*
\selectlanguage{latin}
\end{paracol}


% ---- martyrology/mart02/mart0209.htm
\needspace{10\baselineskip}
\begin{paracol}{2}
\selectlanguage{latin}
\begin{center}{\color{gregoriocolor} Quinto Idus Februárii. 
 Luna\dots\ }\end{center}
\switchcolumn
\selectlanguage{english}
\begin{center}{\color{gregoriocolor} The 9\textsuperscript{th} Day of 
 February. The\dots\ Day of the Moon.}\end{center}
\end{paracol}

\noindent\begin{tabularx}{\linewidth}{*{19}{>{\centering\arraybackslash}X}}
 \textcolor{gregoriocolor}{a} & \textcolor{gregoriocolor}{b} & \textcolor{gregoriocolor}{c} & \textcolor{gregoriocolor}{d} & \textcolor{gregoriocolor}{e} & \textcolor{gregoriocolor}{f} & \textcolor{gregoriocolor}{g} & \textcolor{gregoriocolor}{h} & \textcolor{gregoriocolor}{i} & \textcolor{gregoriocolor}{k} & \textcolor{gregoriocolor}{l} & \textcolor{gregoriocolor}{m} & \textcolor{gregoriocolor}{n} & \textcolor{gregoriocolor}{p} & \textcolor{gregoriocolor}{q} & \textcolor{gregoriocolor}{r} & \textcolor{gregoriocolor}{s} & \textcolor{gregoriocolor}{t} & \textcolor{gregoriocolor}{u} \\
 11 & 12 & 13 & 14 & 15 & 16 & 17 & 18 & 19 & 20 & 21 & 22 & 23 & 24 & 25 & 26 & 27 & 28 & 29 \\
\end{tabularx}
\vspace{0.5\baselineskip}
\noindent\begin{tabularx}{\linewidth}{*{12}{>{\centering\arraybackslash}X}}
 \textcolor{gregoriocolor}{A} & \textcolor{gregoriocolor}{B} & \textcolor{gregoriocolor}{C} & \textcolor{gregoriocolor}{D} & \textcolor{gregoriocolor}{E} & F & \textcolor{gregoriocolor}{F} & \textcolor{gregoriocolor}{G} & \textcolor{gregoriocolor}{H} & \textcolor{gregoriocolor}{M} & \textcolor{gregoriocolor}{N} & \textcolor{gregoriocolor}{P} \\
 1 & 2 & 3 & 4 & 5 & 6 & 5 & 6 & 7 & 8 & 9 & 10 \\
\end{tabularx}

\begin{paracol}{2}
\selectlanguage{latin}
\lettrine[lines=2]{S}{ancti} Cyrílli, Epíscopi Alexandríni, Confessóris et Ecclésiæ 
 Doctóris; cujus dies natális quinto Kaléndas Februárii recensétur.
\switchcolumn
\selectlanguage{english}
\lettrine[lines=2]{S}{t.} Cyril, bishop of Alexandria, confessor and doctor of the Church. 
 His birthday was mentioned on the 28th of January.
\switchcolumn*
\selectlanguage{latin}
Alexandríæ natális sanctæ Apollóniæ, Vírginis 
 et Mártyris, cui persecutóres, sub Décio, dentes omnes primum excussérunt. 
 Deínde, constrúcto ac succénso rogo, iídem commináti sunt, nisi cum eis ímpia verba proférret, vivam se eam incensúros; at illa, cum páululum intra 
 semetípsam deliberásset, repénte se de mánibus impiórum prorípuit, et in 
 ignem, quem paráverant, majóre Sancti Spíritus flamma intus æstuans, sponte 
 ita prosilívit, ut perterreréntur étiam ipsi crudelitátis auctóres, quod 
 prómptior invénta esset ad mortem fémina quam persecútor ad pœnam.
\switchcolumn
\selectlanguage{english}
At Alexandria, in the reign of Decius, the birthday of St. Apollonia, 
 virgin, who had all her teeth broken out by the persecutors; then, having 
 constructed and lighted a pyre, they threatened to burn her alive unless she 
 uttered with them certain impious words. Deliberating a while within 
 herself, she suddenly slipped from their grasp, and prompted by the greater 
 fire of the Holy Ghost within her, she rushed voluntarily into the fire 
 which they had prepared. Those responsible for her death were struck 
 with terror at the sight of a woman who was more willing to die than they to 
 kill her.
\switchcolumn*
\selectlanguage{latin}
Romæ pássio sanctórum Mártyrum Alexándri et 
 aliórum trigínta octo coronatórum.
\switchcolumn
\selectlanguage{english}
At Rome, the passion of the holy martyrs Alexander and thirty-eight others 
 crowned with him.
\switchcolumn*
\selectlanguage{latin}
In castéllo Lemelénsi, in Africa, sanctórum Mártyrum Primi et Donáti 
 Diaconórum, qui, cum altáre in Ecclésia tutaréntur, a Donatístis occísi sunt.
\switchcolumn
\selectlanguage{english}
In the village of Lamelum in Africa, the holy martyrs Primus and Donatus, 
 deacons, who were killed by the Donatists as they guarded the altar in the 
 church.
\switchcolumn*
\selectlanguage{latin}
Solis, in Cypro, sanctórum Mártyrum Ammónii et Alexándri.
\switchcolumn
\selectlanguage{english}
At Solum in Cyprus, the holy martyrs Ammonius and Alexander.
\switchcolumn*
\selectlanguage{latin}
Antiochíæ sancti Nicéphori Mártyris, qui sub 
 Valeriáno Imperatóre, cápite cæsus, martyrii corónam accépit.
\switchcolumn
\selectlanguage{english}
At Antioch, under Emperor Valerian, St. Nicephorus, martyr, who was beheaded 
 and thus received the crown of martyrdom.
\switchcolumn*
\selectlanguage{latin}
In monastério Fontanéllæ, in Gállia, sancti 
 Ansbérti, Rotomagénsis Epíscopi.
\switchcolumn
\selectlanguage{english}
In the monastery of Fontanelle in France, St. Ansbert, bishop of Rouen.
\switchcolumn*
\selectlanguage{latin}
Canúsii, in Apúlia, sancti Sabíni, Epíscopi et Confessóris; qui (ut beátus 
 Gregórius Papa refert), prophetíæ spíritu ac 
 miraculórum dono præditus, sibi jam cæco exhíbitum a fámulo, præmiis 
 corrúpto, venéni póculum divino agnóvit instínctu, sed, prænuntiáta mox a 
 Deo suménda de corruptóre vindícta signóque Crucis facto, venénum secúrus 
 ebíbit ac nullum ex eo nocuméntum accépit.
\switchcolumn
\selectlanguage{english}
At Canossa in Apulia, St. Sabinus, bishop and confessor. Blessed Pope 
 Gregory tells that he was endowed with the spirit of prophecy and the power 
 of miracles. After he had become blind, when a cup of poison was 
 offered to him by a servant who was bribed, he knew it by divine instinct. 
 He, however, declared that God would punish the one who had bribed the 
 servant, and, making the sign of the cross, he drank the poison without 
 anxiety and without harmful effect.
\switchcolumn*
\selectlanguage{latin}
\end{paracol}


% ---- martyrology/mart02/mart0210.htm
\needspace{10\baselineskip}
\begin{paracol}{2}
\selectlanguage{latin}
\begin{center}{\color{gregoriocolor} Quarto Idus Februárii. 
 Luna\dots\ }\end{center}
\switchcolumn
\selectlanguage{english}
\begin{center}{\color{gregoriocolor} The 10\textsuperscript{th} Day of 
 February. The\dots\ Day of the Moon.}\end{center}
\end{paracol}

\noindent\begin{tabularx}{\linewidth}{*{19}{>{\centering\arraybackslash}X}}
 \textcolor{gregoriocolor}{a} & \textcolor{gregoriocolor}{b} & \textcolor{gregoriocolor}{c} & \textcolor{gregoriocolor}{d} & \textcolor{gregoriocolor}{e} & \textcolor{gregoriocolor}{f} & \textcolor{gregoriocolor}{g} & \textcolor{gregoriocolor}{h} & \textcolor{gregoriocolor}{i} & \textcolor{gregoriocolor}{k} & \textcolor{gregoriocolor}{l} & \textcolor{gregoriocolor}{m} & \textcolor{gregoriocolor}{n} & \textcolor{gregoriocolor}{p} & \textcolor{gregoriocolor}{q} & \textcolor{gregoriocolor}{r} & \textcolor{gregoriocolor}{s} & \textcolor{gregoriocolor}{t} & \textcolor{gregoriocolor}{u} \\
 12 & 13 & 14 & 15 & 16 & 17 & 18 & 19 & 20 & 21 & 22 & 23 & 24 & 25 & 26 & 27 & 28 & 29 & 1 \\
\end{tabularx}
\vspace{0.5\baselineskip}
\noindent\begin{tabularx}{\linewidth}{*{12}{>{\centering\arraybackslash}X}}
 \textcolor{gregoriocolor}{A} & \textcolor{gregoriocolor}{B} & \textcolor{gregoriocolor}{C} & \textcolor{gregoriocolor}{D} & \textcolor{gregoriocolor}{E} & F & \textcolor{gregoriocolor}{F} & \textcolor{gregoriocolor}{G} & \textcolor{gregoriocolor}{H} & \textcolor{gregoriocolor}{M} & \textcolor{gregoriocolor}{N} & \textcolor{gregoriocolor}{P} \\
 2 & 3 & 4 & 5 & 6 & 7 & 6 & 7 & 8 & 9 & 10 & 11 \\
\end{tabularx}

\begin{paracol}{2}
\selectlanguage{latin}
\lettrine[lines=2]{A}{pud} montem Cassínum sanctæ Scholásticæ 
 Vírginis, soróris sancti Benedícti Abbátis; qui ejus ánimam, instar colúmbæ, 
 migrántem e córpore in cælum ascéndere vidit.
\switchcolumn
\selectlanguage{english}
\lettrine[lines=2]{O}{n} Monte Cassino, St. Scholastica, virgin, whose soul was seen by her 
 brother, St. Benedict, abbot, leaving her body in the form of a dove, and 
 ascending into heaven.
\switchcolumn*
\selectlanguage{latin}
Romæ sanctórum Mártyrum Zótici, Irenǽi, 
 Hyacínthi et Amántii.
\switchcolumn
\selectlanguage{english}
At Rome, the holy martyrs Zoticus, Irenaeus, Hyacinth, and Amantius.
\switchcolumn*
\selectlanguage{latin}
Ibídem, via Lavicána, sanctórum decem mílitum Mártyrum.
\switchcolumn
\selectlanguage{english}
In the same place, on the Via Lavicana, ten holy soldiers, martyrs.
\switchcolumn*
\selectlanguage{latin}
Item Romæ, via Appia, sanctæ Sotéris, Vírginis 
 et Mártyris; quæ (ut scribit sanctus Ambrósius), nóbili génere nata, paréntum 
 Consulátus et Præfectúras ob Christum contémpsit. Hæc, jussa idolis 
 immoláre, et non acquiéscens, gráviter et diutíssime álapis cæsa est; et, 
 cum cétera quoque pœnárum génera vicísset, demum, percússa gládio, læta 
 migrávit ad Sponsum.
\switchcolumn
\selectlanguage{english}
Also at Rome, on the Appian Way, St. Soter, virgin and martyr, descended of 
 a noble family, but as St. Ambrose mentions, for the love of Christ she set 
 at naught the consular and other dignitaries of her people. Upon her 
 refusal to sacrifice to the gods, she was for a long time cruelly scourged. 
 She overcame these and various other torments, then was struck with the 
 sword; and joyfully went to her heavenly spouse.
\switchcolumn*
\selectlanguage{latin}
In Campánia sancti Silváni, Epíscopi et Confessóris.
\switchcolumn
\selectlanguage{english}
In Campania, St. Silvanus, bishop and confessor.
\switchcolumn*
\selectlanguage{latin}
In Stábulo Rhodis, in territorio Senénsi, sancti Guiliélmi Eremítæ.
\switchcolumn
\selectlanguage{english}
At Malavalle, near Siena, St. William, hermit.
\switchcolumn*
\selectlanguage{latin}
In pago Rotomagénsi sanctæ Austrebértæ Vírginis, 
 miráculis célebris.
\switchcolumn
\selectlanguage{english}
In the diocese of Rouen, St. Austreberta, virgin, renowned for miracles.
\switchcolumn*
\selectlanguage{latin}
\end{paracol}


% ---- martyrology/mart02/mart0211.htm
\needspace{10\baselineskip}
\begin{paracol}{2}
\selectlanguage{latin}
\begin{center}{\color{gregoriocolor} Tértio Idus Februárii. 
 Luna\dots\ }\end{center}
\switchcolumn
\selectlanguage{english}
\begin{center}{\color{gregoriocolor} The 11\textsuperscript{th} Day of 
 February. The\dots\ Day of the Moon.}\end{center}
\end{paracol}

\noindent\begin{tabularx}{\linewidth}{*{19}{>{\centering\arraybackslash}X}}
 \textcolor{gregoriocolor}{a} & \textcolor{gregoriocolor}{b} & \textcolor{gregoriocolor}{c} & \textcolor{gregoriocolor}{d} & \textcolor{gregoriocolor}{e} & \textcolor{gregoriocolor}{f} & \textcolor{gregoriocolor}{g} & \textcolor{gregoriocolor}{h} & \textcolor{gregoriocolor}{i} & \textcolor{gregoriocolor}{k} & \textcolor{gregoriocolor}{l} & \textcolor{gregoriocolor}{m} & \textcolor{gregoriocolor}{n} & \textcolor{gregoriocolor}{p} & \textcolor{gregoriocolor}{q} & \textcolor{gregoriocolor}{r} & \textcolor{gregoriocolor}{s} & \textcolor{gregoriocolor}{t} & \textcolor{gregoriocolor}{u} \\
 13 & 14 & 15 & 16 & 17 & 18 & 19 & 20 & 21 & 22 & 23 & 24 & 25 & 26 & 27 & 28 & 29 & 1 & 2 \\
\end{tabularx}
\vspace{0.5\baselineskip}
\noindent\begin{tabularx}{\linewidth}{*{12}{>{\centering\arraybackslash}X}}
 \textcolor{gregoriocolor}{A} & \textcolor{gregoriocolor}{B} & \textcolor{gregoriocolor}{C} & \textcolor{gregoriocolor}{D} & \textcolor{gregoriocolor}{E} & F & \textcolor{gregoriocolor}{F} & \textcolor{gregoriocolor}{G} & \textcolor{gregoriocolor}{H} & \textcolor{gregoriocolor}{M} & \textcolor{gregoriocolor}{N} & \textcolor{gregoriocolor}{P} \\
 3 & 4 & 5 & 6 & 7 & 8 & 7 & 8 & 9 & 10 & 11 & 12 \\
\end{tabularx}

\begin{paracol}{2}
\selectlanguage{latin}
\lettrine[lines=2]{L}{apúrdi,} in Gállia, Apparítio beátæ Maríæ 
 Vírginis Immaculátæ.
\switchcolumn
\selectlanguage{english}
\lettrine[lines=2]{A}{t} Lourdes in France, the apparition of Blessed Mary, Virgin Immaculate.
\switchcolumn*
\selectlanguage{latin}
Hadrianópoli, in Thrácia, sanctórum Mártyrum Lúcii Epíscopi, et Sociórum 
 ejus, sub Constántio. Ipse Lúcius, ab Ariánis multa perpéssus, in 
 vínculis martyrium consummávit; céteri vero nobilióres cívium, cum Ariános, 
 in Sardicénsi Concílio tunc damnátos, recípere noluíssent, a Philágrio 
 Cómite capitálem senténtiam excepérunt.
\switchcolumn
\selectlanguage{english}
At Adrianople, the holy martyrs Lucius, bishop, and his companions. 
 Lucius suffered much from the Arians under Constantius, and completed his 
 martyrdom in prison. The others, among the foremost citizens, refusing 
 to communicate with the Arians, who were just condemned in the Council of 
 Sardica, were sentenced to capital punishment by the count Philagrius.
\switchcolumn*
\selectlanguage{latin}
In Africa natális sanctórum Mártyrum Saturníni Presbyteri, Datívi, Felícis, 
 Ampélii et Sociórum, qui, in persecutióne Diocletiáni, cum ad Domínicum ex 
 more celebrándum conveníssent, idcírco, a milítibus comprehénsi, sub Anolíno 
 Procónsule passi sunt.
\switchcolumn
\selectlanguage{english}
In Africa, during the persecution of Diocletian, the birthday of the holy 
 martyrs Saturninus, a priest, Davitus, Felix, Ampelius, and their 
 companions. They had, as was their custom, assembled for Mass when 
 they were seized by the soldiers and put to death, under the proconsul 
 Anolinus.
\switchcolumn*
\selectlanguage{latin}
In Numídia commemorátio plurimórum sanctórum Mártyrum, qui, in 
 eádem 
 persecutióne comprehénsi sunt, et, cum juxta Imperatóris edíctum divínas 
 Scriptúras trádere noluíssent, gravíssimis excruciáti sunt supplíciis, ac 
 tandem occísi.
\switchcolumn
\selectlanguage{english}
In Numidia, in the same persecution, the commemoration of many holy martyrs, 
 who, refusing after their apprehension to deliver the holy Scriptures in 
 conformity with an imperial edict, were given over to most painful torments 
 and slain.
\switchcolumn*
\selectlanguage{latin}
Romæ sancti Gregórii Papæ Secúndi, qui Leónis 
 Isáurici impietáti acérrime réstitit, et sanctum Bonifátium ad prædicándum 
 Evangélium in Germániam misit.
\switchcolumn
\selectlanguage{english}
At Rome, Pope St. Gregory II, who courageously withstood the impiety of Leo 
 the Isaurian, and sent St. Boniface to preach the Gospel in Germany.
\switchcolumn*
\selectlanguage{latin}
Item Romæ sancti Paschális Papæ Primi, qui 
 plúrima sanctórum Mártyrum córpora levávit e cryptis, éaque in divérsis 
 Urbis Ecclésiis honorífice collocávit.
\switchcolumn
\selectlanguage{english}
Also at Rome, Pope St. Paschal I, who raised many bodies of the holy martyrs 
 from their crypts, and buried them with honour in various churches in the 
 city.
\switchcolumn*
\selectlanguage{latin}
Ravénnæ sancti Calóceri, Epíscopi et 
 Confessóris.
\switchcolumn
\selectlanguage{english}
At Ravenna, St. Calocerus, bishop and confessor.
\switchcolumn*
\selectlanguage{latin}
Medioláni sancti Lázari Epíscopi.
\switchcolumn
\selectlanguage{english}
At Milan, St. Lazarus, bishop.
\switchcolumn*
\selectlanguage{latin}
Cápuæ sancti Castrénsis Epíscopi.
\switchcolumn
\selectlanguage{english}
At Capua, St. Castrensis, bishop.
\switchcolumn*
\selectlanguage{latin}
In Castro Nantoniénsi, in Gállia, sancti Severíni, qui fuit Abbas monastérii 
 Agaunénsis, suísque précibus cultórem Dei, Regem Clodovéum, 
 a diútina infirmitáte liberávit.
\switchcolumn
\selectlanguage{english}
At Chateau Landon in France, St. Severin, abbot of the monastery of Agaune, 
 by whose prayers the Christian king Clovis was delivered from a long 
 sickness.
\switchcolumn*
\selectlanguage{latin}
In Ægypto sancti Jonæ Mónachi, virtútibus clari.
\switchcolumn
\selectlanguage{english}
In Egypt, St. Jonas, a monk, eminent for his virtues.
\switchcolumn*
\selectlanguage{latin}
Viénnæ, in Gállia, Translátio córporis sancti 
 Desidérii, Epíscopi et Mártyris, ex território Lugdunénsi, in quo ipse olim 
 passus fúerat décimo Kaléndas Júnii.
\switchcolumn
\selectlanguage{english}
At Vienne in France, the translation of the body of St. Desiderius, bishop 
 and martyr, from the district of Lyons where he had died on the 23rd of May.
\switchcolumn*
\selectlanguage{latin}
\end{paracol}


% ---- martyrology/mart02/mart0212.htm
\needspace{10\baselineskip}
\begin{paracol}{2}
\selectlanguage{latin}
\begin{center}{\color{gregoriocolor} Prídie Idus Februárii. 
 Luna\dots\ }\end{center}
\switchcolumn
\selectlanguage{english}
\begin{center}{\color{gregoriocolor} The 12\textsuperscript{th} Day of 
 February. The\dots\ Day of the Moon.}\end{center}
\end{paracol}

\noindent\begin{tabularx}{\linewidth}{*{19}{>{\centering\arraybackslash}X}}
 \textcolor{gregoriocolor}{a} & \textcolor{gregoriocolor}{b} & \textcolor{gregoriocolor}{c} & \textcolor{gregoriocolor}{d} & \textcolor{gregoriocolor}{e} & \textcolor{gregoriocolor}{f} & \textcolor{gregoriocolor}{g} & \textcolor{gregoriocolor}{h} & \textcolor{gregoriocolor}{i} & \textcolor{gregoriocolor}{k} & \textcolor{gregoriocolor}{l} & \textcolor{gregoriocolor}{m} & \textcolor{gregoriocolor}{n} & \textcolor{gregoriocolor}{p} & \textcolor{gregoriocolor}{q} & \textcolor{gregoriocolor}{r} & \textcolor{gregoriocolor}{s} & \textcolor{gregoriocolor}{t} & \textcolor{gregoriocolor}{u} \\
 14 & 15 & 16 & 17 & 18 & 19 & 20 & 21 & 22 & 23 & 24 & 25 & 26 & 27 & 28 & 29 & 1 & 2 & 3 \\
\end{tabularx}
\vspace{0.5\baselineskip}
\noindent\begin{tabularx}{\linewidth}{*{12}{>{\centering\arraybackslash}X}}
 \textcolor{gregoriocolor}{A} & \textcolor{gregoriocolor}{B} & \textcolor{gregoriocolor}{C} & \textcolor{gregoriocolor}{D} & \textcolor{gregoriocolor}{E} & F & \textcolor{gregoriocolor}{F} & \textcolor{gregoriocolor}{G} & \textcolor{gregoriocolor}{H} & \textcolor{gregoriocolor}{M} & \textcolor{gregoriocolor}{N} & \textcolor{gregoriocolor}{P} \\
 4 & 5 & 6 & 7 & 8 & 9 & 8 & 9 & 10 & 11 & 12 & 13 \\
\end{tabularx}

\begin{paracol}{2}
\selectlanguage{latin}
\lettrine[lines=2]{S}{anctórum} septem Fundatórum Ordinis Servórum beátæ 
 Maríæ Vírginis, Confessórum, quorum deposítio respectívis diébus recólitur. 
 Quos autem in vita unus veræ fraternitátis spíritus sociávit, et indivísa 
 post óbitum venerátio pópuli prosecúta est, eos Leo Décimus tértius, 
 Póntifex Máximus, una páriter Sanctórum fastis accénsuit.
\switchcolumn
\selectlanguage{english}
\lettrine[lines=2]{T}{he} seven Holy Founders of the Order of Servites of the Blessed Virgin Mary, 
 whose deaths are noted on their respective days. As one spirit of true 
 fraternal love united them in life, and as the people joined them together 
 in the same veneration after death, Pope Leo XIII placed them together in 
 the catalogue of the saints.
\switchcolumn*
\selectlanguage{latin}
In Africa sancti Damiáni, mílitis et Mártyris.
\switchcolumn
\selectlanguage{english}
In Africa, St. Damian, soldier and martyr.
\switchcolumn*
\selectlanguage{latin}
Carthágine sanctórum Mártyrum Modésti et Juliáni.
\switchcolumn
\selectlanguage{english}
At Carthage, the holy martyrs Modestus and Julian.
\switchcolumn*
\selectlanguage{latin}
Alexandríæ sanctórum Mártyrum Modésti et 
 Ammónii infántum.
\switchcolumn
\selectlanguage{english}
At Alexandria, the holy children Modestus and Ammonius, martyrs.
\switchcolumn*
\selectlanguage{latin}
Barcinóne, in Hispánia, sanctæ Euláliæ Vírginis, 
 quæ, témpore Diocletiáni Imperatóris, equúleum, úngulas flammásque perpéssa, 
 demum, cruci affíxa, gloriósam martyrii corónam accépit.
\switchcolumn
\selectlanguage{english}
At Barcelona in Spain, in the time of Emperor Diocletian, St. Eulalia, 
 virgin, who, being racked, torn with iron hooks, cast into the fire, and 
 crucified, received the glorious crown of martyrdom.
\switchcolumn*
\selectlanguage{latin}
Constantinópoli sancti Melétii, Epíscopi Antiochéni, qui, pro fide cathólica 
 sæpe exsílium passus, demum in eádem urbe 
 migrávit ad Dóminum. Ejus virtútes sanctus Joánnes Chrysóstomus et 
 sanctus Gregórius Nyssénus summis láudibus celebrárunt.
\switchcolumn
\selectlanguage{english}
At Constantinople, St. Meletius, bishop of Antioch, who often suffered exile 
 for the Catholic faith, and finally died at Constantinople and went to his 
 reward. His virtues have been extolled by St. John Chrysostom and St. 
 Gregory of Nyssa.
\switchcolumn*
\selectlanguage{latin}
Item Constantinópoli sancti Antónii Epíscopi, témpore Leónis sexti 
 Imperatóris.
\switchcolumn
\selectlanguage{english}
Also at Constantinople, St. Anthony, a bishop in the time of Emperor Leo VI.
\switchcolumn*
\selectlanguage{latin}
Verónæ sancti Gaudéntii, Epíscopi et 
 Confessóris.
\switchcolumn
\selectlanguage{english}
At Verona, St. Gaudentius, bishop and confessor.
\switchcolumn*
\selectlanguage{latin}
\end{paracol}


% ---- martyrology/mart02/mart0213.htm
\needspace{10\baselineskip}
\begin{paracol}{2}
\selectlanguage{latin}
\begin{center}{\color{gregoriocolor} Idibus Februárii. 
 Luna\dots\ }\end{center}
\switchcolumn
\selectlanguage{english}
\begin{center}{\color{gregoriocolor} The 13\textsuperscript{th} Day of 
 February. The\dots\ Day of the Moon.}\end{center}
\end{paracol}

\noindent\begin{tabularx}{\linewidth}{*{19}{>{\centering\arraybackslash}X}}
 \textcolor{gregoriocolor}{a} & \textcolor{gregoriocolor}{b} & \textcolor{gregoriocolor}{c} & \textcolor{gregoriocolor}{d} & \textcolor{gregoriocolor}{e} & \textcolor{gregoriocolor}{f} & \textcolor{gregoriocolor}{g} & \textcolor{gregoriocolor}{h} & \textcolor{gregoriocolor}{i} & \textcolor{gregoriocolor}{k} & \textcolor{gregoriocolor}{l} & \textcolor{gregoriocolor}{m} & \textcolor{gregoriocolor}{n} & \textcolor{gregoriocolor}{p} & \textcolor{gregoriocolor}{q} & \textcolor{gregoriocolor}{r} & \textcolor{gregoriocolor}{s} & \textcolor{gregoriocolor}{t} & \textcolor{gregoriocolor}{u} \\
 15 & 16 & 17 & 18 & 19 & 20 & 21 & 22 & 23 & 24 & 25 & 26 & 27 & 28 & 29 & 1 & 2 & 3 & 4 \\
\end{tabularx}
\vspace{0.5\baselineskip}
\noindent\begin{tabularx}{\linewidth}{*{12}{>{\centering\arraybackslash}X}}
 \textcolor{gregoriocolor}{A} & \textcolor{gregoriocolor}{B} & \textcolor{gregoriocolor}{C} & \textcolor{gregoriocolor}{D} & \textcolor{gregoriocolor}{E} & F & \textcolor{gregoriocolor}{F} & \textcolor{gregoriocolor}{G} & \textcolor{gregoriocolor}{H} & \textcolor{gregoriocolor}{M} & \textcolor{gregoriocolor}{N} & \textcolor{gregoriocolor}{P} \\
 5 & 6 & 7 & 8 & 9 & 10 & 9 & 10 & 11 & 12 & 13 & 14 \\
\end{tabularx}

\begin{paracol}{2}
\selectlanguage{latin}
\lettrine[lines=2]{A}{ntiochíæ} natális sancti Agábi Prophétæ, de quo 
 beátus Lucas in Actibus Apostólicis scribit.
\switchcolumn
\selectlanguage{english}
\lettrine[lines=2]{A}{t} Antioch, the birthday of St. Agabus, prophet, of whom mention is made by 
 St. Luke in the Acts of the Apostles.
\switchcolumn*
\selectlanguage{latin}
Tudérti, in Umbria, sancti Benígni, Presbyteri et Mártyris; qui, Diocletiáni 
 et Maximiáni Imperatórum témpore, cum fidem Christiánam verbo et exémplo 
 propagáre non desísteret, ab idolórum cultóribus captus est, ac, váriis 
 afféctus supplíciis, sacerdotále munus honóre martyrii cumulávit.
\switchcolumn
\selectlanguage{english}
At Todi in Umbria, St. Benignus, priest and martyr, who would not cease 
 spreading the Christian faith. In the reign of Emperors Diocletian and 
 Maximian he was taken by the pagans, suffered various tortures, and finally 
 reached the perfection of his priestly office with the honour of martyrdom.
\switchcolumn*
\selectlanguage{latin}
Melitínæ, in Arménia, sancti Polyéucti Mártyris, 
 qui in persecutióne Décii multa passus, martyrii corónam adéptus est.
\switchcolumn
\selectlanguage{english}
At Meletine in Armenia, in the persecution of Decius, St. Polyeuctus, who, 
 after many sufferings, obtained the crown of martyrdom.
\switchcolumn*
\selectlanguage{latin}
Lugdúni, in Gállia, sancti Juliáni Mártyris.
\switchcolumn
\selectlanguage{english}
At Lyons in France, St. Julian, martyr.
\switchcolumn*
\selectlanguage{latin}
Ravénnæ sanctárum Fuscæ Vírginis, ejúsque 
 nutrícis Mauræ; quæ, Décio imperánte, multa sub Quinctiáno Præside perpéssæ, 
 demum, gládio transfíxæ, martyrium consummárunt.
\switchcolumn
\selectlanguage{english}
At Ravenna, in the time of Emperor Decius and the governor Quinctian, the 
 Saints Fusca, virgin, and Maura, her nurse. They endured many 
 afflictions, but were finally transfixed with a sword, and thus ended their 
 martyrdom.
\switchcolumn*
\selectlanguage{latin}
Lugdúni, in Gállia, sancti Stéphani, Epíscopi et Confessóris.
\switchcolumn
\selectlanguage{english}
At Lyons in France, St. Stephen, bishop and confessor.
\switchcolumn*
\selectlanguage{latin}
Reáte sancti Stéphani Abbátis, miræ patiéntiæ 
 viri; in cujus tránsitu (ut refert beátus Gregórius Papa) sancti Angeli, 
 céteris étiam vidéntibus, adfuérunt.
\switchcolumn
\selectlanguage{english}
At Rieti, the abbot St. Stephen, a man of wonderful patience, at whose 
 death, as is related by blessed Pope Gregory, the holy angels were present 
 and visible to all.
\switchcolumn*
\selectlanguage{latin}
\end{paracol}


% ---- martyrology/mart02/mart0214.htm
\needspace{10\baselineskip}
\begin{paracol}{2}
\selectlanguage{latin}
\begin{center}{\color{gregoriocolor} Sextodécimo Kaléndas Mártii. 
 Luna\dots\ }\end{center}
\switchcolumn
\selectlanguage{english}
\begin{center}{\color{gregoriocolor} The 14\textsuperscript{th} Day of 
 February. The\dots\ Day of the Moon.}\end{center}
\end{paracol}

\noindent\begin{tabularx}{\linewidth}{*{19}{>{\centering\arraybackslash}X}}
 \textcolor{gregoriocolor}{a} & \textcolor{gregoriocolor}{b} & \textcolor{gregoriocolor}{c} & \textcolor{gregoriocolor}{d} & \textcolor{gregoriocolor}{e} & \textcolor{gregoriocolor}{f} & \textcolor{gregoriocolor}{g} & \textcolor{gregoriocolor}{h} & \textcolor{gregoriocolor}{i} & \textcolor{gregoriocolor}{k} & \textcolor{gregoriocolor}{l} & \textcolor{gregoriocolor}{m} & \textcolor{gregoriocolor}{n} & \textcolor{gregoriocolor}{p} & \textcolor{gregoriocolor}{q} & \textcolor{gregoriocolor}{r} & \textcolor{gregoriocolor}{s} & \textcolor{gregoriocolor}{t} & \textcolor{gregoriocolor}{u} \\
 16 & 17 & 18 & 19 & 20 & 21 & 22 & 23 & 24 & 25 & 26 & 27 & 28 & 29 & 1 & 2 & 3 & 4 & 5 \\
\end{tabularx}
\vspace{0.5\baselineskip}
\noindent\begin{tabularx}{\linewidth}{*{12}{>{\centering\arraybackslash}X}}
 \textcolor{gregoriocolor}{A} & \textcolor{gregoriocolor}{B} & \textcolor{gregoriocolor}{C} & \textcolor{gregoriocolor}{D} & \textcolor{gregoriocolor}{E} & F & \textcolor{gregoriocolor}{F} & \textcolor{gregoriocolor}{G} & \textcolor{gregoriocolor}{H} & \textcolor{gregoriocolor}{M} & \textcolor{gregoriocolor}{N} & \textcolor{gregoriocolor}{P} \\
 6 & 7 & 8 & 9 & 10 & 11 & 10 & 11 & 12 & 13 & 14 & 15 \\
\end{tabularx}

\begin{paracol}{2}
\selectlanguage{latin}
% NOTE: this is lines=1 because of a bad page break on p. 42... Can we fix this?
\lettrine[lines=1]{R}{omæ,} via Flamínia, natális sancti Valentíni, 
 Presbyteri et Mártyris, qui, post multa sanitátum et doctrínæ insígnia, 
 fústibus cæsus et decollátus est, sub Cláudio Cǽsare.
\switchcolumn
\selectlanguage{english}
\lettrine[lines=1]{A}{t} Rome, on the Flaminian Way, in the time of Emperor Claudius, the birthday 
 of St. Valentine, priest and martyr, who after having cured and instructed 
 many persons, was beaten with clubs and beheaded.
\switchcolumn*
\selectlanguage{latin}
Ibídem deposítio sancti Cyrílli, Epíscopi et Confessóris; qui, una cum 
 sancto Methódio, simíliter Epíscopo et fratre suo, cujus dies natális octávo 
 Idus Aprílis recensétur, multas Slávicas gentes 
 earúmque Reges ad fidem Christi perdúxit. Horum tamen Sanctórum 
 festívitas Nonis Júlii celebrátur.
\switchcolumn
\selectlanguage{english}
In the same place, St. Cyril, bishop, who together with his brother 
 Methodius, also a bishop, whose birthday is the 6th of April, brought many 
 people and the rulers of Moravia to the faith of Christ. Their feast 
 is celebrated on the 7th of July.
\switchcolumn*
\selectlanguage{latin}
Item Romæ sanctórum Mártyrum Vitális, Felículæ 
 et Zenónis.
\switchcolumn
\selectlanguage{english}
Also at Rome, the holy martyrs Vitalis, Felicula and Zeno.
\switchcolumn*
\selectlanguage{latin}
Iterámnæ sancti Valentíni, Epíscopi et Mártyris, 
 qui, post diútinam cædem mancipátus custódiæ, et, cum superári non posset, 
 tandem, médiæ noctis siléntio ejéctus de cárcere, decollátus est, jussu 
 Præfécti urbis Plácidi.
\switchcolumn
\selectlanguage{english}
At Teramo, St. Valentine, bishop and martyr, who was scourged, committed to 
 prison, and, because he remained unshaken in his faith, was taken out of his 
 dungeon in the dead of night and beheaded by order of Placidus, prefect of 
 the city.
\switchcolumn*
\selectlanguage{latin}
Alexandríæ sanctórum Mártyrum Cyriónis, 
 Presbyteri, Bassiáni Lectóris, Agathónis Exorcístæ, et Móysis; qui omnes, 
 igne combústi, evolavérunt ad cælum.
\switchcolumn
\selectlanguage{english}
At Alexandria, the holy martyrs Cyrion, priest; Bassian, lector; Agatho, 
 exorcist; and Moses, who perished in the flames and took their flight to 
 heaven.
\switchcolumn*
\selectlanguage{latin}
Interámnæ sanctórum Próculi, Ephébi et 
 Apollónii Mártyrum, qui, cum ad sancti Valentíni corpus vigílias ágerent, 
 Leóntii Consuláris jussu comprehénsi sunt, et gládio cæsi.
\switchcolumn
\selectlanguage{english}
At Teramo, the holy martyrs Proculus, Ephebus, and Apollonius, who, while 
 keeping watch at the body of St. Valentine, were arrested and put to the 
 sword by command of the consular officer, Leontius.
\switchcolumn*
\selectlanguage{latin}
Alexandríæ sanctórum Mártyrum Bassi, Antónii et 
 Protólici, qui demérsi sunt in mare.
\switchcolumn
\selectlanguage{english}
At Alexandria, the holy martyrs Bassus, Anthony, and Protolicus, who were 
 drowned in the sea.
\switchcolumn*
\selectlanguage{latin}
Item Alexandríæ sanctórum Dionysii et Ammónii 
 decollatórum.
\switchcolumn
\selectlanguage{english}
Also at Alexandria, the Saints Denis and Ammonius, who were beheaded.
\switchcolumn*
\selectlanguage{latin}
Neápoli, in Campánia, sancti Nostriáni Epíscopi, qui in cathólica fide 
 contra hæréticam pravitátem tuénda éxstitit 
 insígnis.
\switchcolumn
\selectlanguage{english}
At Naples, in Campania, St. Nostrian, bishop, who was outstanding for his 
 defence of the Catholic faith against heretical errors.
\switchcolumn*
\selectlanguage{latin}
Ravénnæ sancti Eleuchádii, Epíscopi et 
 Confessóris.
\switchcolumn
\selectlanguage{english}
At Ravenna, St. Eleuchadius, bishop and confessor.
\switchcolumn*
\selectlanguage{latin}
In Bithynia sancti Auxéntii Abbátis.
\switchcolumn
\selectlanguage{english}
In Bithynia, St. Auxentius, abbot.
\switchcolumn*
\selectlanguage{latin}
Apud Surréntum sancti Antoníni Abbátis, qui e monastério Cassinénsi, a 
 Longobárdis devastáto, in solitúdinem ejúsdem urbis secéssit; ibíque, 
 sanctitáte célebris, obdormívit in Dómino. Ipsíus corpus multis 
 quotídie miráculis, et præsértim in energúmenis 
 liberándis effúlget.
\switchcolumn
\selectlanguage{english}
At Sorrento, St. Anthony, abbot, who, when the monastery of Monte Cassino 
 was devastated by the Lombards, withdrew into a solitude of the 
 neighbourhood, where, celebrated for his holiness, he went calmly to his 
 repose in God. His body is daily glorified by many miracles, 
 especially by the deliverance of possessed persons.
\switchcolumn*
\selectlanguage{latin}
\end{paracol}


% ---- martyrology/mart02/mart0215.htm
\needspace{10\baselineskip}
\begin{paracol}{2}
\selectlanguage{latin}
\begin{center}{\color{gregoriocolor} Quintodécimo Kaléndas Mártii. 
 Luna\dots\ }\end{center}
\switchcolumn
\selectlanguage{english}
\begin{center}{\color{gregoriocolor} The 15\textsuperscript{th} Day of 
 February. The\dots\ Day of the Moon.}\end{center}
\end{paracol}

\noindent\begin{tabularx}{\linewidth}{*{19}{>{\centering\arraybackslash}X}}
 \textcolor{gregoriocolor}{a} & \textcolor{gregoriocolor}{b} & \textcolor{gregoriocolor}{c} & \textcolor{gregoriocolor}{d} & \textcolor{gregoriocolor}{e} & \textcolor{gregoriocolor}{f} & \textcolor{gregoriocolor}{g} & \textcolor{gregoriocolor}{h} & \textcolor{gregoriocolor}{i} & \textcolor{gregoriocolor}{k} & \textcolor{gregoriocolor}{l} & \textcolor{gregoriocolor}{m} & \textcolor{gregoriocolor}{n} & \textcolor{gregoriocolor}{p} & \textcolor{gregoriocolor}{q} & \textcolor{gregoriocolor}{r} & \textcolor{gregoriocolor}{s} & \textcolor{gregoriocolor}{t} & \textcolor{gregoriocolor}{u} \\
 17 & 18 & 19 & 20 & 21 & 22 & 23 & 24 & 25 & 26 & 27 & 28 & 29 & 1 & 2 & 3 & 4 & 5 & 6 \\
\end{tabularx}
\vspace{0.5\baselineskip}
\noindent\begin{tabularx}{\linewidth}{*{12}{>{\centering\arraybackslash}X}}
 \textcolor{gregoriocolor}{A} & \textcolor{gregoriocolor}{B} & \textcolor{gregoriocolor}{C} & \textcolor{gregoriocolor}{D} & \textcolor{gregoriocolor}{E} & F & \textcolor{gregoriocolor}{F} & \textcolor{gregoriocolor}{G} & \textcolor{gregoriocolor}{H} & \textcolor{gregoriocolor}{M} & \textcolor{gregoriocolor}{N} & \textcolor{gregoriocolor}{P} \\
 7 & 8 & 9 & 10 & 11 & 12 & 11 & 12 & 13 & 14 & 15 & 16 \\
\end{tabularx}

\begin{paracol}{2}
\selectlanguage{latin}
\lettrine[lines=2]{B}{ríxiæ} natális sanctórum Mártyrum Faustíni et 
 Jovítæ fratrum, qui sub Hadriáno Imperatóre, post multa præclára ob Christi 
 fidem suscépta certámina, victrícem martyrii corónam accepérunt.
\switchcolumn
\selectlanguage{english}
\lettrine[lines=2]{A}{t} Brescia, in the time of Emperor Adrian, the birthday of the holy martyrs 
 Faustinus and Jovita, who received the triumphant crown of martyrdom after 
 many glorious combats for the faith of Christ.
\switchcolumn*
\selectlanguage{latin}
Romæ sancti Cratónis Mártyris, qui, cum uxóre 
 sua et univérsa domo a beáto Valentíno Epíscopo baptizátus, non multo post, 
 una cum illis, martyrio consummátus est.
\switchcolumn
\selectlanguage{english}
At Rome, St. Craton, martyr. A short time after being baptized with 
 his wife and all his household by the holy bishop Valentine, he was put to 
 death with them.
\switchcolumn*
\selectlanguage{latin}
Interámnæ natális sanctórum Mártyrum Saturníni, 
 Cástuli, Magni et Lúcii.
\switchcolumn
\selectlanguage{english}
At Teramo, the birthday of the holy martyrs Saturninus, Castulus, Magnus, 
 and Lucius.
\switchcolumn*
\selectlanguage{latin}
Ibídem sanctæ Agapis, Vírginis et Mártyris.
\switchcolumn
\selectlanguage{english}
In the same place, St. Agape, virgin and martyr.
\switchcolumn*
\selectlanguage{latin}
Vasióne, in Gálliis, sancti Quinídii Epíscopi, cujus mortem in conspéctu 
 Dómini pretiósam mirácula crebra testántur.
\switchcolumn
\selectlanguage{english}
At Vaison in France, St. Quinidius, bishop, whose death was precious in the 
 sight of God, as is shewn by frequent miracles.
\switchcolumn*
\selectlanguage{latin}
Cápuæ sancti Decorósi, Epíscopi et Confessóris.
\switchcolumn
\selectlanguage{english}
At Capua, St. Decorosus, bishop and confessor.
\switchcolumn*
\selectlanguage{latin}
In Província Valériæ sancti Sevéri Presbyteri, 
 qui (ut beátus Gregórius Papa scribit), fusis lácrimis, defúnctum revocávit 
 ad vitam.
\switchcolumn
\selectlanguage{english}
In the province of Valeria, St. Severus, priest, of whom St. Gregory says 
 that by his tears he recalled a dead man to life.
\switchcolumn*
\selectlanguage{latin}
Antiochíæ sancti Joséphi Diáconi.
\switchcolumn
\selectlanguage{english}
At Antioch, St. Joseph, deacon.
\switchcolumn*
\selectlanguage{latin}
Arvérnis, in Gállia, sanctæ Geórgiæ Vírginis.
\switchcolumn
\selectlanguage{english}
In Auvergne in France, St. Georgia, virgin.
\switchcolumn*
\selectlanguage{latin}
\end{paracol}


% ---- martyrology/mart02/mart0216.htm
\needspace{10\baselineskip}
\begin{paracol}{2}
\selectlanguage{latin}
\begin{center}{\color{gregoriocolor} Quartodécimo Kaléndas Mártii. 
 Luna\dots\ }\end{center}
\switchcolumn
\selectlanguage{english}
\begin{center}{\color{gregoriocolor} The 16\textsuperscript{th} Day of 
 February. The\dots\ Day of the Moon.}\end{center}
\end{paracol}

\noindent\begin{tabularx}{\linewidth}{*{19}{>{\centering\arraybackslash}X}}
 \textcolor{gregoriocolor}{a} & \textcolor{gregoriocolor}{b} & \textcolor{gregoriocolor}{c} & \textcolor{gregoriocolor}{d} & \textcolor{gregoriocolor}{e} & \textcolor{gregoriocolor}{f} & \textcolor{gregoriocolor}{g} & \textcolor{gregoriocolor}{h} & \textcolor{gregoriocolor}{i} & \textcolor{gregoriocolor}{k} & \textcolor{gregoriocolor}{l} & \textcolor{gregoriocolor}{m} & \textcolor{gregoriocolor}{n} & \textcolor{gregoriocolor}{p} & \textcolor{gregoriocolor}{q} & \textcolor{gregoriocolor}{r} & \textcolor{gregoriocolor}{s} & \textcolor{gregoriocolor}{t} & \textcolor{gregoriocolor}{u} \\
 18 & 19 & 20 & 21 & 22 & 23 & 24 & 25 & 26 & 27 & 28 & 29 & 1 & 2 & 3 & 4 & 5 & 6 & 7 \\
\end{tabularx}
\vspace{0.5\baselineskip}
\noindent\begin{tabularx}{\linewidth}{*{12}{>{\centering\arraybackslash}X}}
 \textcolor{gregoriocolor}{A} & \textcolor{gregoriocolor}{B} & \textcolor{gregoriocolor}{C} & \textcolor{gregoriocolor}{D} & \textcolor{gregoriocolor}{E} & F & \textcolor{gregoriocolor}{F} & \textcolor{gregoriocolor}{G} & \textcolor{gregoriocolor}{H} & \textcolor{gregoriocolor}{M} & \textcolor{gregoriocolor}{N} & \textcolor{gregoriocolor}{P} \\
 8 & 9 & 10 & 11 & 12 & 13 & 12 & 13 & 14 & 15 & 16 & 17 \\
\end{tabularx}

\begin{paracol}{2}
\selectlanguage{latin}
\lettrine[lines=2]{R}{omæ} beáti Onésimi, de quo sanctus Paulus 
 Apóstolus ad Philémonem scribit; quem étiam, post sanctum Timótheum, 
 Ephesiórum Epíscopum ordinávit, prædicationísque verbum illi commisit. 
 Ipse autem Onésimus, vinctus Romam perdúctus ac pro fide Christi lapidátus, 
 primo ibídem sepúltus fuit; inde ad locum ubi Epíscopus fúerat ordinátus, 
 corpus ejus delátum est.
\switchcolumn
\selectlanguage{english}
\lettrine[lines=2]{A}{t} Rome, blessed Onesimus, concerning whom the apostle St. Paul wrote to 
 Philemon. He made him bishop of Ephesus after St. Timothy, and 
 committed to him the office of preaching. Being led a prisoner to 
 Rome, and stoned to death for the faith of Christ, he was first buried 
 there, but his body was afterwards taken to the place where he had been 
 bishop.
\switchcolumn*
\selectlanguage{latin}
In Ægypto sancti Juliáni Mártyris, cum áliis 
 quinque míllibus.
\switchcolumn
\selectlanguage{english}
In Egypt, St. Julian, martyr, with five thousand other Christians.
\switchcolumn*
\selectlanguage{latin}
Cæsaréæ, in Palæstína, sanctórum Mártyrum 
 Ægyptiórum Elíæ, Jeremíæ, Isaíæ, Samuélis et Daniélis; qui, cum spontánee 
 ministrássent Confessóribus in Cilícia ad metálla damnátis, et inde 
 reverteréntur, sunt comprehénsi, et a Firmiliáno Præside, sub Galério 
 Maximiáno Imperatóre, sævíssime torti, gládio demum percússi sunt. 
 Post eos sanctus Porphyrius, Pámphili Mártyris fámulus, et sanctus Seléucus 
 Cáppadox, qui iterátis certamínibus sæpe vícerant, rursus cruciáti sunt, 
 atque alter incéndio, gládio alter corónam martyrii accepérunt.
\switchcolumn
\selectlanguage{english}
At Caesarea, in Palestine, the holy martyrs Elias, Jeremias, Isaias, Samuel, 
 and Daniel. These Egyptians of their own accord ministered to the 
 confessors condemned to labour in the mines of Cilicia, but were arrested 
 upon their return, and after being cruelly tortured by the governor 
 Firmilian, under Emperor Galerius Maximian, were put to the sword. 
 After them, St. Porphyry, servant of the martyr Pamphilus, and St. Seleucus 
 the Cappadocian, who had been triumphant in several previous tests, being 
 again tortured, now won the crown of martyrdom, the one by fire, the other 
 by the sword.
\switchcolumn*
\selectlanguage{latin}
Nicomedíæ sanctæ Juliánæ, Vírginis et Mártyris; 
 quæ, sub Maximiáno Imperatóre, primum a patre suo Africáno gráviter cæsa, 
 deínde ab Evilásio Præfécto, cui núbere recusáverat, várie cruciáta, et 
 póstmodum in cárcerem detrúsa, ubi palam cum diábolo conflíxit, demum, cum 
 flammas ígnium et ollam fervéntem superásset, cápitis decollatióne martyrium 
 consummávit. Ipsíus autem corpus póstea Cumas, in Campánia translátum 
 est.
\switchcolumn
\selectlanguage{english}
At Nicomedia, St. Juliana, virgin and martyr. Under Emperor Maximian, 
 she was first severely scourged by her own father, Africanus, and then made 
 to suffer many torments by the prefect Evilasius, whom she had refused to 
 marry. Later thrown into prison, she encountered the evil spirit in a 
 visible manner. Finally, because the fiery furnace and a caldron of 
 boiling oil could do her no injury, her martyrdom was fulfilled by 
 beheading. Her body was later transferred to Cumi in Campania.
\switchcolumn*
\selectlanguage{latin}
Bríxiæ sancti Faustíni, Epíscopi et Confessóris.
\switchcolumn
\selectlanguage{english}
At Brescia, St. Faustinus, bishop and confessor.
\switchcolumn*
\selectlanguage{latin}
\end{paracol}


% ---- martyrology/mart02/mart0217.htm
\needspace{10\baselineskip}
\begin{paracol}{2}
\selectlanguage{latin}
\begin{center}{\color{gregoriocolor} Tertiodécimo Kaléndas Mártii. 
 Luna\dots\ }\end{center}
\switchcolumn
\selectlanguage{english}
\begin{center}{\color{gregoriocolor} The 17\textsuperscript{th} Day of 
 February. The\dots\ Day of the Moon.}\end{center}
\end{paracol}

\noindent\begin{tabularx}{\linewidth}{*{19}{>{\centering\arraybackslash}X}}
 \textcolor{gregoriocolor}{a} & \textcolor{gregoriocolor}{b} & \textcolor{gregoriocolor}{c} & \textcolor{gregoriocolor}{d} & \textcolor{gregoriocolor}{e} & \textcolor{gregoriocolor}{f} & \textcolor{gregoriocolor}{g} & \textcolor{gregoriocolor}{h} & \textcolor{gregoriocolor}{i} & \textcolor{gregoriocolor}{k} & \textcolor{gregoriocolor}{l} & \textcolor{gregoriocolor}{m} & \textcolor{gregoriocolor}{n} & \textcolor{gregoriocolor}{p} & \textcolor{gregoriocolor}{q} & \textcolor{gregoriocolor}{r} & \textcolor{gregoriocolor}{s} & \textcolor{gregoriocolor}{t} & \textcolor{gregoriocolor}{u} \\
 19 & 20 & 21 & 22 & 23 & 24 & 25 & 26 & 27 & 28 & 29 & 1 & 2 & 3 & 4 & 5 & 6 & 7 & 8 \\
\end{tabularx}
\vspace{0.5\baselineskip}
\noindent\begin{tabularx}{\linewidth}{*{12}{>{\centering\arraybackslash}X}}
 \textcolor{gregoriocolor}{A} & \textcolor{gregoriocolor}{B} & \textcolor{gregoriocolor}{C} & \textcolor{gregoriocolor}{D} & \textcolor{gregoriocolor}{E} & F & \textcolor{gregoriocolor}{F} & \textcolor{gregoriocolor}{G} & \textcolor{gregoriocolor}{H} & \textcolor{gregoriocolor}{M} & \textcolor{gregoriocolor}{N} & \textcolor{gregoriocolor}{P} \\
 9 & 10 & 11 & 12 & 13 & 14 & 13 & 14 & 15 & 16 & 17 & 18 \\
\end{tabularx}

\begin{paracol}{2}
\selectlanguage{latin}
\lettrine[lines=2]{F}{loréntiæ} natális sancti Aléxii Falconérii 
 Confessóris, e septem Fundatóribus Ordinis Servórum beátæ Maríæ Vírginis; 
 qui, décimo supra centésimum vitæ suæ anno, Christi Jesu et Angelórum 
 præséntia recreátus, beáto fine quiévit. Ipsíus tamen ac Sociórum 
 festum prídie Idus Februárii celebrátur.
\switchcolumn
\selectlanguage{english}
\lettrine[lines=2]{I}{n} Florence, the birthday of St. Alexis Falconieri, confessor, one of the 
 seven founders of the Order of the Servites of the Blessed Virgin Mary. 
 In the one hundred and tenth year of his age, he ended his blessed career in 
 the consoling presence of Christ Jesus and the angels. His feast, with 
 that of his companions, is kept on the 12th of February.
\switchcolumn*
\selectlanguage{latin}
Romæ pássio sancti Faustíni, quem álii 
 quadragínta quátuor secúti sunt ad corónam.
\switchcolumn
\selectlanguage{english}
At Rome, the passion of St. Faustinus, whom forty-four others followed to 
 receive the crown of martyrdom.
\switchcolumn*
\selectlanguage{latin}
In Pérside natális beáti Polychrónii, Epíscopi Babylónis, qui, in 
 persecutióne Décii, ore lapídibus cæso, mánibus 
 exténsis, ad cælum óculos élevans, emísit spíritum.
\switchcolumn
\selectlanguage{english}
In Persia, during the persecution of Decius, the birthday of blessed 
 Polychronius, bishop of Babylon, who, being struck in the mouth with stones, 
 died with hands outstretched and eyes lifted towards heaven.
\switchcolumn*
\selectlanguage{latin}
Concórdiæ, in Venetórum fínibus, sanctórum 
 Mártyrum Donáti, Secundiáni et Rómuli, cum áliis octogínta sex, ejúsdem 
 corónæ consórtibus.
\switchcolumn
\selectlanguage{english}
At Concordia, the holy martyrs Donatus, Secundian, and Romulus, with 
 eighty-six others, partakers of the same crown.
\switchcolumn*
\selectlanguage{latin}
Cæsaréæ, in Palæstína, sancti Theodúli senis, 
 qui, cum esset ex família Prǽsidis Firmiliáni, et, Mártyrum excitátus 
 exémplo, Christum constánter confiterétur, martyrii palmam, cruci affíxus, 
 nóbili triúmpho proméruit.
\switchcolumn
\selectlanguage{english}
At Caesarea in Palestine, the death of St. Theodulus, in the service of the 
 governor Firmilian, at a great age. Prompted by the example of the 
 martyrs, he confessed Christ with constancy, and was nailed to a cross. 
 By this noble victory he merited the palm of martyrdom.
\switchcolumn*
\selectlanguage{latin}
Ibidem sancti Juliáni Cappádocis, qui, 
 exosculátus necatórum Mártyrum córpora, et ídeo ut Christiánus delátus et ad 
 Prǽsidem ductus, lento igne jussus est combúri.
\switchcolumn
\selectlanguage{english}
In the same place, St. Julian the Cappadocian, who, because he had kissed 
 the relics of martyrs, was denounced as a Christian. Being taken to 
 the governor, he was ordered to be burned to death over a slow fire.
\switchcolumn*
\selectlanguage{latin}
In pago Tarvanénsi, in Gállia, sancti Silvíni, Epíscopi Tolosáni.
\switchcolumn
\selectlanguage{english}
In the territory of Terouanne in France, St. Silvinus, bishop of Toulouse.
\switchcolumn*
\selectlanguage{latin}
In monastério Cluain-ednechénsi, in Hibérnia, sancti Fintáni, Presbyteri et 
 Abbátis.
\switchcolumn
\selectlanguage{english}
In the monastery of Cluainedhech in Ireland, St. Fintan, abbot.
\switchcolumn*
\selectlanguage{latin}
\end{paracol}


% ---- martyrology/mart02/mart0218.htm
\needspace{10\baselineskip}
\begin{paracol}{2}
\selectlanguage{latin}
\begin{center}{\color{gregoriocolor} Duodécimo Kaléndas Mártii. 
 Luna\dots\ }\end{center}
\switchcolumn
\selectlanguage{english}
\begin{center}{\color{gregoriocolor} The 18\textsuperscript{th} Day of 
 February. The\dots\ Day of the Moon.}\end{center}
\end{paracol}

\noindent\begin{tabularx}{\linewidth}{*{19}{>{\centering\arraybackslash}X}}
 \textcolor{gregoriocolor}{a} & \textcolor{gregoriocolor}{b} & \textcolor{gregoriocolor}{c} & \textcolor{gregoriocolor}{d} & \textcolor{gregoriocolor}{e} & \textcolor{gregoriocolor}{f} & \textcolor{gregoriocolor}{g} & \textcolor{gregoriocolor}{h} & \textcolor{gregoriocolor}{i} & \textcolor{gregoriocolor}{k} & \textcolor{gregoriocolor}{l} & \textcolor{gregoriocolor}{m} & \textcolor{gregoriocolor}{n} & \textcolor{gregoriocolor}{p} & \textcolor{gregoriocolor}{q} & \textcolor{gregoriocolor}{r} & \textcolor{gregoriocolor}{s} & \textcolor{gregoriocolor}{t} & \textcolor{gregoriocolor}{u} \\
 20 & 21 & 22 & 23 & 24 & 25 & 26 & 27 & 28 & 29 & 1 & 2 & 3 & 4 & 5 & 6 & 7 & 8 & 9 \\
\end{tabularx}
\vspace{0.5\baselineskip}
\noindent\begin{tabularx}{\linewidth}{*{12}{>{\centering\arraybackslash}X}}
 \textcolor{gregoriocolor}{A} & \textcolor{gregoriocolor}{B} & \textcolor{gregoriocolor}{C} & \textcolor{gregoriocolor}{D} & \textcolor{gregoriocolor}{E} & F & \textcolor{gregoriocolor}{F} & \textcolor{gregoriocolor}{G} & \textcolor{gregoriocolor}{H} & \textcolor{gregoriocolor}{M} & \textcolor{gregoriocolor}{N} & \textcolor{gregoriocolor}{P} \\
 10 & 11 & 12 & 13 & 14 & 15 & 14 & 15 & 16 & 17 & 18 & 19 \\
\end{tabularx}

\begin{paracol}{2}
\selectlanguage{latin}
\lettrine[lines=2]{H}{ierosólymis} natális sancti Simeónis, Epíscopi et Mártyris; qui fílius 
 Cléophæ et propínquus Salvatóris secúndum 
 carnem fuísse tráditur. Hic, Hierosolymórum Epíscopus post Jacóbum, 
 fratrem Dómini, ordinátus, et, in Trajáni persecutióne, multis supplíciis 
 afféctus, martyrio consummátus est, ómnibus qui áderant et Júdice ipso 
 mirántibus ut centum vigínti annórum senex fórtiter constantérque supplícium 
 crucis pertulísset.
\switchcolumn
\selectlanguage{english}
\lettrine[lines=2]{A}{t} Jerusalem, the birthday of St. Simeon, bishop and martyr, who is said to 
 have been the son of Cleophas, and a relative of the Saviour according to 
 the flesh. He was consecrated bishop of Jerusalem after St. James, the 
 cousin of our Lord. In the persecution of Trajan, after having endured 
 many torments, his martyrdom was completed. All who were present, even 
 the judge himself, were astonished that a man one hundred and twenty years 
 of age could bear the torment of crucifixion with such fortitude and 
 constancy.
\switchcolumn*
\selectlanguage{latin}
Apud Ostia Tiberína sanctórum Mártyrum Máximi et Cláudii fratrum, et Præpedígnæ, 
 uxóris Cláudii, cum duóbus fíliis Alexándro et Cútia; qui, cum essent 
 præclaríssimi géneris, omnes, jubénte Diocletiáno, tenti atque in exsílium 
 deportáti sunt, ac deínde, incéndio concremáti, Deo ipsi odoríferum martyrii 
 sacrifícium obtulérunt. Eórum relíquiæ, in flumen projéctæ et a 
 Christiánis perquisítæ, juxta eándem civitátem sepúltæ sunt.
\switchcolumn
\selectlanguage{english}
At Ostia, the holy martyrs Maximus and his brother Claudius, and Praepedigna, 
 the wife of Claudius, with her two sons Alexander and Cutias. Although 
 all of a noble birth, by the order of Diocletian, they were apprehended and 
 sent into exile. Afterwards being burned alive, they offered to God 
 the sweet sacrifice of martyrdom. Their remains were cast into the 
 river, but the Christians found them and buried them near the city.
\switchcolumn*
\selectlanguage{latin}
In Africa sanctórum Mártyrum Lúcii, Silváni, Rútuli, 
 Clássici, Secundíni, Frúctuli et Máximi.
\switchcolumn
\selectlanguage{english}
In Africa, the holy martyrs Lucius, Sylvanus, Rutulus, Classicus, Secundinus, 
 Fructulus, and Maximus.
\switchcolumn*
\selectlanguage{latin}
Constantinópoli sancti Flaviáni Epíscopi, qui, cum fidem cathólicam Ephesi 
 propugnáret, ab ímpii Dióscori factióne pugnis 
 et cálcibus percússus est, et, in exsílium actus, ibídem post tríduum vitam 
 finívit.
\switchcolumn
\selectlanguage{english}
At Constantinople, St. Flavian, bishop, who, for having defended the 
 Catholic faith at Ephesus, was attacked with slaps and kicks by the faction 
 of the impious Dioscorus, and then driven into exile where he died within 
 three days.
\switchcolumn*
\selectlanguage{latin}
Toléti, in Hispánia, sancti Helládii, Epíscopi et Confessóris, qui a sancto 
 Ildefónso, Toletáno Episcopo, multis láudibus celebrátur.
\switchcolumn
\selectlanguage{english}
At Toledo, Spain, St. Helladius, bishop and confessor, who received much 
 praise from St. Ildefonse, Bishop of Toledo.
\switchcolumn*
\selectlanguage{latin}
\end{paracol}


% ---- martyrology/mart02/mart0219.htm
\needspace{10\baselineskip}
\begin{paracol}{2}
\selectlanguage{latin}
\begin{center}{\color{gregoriocolor} Undécimo Kaléndas Mártii. 
 Luna\dots\ }\end{center}
\switchcolumn
\selectlanguage{english}
\begin{center}{\color{gregoriocolor} The 19\textsuperscript{th} Day of 
 February. The\dots\ Day of the Moon.}\end{center}
\end{paracol}

\noindent\begin{tabularx}{\linewidth}{*{19}{>{\centering\arraybackslash}X}}
 \textcolor{gregoriocolor}{a} & \textcolor{gregoriocolor}{b} & \textcolor{gregoriocolor}{c} & \textcolor{gregoriocolor}{d} & \textcolor{gregoriocolor}{e} & \textcolor{gregoriocolor}{f} & \textcolor{gregoriocolor}{g} & \textcolor{gregoriocolor}{h} & \textcolor{gregoriocolor}{i} & \textcolor{gregoriocolor}{k} & \textcolor{gregoriocolor}{l} & \textcolor{gregoriocolor}{m} & \textcolor{gregoriocolor}{n} & \textcolor{gregoriocolor}{p} & \textcolor{gregoriocolor}{q} & \textcolor{gregoriocolor}{r} & \textcolor{gregoriocolor}{s} & \textcolor{gregoriocolor}{t} & \textcolor{gregoriocolor}{u} \\
 21 & 22 & 23 & 24 & 25 & 26 & 27 & 28 & 29 & 1 & 2 & 3 & 4 & 5 & 6 & 7 & 8 & 9 & 10 \\
\end{tabularx}
\vspace{0.5\baselineskip}
\noindent\begin{tabularx}{\linewidth}{*{12}{>{\centering\arraybackslash}X}}
 \textcolor{gregoriocolor}{A} & \textcolor{gregoriocolor}{B} & \textcolor{gregoriocolor}{C} & \textcolor{gregoriocolor}{D} & \textcolor{gregoriocolor}{E} & F & \textcolor{gregoriocolor}{F} & \textcolor{gregoriocolor}{G} & \textcolor{gregoriocolor}{H} & \textcolor{gregoriocolor}{M} & \textcolor{gregoriocolor}{N} & \textcolor{gregoriocolor}{P} \\
 11 & 12 & 13 & 14 & 15 & 16 & 15 & 16 & 17 & 18 & 19 & 20 \\
\end{tabularx}

\begin{paracol}{2}
\selectlanguage{latin}
\lettrine[lines=2]{R}{omæ} natális sancti Gabíni, Presbyteri et 
 Mártyris, qui fuit frater beáti Caji Papæ, atque, a Diocletiáno diu in 
 custódia vínculis afflíctus, pretiósa morte sibi cæli gáudia comparávit.
\switchcolumn
\selectlanguage{english}
\lettrine[lines=2]{A}{t} Rome, the birthday of St. Gavinus, priest and martyr, brother of blessed 
 Pope Caius. After being chained in prison for a long time by 
 Diocletian, he obtained the joys of heaven by his esteemed death.
\switchcolumn*
\selectlanguage{latin}
In Africa sanctórum Mártyrum Públii, Juliáni, Marcélli et aliórum.
\switchcolumn
\selectlanguage{english}
In Africa, the holy martyrs Publius, Julian, Marcellus, and others.
\switchcolumn*
\selectlanguage{latin}
In Palæstína commemorátio sanctórum Monachórum, 
 et aliórum Mártyrum, qui a Saracénis, sub Duce Alamúndaro, ob Christi fidem, 
 sævíssime cæsi sunt.
\switchcolumn
\selectlanguage{english}
In Palestine, the commemoration of the holy monks and other martyrs who were 
 barbarously massacred for the faith of Christ by the Saracens, under their 
 leader Almondhar.
\switchcolumn*
\selectlanguage{latin}
Neápoli, in Campánia, sancti Quod-vult-Deus, Carthaginénsis Epíscopi, qui, 
 una cum Clero, a Rege Ariáno Genseríco in 
 fractas et absque remígiis ac velis naves impósitus, præter spem Neápolim áppulit, ibíque, in exsílio pósitus, Conféssor occúbuit.
\switchcolumn
\selectlanguage{english}
At Naples in Campania, St. Quodvultdeus, bishop of Carthage. The Arian 
 king Genseric placed him together with his clergy into boats which were 
 broken and without oars and sails, but they unexpectedly reached Naples. 
 He died in exile as a confessor.
\switchcolumn*
\selectlanguage{latin}
Hierosólymis sancti Zambdæ Epíscopi.
\switchcolumn
\selectlanguage{english}
At Jerusalem, St. Zambdas, bishop.
\switchcolumn*
\selectlanguage{latin}
Solis, in Cypro, sancti Auxíbii Epíscopi.
\switchcolumn
\selectlanguage{english}
At Soli in Cyprus, St. Auxibius, bishop.
\switchcolumn*
\selectlanguage{latin}
Apud Benevéntum sancti Barbáti Epíscopi, qui, sanctitáte célebris, 
 Longobárdos et eórum Ducem convértit ad Christum.
\switchcolumn
\selectlanguage{english}
At Benevento, St. Barbatus, a bishop illustrious for sanctity, who converted 
 the Lombards and their chief to the faith of Christ.
\switchcolumn*
\selectlanguage{latin}
Medioláni sancti Mansuéti, Epíscopi et Confessóris.
\switchcolumn
\selectlanguage{english}
At Milan, St. Mansuetus, bishop and confessor.
\switchcolumn*
\selectlanguage{latin}
\end{paracol}


% ---- martyrology/mart02/mart0220.htm
\needspace{10\baselineskip}
\begin{paracol}{2}
\selectlanguage{latin}
\begin{center}{\color{gregoriocolor} Décimo Kaléndas Mártii. 
 Luna\dots\ }\end{center}
\switchcolumn
\selectlanguage{english}
\begin{center}{\color{gregoriocolor} The 20\textsuperscript{th} Day of 
 February. The\dots\ Day of the Moon.}\end{center}
\end{paracol}

\noindent\begin{tabularx}{\linewidth}{*{19}{>{\centering\arraybackslash}X}}
 \textcolor{gregoriocolor}{a} & \textcolor{gregoriocolor}{b} & \textcolor{gregoriocolor}{c} & \textcolor{gregoriocolor}{d} & \textcolor{gregoriocolor}{e} & \textcolor{gregoriocolor}{f} & \textcolor{gregoriocolor}{g} & \textcolor{gregoriocolor}{h} & \textcolor{gregoriocolor}{i} & \textcolor{gregoriocolor}{k} & \textcolor{gregoriocolor}{l} & \textcolor{gregoriocolor}{m} & \textcolor{gregoriocolor}{n} & \textcolor{gregoriocolor}{p} & \textcolor{gregoriocolor}{q} & \textcolor{gregoriocolor}{r} & \textcolor{gregoriocolor}{s} & \textcolor{gregoriocolor}{t} & \textcolor{gregoriocolor}{u} \\
 22 & 23 & 24 & 25 & 26 & 27 & 28 & 29 & 1 & 2 & 3 & 4 & 5 & 6 & 7 & 8 & 9 & 10 & 11 \\
\end{tabularx}
\vspace{0.5\baselineskip}
\noindent\begin{tabularx}{\linewidth}{*{12}{>{\centering\arraybackslash}X}}
 \textcolor{gregoriocolor}{A} & \textcolor{gregoriocolor}{B} & \textcolor{gregoriocolor}{C} & \textcolor{gregoriocolor}{D} & \textcolor{gregoriocolor}{E} & F & \textcolor{gregoriocolor}{F} & \textcolor{gregoriocolor}{G} & \textcolor{gregoriocolor}{H} & \textcolor{gregoriocolor}{M} & \textcolor{gregoriocolor}{N} & \textcolor{gregoriocolor}{P} \\
 12 & 13 & 14 & 15 & 16 & 17 & 16 & 17 & 18 & 19 & 20 & 21 \\
\end{tabularx}

\begin{paracol}{2}
\selectlanguage{latin}
\lettrine[lines=2]{T}{yri,} in Phœnícia, commemorátio beatórum 
 Mártyrum, quorum númerum solíus sciéntia Dei cólligit. Hi omnes, sub 
 Diocletiáno Imperatóre, a Vetúrio, mílitum magístro, multis tormentórum 
 genéribus, sibi ínvicem succedéntibus, occísi sunt; nam, primo quidem 
 flagris toto córpore dilaniáti, inde divérsis bestiárum genéribus tráditi, 
 sed ab illis divína virtúte nil læsi, post, áddita feritáte ignis ac ferri, 
 martyrium consummárunt. Eórum vero gloriósam multitúdinem ad victóriam 
 incitábant Epíscopi Tyránnio, Silvánus, Péleus et Nilus, ac Presbyter 
 Zenóbius, qui, felíci agóne, una cum illis, martyrii palmam adépti sunt.
\switchcolumn
\selectlanguage{english}
\lettrine[lines=2]{A}{t} Tyre in Phoenicia, the commemoration of many blessed martyrs, whose 
 number is known to God alone. Under Emperor Diocletian, they were put 
 to death after a long and varied series of torments by the military 
 commander Veturius. They first had their bodies torn with scourges, 
 then delivered to several different kinds of beasts. Providence 
 prevented their injury throughout all this, but their martyrdom was granted 
 by means of fire and the sword. Tyrannio, Sylvanus, Peleus, and Nilus, 
 all bishops, and Zenobius, a priest, urged the gloriously assembled 
 multitude to victory, and they all endured the test successfully to win the 
 palm of martyrdom.
\switchcolumn*
\selectlanguage{latin}
Constantinópoli sancti Eleuthérii, Epíscopi et Mártyris.
\switchcolumn
\selectlanguage{english}
At Constantinople, St. Eleutherius, bishop and martyr.
\switchcolumn*
\selectlanguage{latin}
In Pérside natális sancti Sadoth Epíscopi, et aliórum centum vigínti octo; 
 qui, sub Rege Persárum Sápore, cum Solem 
 adoráre renuíssent, crudéli nece præcláras sibi corónas comparárunt.
\switchcolumn
\selectlanguage{english}
In Persia, in the time of King Sapor, the birthday of St. Sadoth, bishop, 
 and one hundred and twenty-eight others who refused to adore the sun, but 
 who by a cruel death purchased shining crowns.
\switchcolumn*
\selectlanguage{latin}
In Cypro sanctórum Mártyrum Potámii et Nemésii.
\switchcolumn
\selectlanguage{english}
In the island of Cyprus, the holy martyrs Pothamius and Nemesius.
\switchcolumn*
\selectlanguage{latin}
Cátanæ, in Sicília, sancti Leónis Epíscopi, qui 
 virtútibus atque miráculis coruscávit.
\switchcolumn
\selectlanguage{english}
At Catania in Sicily, St. Leo, bishop, illustrious for virtues and miracles.
\switchcolumn*
\selectlanguage{latin}
Eódem die sancti Euchérii, Aurelianénsis Epíscopi, qui eo magis miráculis 
 cláruit, pro plúribus invidórum calúmniis fuit oppréssus.
\switchcolumn
\selectlanguage{english}
The same day, St. Eucherius, bishop of Orleans, who, the more he was 
 oppressed by the calumnies of the envious, the more he impressed them with 
 his miracles.
\switchcolumn*
\selectlanguage{latin}
Tornáci, in Gálliis, sancti Eleuthérii, Epíscopi et Confessóris.
\switchcolumn
\selectlanguage{english}
At Tournai in Belgium, St. Eleutherius, bishop and confessor.
\switchcolumn*
\selectlanguage{latin}
\end{paracol}


% ---- martyrology/mart02/mart0221.htm
\needspace{10\baselineskip}
\begin{paracol}{2}
\selectlanguage{latin}
\begin{center}{\color{gregoriocolor} Nono Kaléndas Mártii. 
 Luna\dots\ }\end{center}
\switchcolumn
\selectlanguage{english}
\begin{center}{\color{gregoriocolor} The 21\textsuperscript{st} Day of 
 February. The\dots\ Day of the Moon.}\end{center}
\end{paracol}

\noindent\begin{tabularx}{\linewidth}{*{19}{>{\centering\arraybackslash}X}}
 \textcolor{gregoriocolor}{a} & \textcolor{gregoriocolor}{b} & \textcolor{gregoriocolor}{c} & \textcolor{gregoriocolor}{d} & \textcolor{gregoriocolor}{e} & \textcolor{gregoriocolor}{f} & \textcolor{gregoriocolor}{g} & \textcolor{gregoriocolor}{h} & \textcolor{gregoriocolor}{i} & \textcolor{gregoriocolor}{k} & \textcolor{gregoriocolor}{l} & \textcolor{gregoriocolor}{m} & \textcolor{gregoriocolor}{n} & \textcolor{gregoriocolor}{p} & \textcolor{gregoriocolor}{q} & \textcolor{gregoriocolor}{r} & \textcolor{gregoriocolor}{s} & \textcolor{gregoriocolor}{t} & \textcolor{gregoriocolor}{u} \\
 23 & 24 & 25 & 26 & 27 & 28 & 29 & 1 & 2 & 3 & 4 & 5 & 6 & 7 & 8 & 9 & 10 & 11 & 12 \\
\end{tabularx}
\vspace{0.5\baselineskip}
\noindent\begin{tabularx}{\linewidth}{*{12}{>{\centering\arraybackslash}X}}
 \textcolor{gregoriocolor}{A} & \textcolor{gregoriocolor}{B} & \textcolor{gregoriocolor}{C} & \textcolor{gregoriocolor}{D} & \textcolor{gregoriocolor}{E} & F & \textcolor{gregoriocolor}{F} & \textcolor{gregoriocolor}{G} & \textcolor{gregoriocolor}{H} & \textcolor{gregoriocolor}{M} & \textcolor{gregoriocolor}{N} & \textcolor{gregoriocolor}{P} \\
 13 & 14 & 15 & 16 & 17 & 18 & 17 & 18 & 19 & 20 & 21 & 22 \\
\end{tabularx}

\begin{paracol}{2}
\selectlanguage{latin}
\lettrine[lines=2]{S}{cythópoli,} in Palæstína, sancti Severiáni, 
 Epíscopi et Mártyris, qui, Eutychiánis acérrime se oppónens, gládio 
 perémptus est.
\switchcolumn
\selectlanguage{english}
\lettrine[lines=2]{A}{t} Scythopolis in Palestine, St. Severian, bishop and martyr, who was 
 beheaded by the Eutychians because he opposed them so zealously.
\switchcolumn*
\selectlanguage{latin}
In Sicília natális sanctórum Mártyrum septuagínta novem, qui sub Diocletiáno, 
 per divérsa torménta, confessiónis suæ corónam 
 percípere meruérunt.
\switchcolumn
\selectlanguage{english}
In Sicily, in the reign of Diocletian, the birthday of seventy-nine holy 
 martyrs, who, by reason of various tortures for their confession of faith, 
 deserved to receive an immortal crown.
\switchcolumn*
\selectlanguage{latin}
Adruméti, in Africa, sanctórum Mártyrum Véruli, 
 Secundíni, Sirícii, Felícis, Sérvuli, Saturníni, Fortunáti et aliórum 
 séxdecim, qui in persecutióne Wandálica, ob cathólicæ fídei confessiónem, 
 martyrio coronáti sunt.
\switchcolumn
\selectlanguage{english}
At Adrumetum in Africa, during the persecution of the Vandals, the holy 
 martyrs, Verulus, Secundinus, Siricius, Felix, Servulus, Saturninus, 
 Fortunatus, and sixteen others, who were crowned with martyrdom for 
 professing the Catholic faith.
\switchcolumn*
\selectlanguage{latin}
Damásci sancti Petri Maviméni, qui, cum díceret Arábibus quibúsdam, ad se
 ægrótum veniéntibus: « Omnis qui fidem 
 Christiánam cathólicam non ampléctitur, damnátus est, sicut et Máhumet, 
 pseudoprophéta vester », ab illis est necatus.
\switchcolumn
\selectlanguage{english}
At Damascus, St. Peter Mavimenus, who was killed by some Arabs who visited 
 him in his sickness, because he said to them: ``Whoever does not embrace the 
 Christian and Catholic faith is lost, like your false prophet Mohammed.''
\switchcolumn*
\selectlanguage{latin}
Metis, in Gállia, sancti Felícis Epíscopi.
\switchcolumn
\selectlanguage{english}
At Metz in France, St. Felix, bishop.
\switchcolumn*
\selectlanguage{latin}
Bríxiæ sancti Patérii Epíscopi.
\switchcolumn
\selectlanguage{english}
At Brescia, St. Paterius, bishop.
\switchcolumn*
\selectlanguage{latin}
\end{paracol}


% ---- martyrology/mart02/mart0222.htm
\needspace{10\baselineskip}
\begin{paracol}{2}
\selectlanguage{latin}
\begin{center}{\color{gregoriocolor} Octávo Kaléndas Mártii. 
 Luna\dots\ }\end{center}
\switchcolumn
\selectlanguage{english}
\begin{center}{\color{gregoriocolor} The 22\textsuperscript{nd} Day of 
 February. The\dots\ Day of the Moon.}\end{center}
\end{paracol}

\noindent\begin{tabularx}{\linewidth}{*{19}{>{\centering\arraybackslash}X}}
 \textcolor{gregoriocolor}{a} & \textcolor{gregoriocolor}{b} & \textcolor{gregoriocolor}{c} & \textcolor{gregoriocolor}{d} & \textcolor{gregoriocolor}{e} & \textcolor{gregoriocolor}{f} & \textcolor{gregoriocolor}{g} & \textcolor{gregoriocolor}{h} & \textcolor{gregoriocolor}{i} & \textcolor{gregoriocolor}{k} & \textcolor{gregoriocolor}{l} & \textcolor{gregoriocolor}{m} & \textcolor{gregoriocolor}{n} & \textcolor{gregoriocolor}{p} & \textcolor{gregoriocolor}{q} & \textcolor{gregoriocolor}{r} & \textcolor{gregoriocolor}{s} & \textcolor{gregoriocolor}{t} & \textcolor{gregoriocolor}{u} \\
 24 & 25 & 26 & 27 & 28 & 29 & 1 & 2 & 3 & 4 & 5 & 6 & 7 & 8 & 9 & 10 & 11 & 12 & 13 \\
\end{tabularx}
\vspace{0.5\baselineskip}
\noindent\begin{tabularx}{\linewidth}{*{12}{>{\centering\arraybackslash}X}}
 \textcolor{gregoriocolor}{A} & \textcolor{gregoriocolor}{B} & \textcolor{gregoriocolor}{C} & \textcolor{gregoriocolor}{D} & \textcolor{gregoriocolor}{E} & F & \textcolor{gregoriocolor}{F} & \textcolor{gregoriocolor}{G} & \textcolor{gregoriocolor}{H} & \textcolor{gregoriocolor}{M} & \textcolor{gregoriocolor}{N} & \textcolor{gregoriocolor}{P} \\
 14 & 15 & 16 & 17 & 18 & 19 & 18 & 19 & 20 & 21 & 22 & 23 \\
\end{tabularx}

\begin{paracol}{2}
\selectlanguage{latin}
\lettrine[lines=2]{A}{ntiochíæ} Cáthedra sancti Petri Apóstoli, ubi 
 primum discípuli cognomináti sunt Christiáni.
\switchcolumn
\selectlanguage{english}
\lettrine[lines=2]{T}{he} Chair of St. Peter at Antioch, where the disciples were first called 
 Christians.
\switchcolumn*
\selectlanguage{latin}
Favéntiæ, in Æmília, natális sancti Petri 
 Damiáni, Cardinális atque Epíscopi Ostiénsis et Confessóris, ex Ordine 
 Camaldulénsi, doctrína et sanctitáte célebris, quem Leo Papa Duodécimus 
 Doctórem universális Ecclésiæ declarávit. Ipsíus autem festum sequénti 
 die celebrátur.
\switchcolumn
\selectlanguage{english}
At Faenza in Emilia, the birthday of St. Peter Damian, cardinal bishop of 
 Ostia and confessor. He was a Camaldolese monk, famous for his 
 learning and sanctity, whom Pope Leo XII declared a doctor of the universal 
 Church. His feast is celebrated tomorrow.
\switchcolumn*
\selectlanguage{latin}
Salamínæ, in Cypro, sancti Aristiónis, qui (ut 
 mox memorándus Pápias testátur) fuit unus de septuagínta duóbus Christi 
 discípulis.
\switchcolumn
\selectlanguage{english}
At Salamis in Cyprus, St. Aristio, who (says Papias, the next to be 
 mentioned) was one of the seventy-two disciples of Christ.
\switchcolumn*
\selectlanguage{latin}
Hierápoli, in Phrygia, beáti Pápiæ, ejúsdem 
 civitátis Epíscopi, qui sancti Joánnis Senióris audítor, Polycárpi autem 
 sodális fuit.
\switchcolumn
\selectlanguage{english}
At Hierapolis in Phrygia, blessed Papias, bishop of that city, who was a 
 companion of Polycarp and a disciple of St. John.
\switchcolumn*
\selectlanguage{latin}
In Arábia commemorátio plurimórum sanctórum Mártyrum, qui, sub Galério 
 Maximiáno Imperatóre, sævíssime cæsi sunt.
\switchcolumn
\selectlanguage{english}
In Arabia, the commemoration of many holy martyrs who were barbarously put 
 to death under Emperor Galerius Maximian.
\switchcolumn*
\selectlanguage{latin}
Alexandríæ sancti Abílii Epíscopi, qui, 
 secúndus post beátum Marcum, factus ejúsdem civitátis Epíscopus, sacerdótium 
 virtúte conspícuus ministrávit.
\switchcolumn
\selectlanguage{english}
At Alexandria, St. Abilias, bishop, who was the second shepherd of that city 
 after St. Mark, and who administered his charge with eminent piety.
\switchcolumn*
\selectlanguage{latin}
Viénnæ, in Gállia, sancti Paschásii Epíscopi, 
 eruditióne et morum sanctitáte præclári.
\switchcolumn
\selectlanguage{english}
At Vienne in France, St. Paschasius, bishop, celebrated for his learning and 
 holy life.
\switchcolumn*
\selectlanguage{latin}
Ravénnæ sancti Maximiáni, Epíscopi et 
 Confessóris.
\switchcolumn
\selectlanguage{english}
At Ravenna, St. Maximian, bishop and confessor.
\switchcolumn*
\selectlanguage{latin}
Cortónæ, in Túscia, sanctæ Margarítæ, ex tértio 
 Ordine sancti Francísci; quæ admirábili pæniténtia et ubérrimis lácrimis 
 máculas anteáctæ vitæ indesinénter abstérsit. Ipsíus corpus, 
 mirabíliter incorrúptum, suávem spirans odórem et crebris miráculis clarum, 
 ibídem magno cum honóre cólitur.
\switchcolumn
\selectlanguage{english}
At Cortona in Tuscany, St. Margaret of the Third Order of St. Francis. 
 By means of commendable penance and fruitful tears, she wiped away the 
 stains of her previous life. Her body miraculously remained incorrupt 
 for more than four centuries, giving forth a sweet odour, and producing 
 frequent miracles. It is honoured in that place with great devotion.
\switchcolumn*
\selectlanguage{latin}
\end{paracol}


% ---- martyrology/mart02/mart0223.htm
\needspace{10\baselineskip}
\begin{paracol}{2}
\selectlanguage{latin}
\begin{center}{\color{gregoriocolor} Séptimo Kaléndas Mártii. 
 Luna\dots\ }\end{center}
\switchcolumn
\selectlanguage{english}
\begin{center}{\color{gregoriocolor} The 23\textsuperscript{rd} Day of 
 February. The\dots\ Day of the Moon.}\end{center}
\end{paracol}

\noindent\begin{tabularx}{\linewidth}{*{19}{>{\centering\arraybackslash}X}}
 \textcolor{gregoriocolor}{a} & \textcolor{gregoriocolor}{b} & \textcolor{gregoriocolor}{c} & \textcolor{gregoriocolor}{d} & \textcolor{gregoriocolor}{e} & \textcolor{gregoriocolor}{f} & \textcolor{gregoriocolor}{g} & \textcolor{gregoriocolor}{h} & \textcolor{gregoriocolor}{i} & \textcolor{gregoriocolor}{k} & \textcolor{gregoriocolor}{l} & \textcolor{gregoriocolor}{m} & \textcolor{gregoriocolor}{n} & \textcolor{gregoriocolor}{p} & \textcolor{gregoriocolor}{q} & \textcolor{gregoriocolor}{r} & \textcolor{gregoriocolor}{s} & \textcolor{gregoriocolor}{t} & \textcolor{gregoriocolor}{u} \\
 25 & 26 & 27 & 28 & 29 & 1 & 2 & 3 & 4 & 5 & 6 & 7 & 8 & 9 & 10 & 11 & 12 & 13 & 14 \\
\end{tabularx}
\vspace{0.5\baselineskip}
\noindent\begin{tabularx}{\linewidth}{*{12}{>{\centering\arraybackslash}X}}
 \textcolor{gregoriocolor}{A} & \textcolor{gregoriocolor}{B} & \textcolor{gregoriocolor}{C} & \textcolor{gregoriocolor}{D} & \textcolor{gregoriocolor}{E} & F & \textcolor{gregoriocolor}{F} & \textcolor{gregoriocolor}{G} & \textcolor{gregoriocolor}{H} & \textcolor{gregoriocolor}{M} & \textcolor{gregoriocolor}{N} & \textcolor{gregoriocolor}{P} \\
 15 & 16 & 17 & 18 & 19 & 20 & 19 & 20 & 21 & 22 & 23 & 24 \\
\end{tabularx}

\begin{paracol}{2}
\selectlanguage{latin}
\lettrine[lines=1]{V}{igília} sancti Matthíæ Apóstoli.
\switchcolumn
\selectlanguage{english}
\lettrine[lines=1]{T}{he} Vigil of St. Matthias the Apostle.
\switchcolumn*
\selectlanguage{latin}
Sancti Petri Damiáni, ex Ordine Camaldulénsi, Cardinális et Epíscopi 
 Ostiénsis, Confessóris et Ecclésiæ Doctóris, 
 qui evolávit in cælum prídie hujus diéi.
\switchcolumn
\selectlanguage{english}
St. Peter Damian, a Camaldolese monk, cardinal bishop of Ostia, confessor 
 and doctor of the Church, who died on the 22nd of February.
\switchcolumn*
\selectlanguage{latin}
Smyrnæ natális sancti Polycárpi, qui, beáti 
 Joánnis Apóstoli discípulus, et ab eo ejúsdem civitátis Epíscopus ordinátus, 
 totíus Asiæ Princeps fuit. Póstea, sub Marco Antoníno et Lúcio Aurélio 
 Cómmodo, sedénte Procónsule et univérso pópulo in theátro advérsus eum 
 personánte, igni tráditus est; et, cum ab igne mínime læderétur, martyrii 
 corónam, gládio confóssus, accépit. Cum illo étiam álii duódecim, qui 
 ex Philadelphía vénerant, in eádem Smyrnénsi urbe, martyrio consummáti sunt. 
 Ipsíus tamen Polycárpi festum séptimo Kaléndas Februárii celebrátur.
\switchcolumn
\selectlanguage{english}
At Smyrna, the birthday of St. Polycarp, a disciple of St. John the Apostle, 
 by whom he was consecrated bishop of that city, and appointed primate of all 
 Asia. Under Marcus Antonius and Lucius Aurelius Commodus, when the 
 proconsul and all those assembled in the amphitheatre cried out against him, 
 he was delivered to the fire, but since it did not harm him, he received the 
 crown of martyrdom by the sword. With him, twelve others who came from 
 Philadelphia met their death by martyrdom in the same city. The feast 
 of St. Polycarp is kept on the 26th of January.
\switchcolumn*
\selectlanguage{latin}
Apud Sírmium beáti Siréni, Mónachi et Mártyris, qui, jubénte Maximiáno 
 Imperatóre, reténtus est, et, cum se Christiánum esse confiterétur, cápite 
 obtruncátus.
\switchcolumn
\selectlanguage{english}
At Sirmio, blessed Sirenus, monk and martyr. He was arrested by order 
 of Emperor Maximian and beheaded for confessing that he was a Christian.
\switchcolumn*
\selectlanguage{latin}
Ibídem natális sanctórum septuagínta duórum Mártyrum, qui, martyrii certámen 
 in præfáta urbe consummántes, mansúra 
 percepérunt regna.
\switchcolumn
\selectlanguage{english}
In the same place, the birthday of seventy-two holy martyrs, who suffered 
 martyrdom in the same city and who took possession of the everlasting 
 kingdom.
\switchcolumn*
\selectlanguage{latin}
In civitáte Asturicénsi, in Hispánia, sanctæ 
 Marthæ, Vírginis et Mártyris, quæ, sub Décio Imperatóre et Patérno 
 Procónsule, dire ob Christi fidem est cruciáta et gládio tandem occísa.
\switchcolumn
\selectlanguage{english}
In the city of Astorga in Spain, St. Martha, virgin and martyr, under 
 Emperor Decius and the proconsul Paternus. She was cruelly tortured 
 for the faith of Christ and was finally slain by the sword.
\switchcolumn*
\selectlanguage{latin}
Constantinópoli sancti Lázari Mónachi, qui, cum 
 sacras Imágines píngeret, idcírco, Imperatóris Iconoclástæ Theóphili jussu, 
 diris suppliciis excruciátur, et ei manus candénti ferro combúritur; sed, 
 Dei virtúte sanátus, abrásas Imágines sanctas pingéndo restítuit, ac demum 
 in pace quiévit.
\switchcolumn
\selectlanguage{english}
At Constantinople, St. Lazarus, monk. The Iconoclast emperor 
 Theophilus commanded him to be tortured with severe punishments because he 
 had painted some sacred pictures. His hand was burned with a hot iron, 
 but it was healed by the power of God, after which he repainted the holy 
 pictures that had been destroyed. He ended his life in peace.
\switchcolumn*
\selectlanguage{latin}
Bríxiæ sancti Felícis Epíscopi.
\switchcolumn
\selectlanguage{english}
At Brescia, St. Felix, bishop.
\switchcolumn*
\selectlanguage{latin}
Romæ sancti Polycárpi Presbyteri, qui, cum 
 beáto Sebastiáno, plúrimos ad Christi fidem convértit, atque ad martyrii 
 glóriam exhortándo perdúxit.
\switchcolumn
\selectlanguage{english}
At Rome, St. Polycarp, priest, who with blessed Sebastian converted many to 
 the faith of Christ, and by his exhortation led them to the glory of 
 martyrdom.
\switchcolumn*
\selectlanguage{latin}
Híspali, in Hispánia, sancti Floréntii 
 Confessóris.
\switchcolumn
\selectlanguage{english}
At Seville in Spain, St. Florentius, confessor.
\switchcolumn*
\selectlanguage{latin}
Tudérti, in Umbria, sanctæ Románæ Vírginis, quæ, 
 a sancto Silvéstro Papa baptizáta, in antris et spelúncis cæléstem vitam 
 duxit, et miraculórum glória cláruit.
\switchcolumn
\selectlanguage{english}
At Todi in Umbria, St. Romana, virgin, who was baptized by Pope St. 
 Sylvester, led a life of holiness in dens and caves, and wrought glorious 
 miracles.
\switchcolumn*
\selectlanguage{latin}
In Anglia sanctæ Milbúrgis Vírginis, fíliæ 
 Regis Merciórum.
\switchcolumn
\selectlanguage{english}
In England, St. Milburga, virgin, the daughter of the king of Mercia.
\switchcolumn*
\selectlanguage{latin}
\end{paracol}


% ---- martyrology/mart02/mart0223leap.htm
\needspace{10\baselineskip}
\begin{paracol}{2}
\selectlanguage{latin}
\begin{center}{\color{gregoriocolor} Séptimo Kaléndas Mártii. 
 Luna\dots\ }\end{center}
\switchcolumn
\selectlanguage{english}
\begin{center}{\color{gregoriocolor} The 23\textsuperscript{rd} Day of 
 February. The\dots\ Day of the Moon.}\end{center}
\end{paracol}

\noindent\begin{tabularx}{\linewidth}{*{19}{>{\centering\arraybackslash}X}}
 \textcolor{gregoriocolor}{a} & \textcolor{gregoriocolor}{b} & \textcolor{gregoriocolor}{c} & \textcolor{gregoriocolor}{d} & \textcolor{gregoriocolor}{e} & \textcolor{gregoriocolor}{f} & \textcolor{gregoriocolor}{g} & \textcolor{gregoriocolor}{h} & \textcolor{gregoriocolor}{i} & \textcolor{gregoriocolor}{k} & \textcolor{gregoriocolor}{l} & \textcolor{gregoriocolor}{m} & \textcolor{gregoriocolor}{n} & \textcolor{gregoriocolor}{p} & \textcolor{gregoriocolor}{q} & \textcolor{gregoriocolor}{r} & \textcolor{gregoriocolor}{s} & \textcolor{gregoriocolor}{t} & \textcolor{gregoriocolor}{u} \\
 25 & 26 & 27 & 28 & 29 & 1 & 2 & 3 & 4 & 5 & 6 & 7 & 8 & 9 & 10 & 11 & 12 & 13 & 14 \\
\end{tabularx}
\vspace{0.5\baselineskip}
\noindent\begin{tabularx}{\linewidth}{*{12}{>{\centering\arraybackslash}X}}
 \textcolor{gregoriocolor}{A} & \textcolor{gregoriocolor}{B} & \textcolor{gregoriocolor}{C} & \textcolor{gregoriocolor}{D} & \textcolor{gregoriocolor}{E} & F & \textcolor{gregoriocolor}{F} & \textcolor{gregoriocolor}{G} & \textcolor{gregoriocolor}{H} & \textcolor{gregoriocolor}{M} & \textcolor{gregoriocolor}{N} & \textcolor{gregoriocolor}{P} \\
 15 & 16 & 17 & 18 & 19 & 20 & 19 & 20 & 21 & 22 & 23 & 24 \\
\end{tabularx}

\begin{paracol}{2}
\selectlanguage{latin}
\lettrine[lines=2]{S}{ancti} Petri Damiáni, ex Ordine Camaldulénsi, Cardinális et Epíscopi 
 Ostiénsis, Confessóris et Ecclésiæ Doctoris, 
 qui evolávit in cælum prídie hujus diéi.
\switchcolumn
\selectlanguage{english}
\lettrine[lines=2]{S}{t.} Peter Damian, a Camaldolese monk, cardinal bishop of Ostia, confessor 
 and doctor of the Church, who died on the 22nd of February.
\switchcolumn*
\selectlanguage{latin}
Smyrnæ natális sancti Polycárpi, qui, beáti 
 Joánnis Apóstoli discípulus, et ab eo ejúsdem civitátis Epíscopus ordinátus, 
 totíus Asiæ Princeps fuit. Póstea, sub Marco Antoníno et Lúcio Aurélio 
 Cómmodo, sedénte Procónsule et univérso pópulo in theátro advérsus eum 
 personánte, igni tráditus est; et, cum ab igne mínime læderétur, martyrii 
 corónam, gládio confóssus, accépit. Cum illo étiam álii duódecim, qui 
 ex Philadelphía vénerant, in eádem Smyrnénsi urbe, martyrio consummáti sunt. 
 Ipsíus tamen Polycárpi festum séptimo Kaléndas Februárii celebrátur.
\switchcolumn
\selectlanguage{english}
At Smyrna, the birthday of St. Polycarp, a disciple of St. John the Apostle, 
 by whom he was consecrated bishop of that city, and appointed primate of all 
 Asia. Under Marcus Antonius and Lucius Aurelius Commodus, when the 
 proconsul and all those assembled in the amphitheatre cried out against him, 
 he was delivered to the fire, but since it did not harm him, he received the 
 crown of martyrdom by the sword. With him, twelve others who came from 
 Philadelphia met their death by martyrdom in the same city. The feast 
 of St. Polycarp is kept on the 26th of January.
\switchcolumn*
\selectlanguage{latin}
Apud Sírmium beáti Siréni, Mónachi et Mártyris, qui, jubénte Maximiáno 
 Imperatóre, reténtus est, et, cum se Christiánum esse confiterétur, cápite 
 obtruncátus.
\switchcolumn
\selectlanguage{english}
At Sirmio, blessed Sirenus, monk and martyr. He was arrested by order 
 of Emperor Maximian and beheaded for confessing that he was a Christian.
\switchcolumn*
\selectlanguage{latin}
Ibídem natális sanctórum septuagínta duórum Mártyrum, qui, martyrii certámen 
 in præfáta urbe consummántes, mansúra 
 percepérunt regna.
\switchcolumn
\selectlanguage{english}
In the same place, the birthday of seventy-two holy martyrs, who suffered 
 martyrdom in the same city and who took possession of the everlasting 
 kingdom.
\switchcolumn*
\selectlanguage{latin}
In civitáte Asturicénsi, in Hispánia, sanctæ 
 Marthæ, Vírginis et Mártyris, quæ, sub Décio Imperatóre et Patérno 
 Procónsule, dire ob Christi fidem est cruciáta et gládio tandem occísa.
\switchcolumn
\selectlanguage{english}
In the city of Astorga in Spain, St. Martha, virgin and martyr, under 
 Emperor Decius and the proconsul Paternus. She was cruelly tortured 
 for the faith of Christ and was finally slain by the sword.
\switchcolumn*
\selectlanguage{latin}
Constantinópoli sancti Lázari Mónachi, qui, cum 
 sacras Imágines píngeret, idcírco, Imperatóris Iconoclástæ Theóphili jussu, 
 diris supplíciis excruciátur, et ei manus candénti ferro combúritur; sed, 
 Dei virtúte sanátus, abrásas Imágines sanctas pingéndo restítuit, ac demum 
 in pace quiévit.
\switchcolumn
\selectlanguage{english}
At Constantinople, St. Lazarus, monk. The Iconoclast emperor 
 Theophilus commanded him to be tortured with severe punishments because he 
 had painted some sacred pictures. His hand was burned with a hot iron, 
 but it was healed by the power of God, after which he repainted the holy 
 pictures that had been destroyed. He ended his life in peace.
\switchcolumn*
\selectlanguage{latin}
Bríxiæ sancti Felícis Epíscopi.
\switchcolumn
\selectlanguage{english}
At Brescia, St. Felix, bishop.
\switchcolumn*
\selectlanguage{latin}
Romæ sancti Polycárpi Presbyteri, qui, cum 
 beáto Sebastiáno, plúrimos ad Christi fidem convértit, atque ad martyrii 
 glóriam exhortándo perdúxit.
\switchcolumn
\selectlanguage{english}
At Rome, St. Polycarp, priest, who with blessed Sebastian converted many to 
 the faith of Christ, and by his exhortation led them to the glory of 
 martyrdom.
\switchcolumn*
\selectlanguage{latin}
Híspali, in Hispánia, sancti Floréntii 
 Confessóris.
\switchcolumn
\selectlanguage{english}
At Seville in Spain, St. Florentius, confessor.
\switchcolumn*
\selectlanguage{latin}
Tudérti, in Umbria, sanctæ Románæ Vírginis, quæ, 
 a sancto Silvéstro Papa baptizáta, in antris et spelúncis cæléstem vitam 
 duxit, et miraculórum glória cláruit.
\switchcolumn
\selectlanguage{english}
At Todi in Umbria, St. Romana, virgin, who was baptized by Pope St. 
 Sylvester, led a life of holiness in dens and caves, and wrought glorious 
 miracles.
\switchcolumn*
\selectlanguage{latin}
In Anglia sanctæ Milbúrgis Vírginis, fíliæ 
 Regis Merciórum.
\switchcolumn
\selectlanguage{english}
In England, St. Milburga, virgin, the daughter of the king of Mercia.
\switchcolumn*
\selectlanguage{latin}
\end{paracol}


% ---- martyrology/mart02/mart0224.htm
\needspace{10\baselineskip}
\begin{paracol}{2}
\selectlanguage{latin}
\begin{center}{\color{gregoriocolor} Sexto Kaléndas Mártii. 
 Luna\dots\ }\end{center}
\switchcolumn
\selectlanguage{english}
\begin{center}{\color{gregoriocolor} The 24\textsuperscript{th} Day of February. The\dots\ Day of the Moon.}\end{center}
\end{paracol}

\noindent\begin{tabularx}{\linewidth}{*{19}{>{\centering\arraybackslash}X}}
 \textcolor{gregoriocolor}{a} & \textcolor{gregoriocolor}{b} & \textcolor{gregoriocolor}{c} & \textcolor{gregoriocolor}{d} & \textcolor{gregoriocolor}{e} & \textcolor{gregoriocolor}{f} & \textcolor{gregoriocolor}{g} & \textcolor{gregoriocolor}{h} & \textcolor{gregoriocolor}{i} & \textcolor{gregoriocolor}{k} & \textcolor{gregoriocolor}{l} & \textcolor{gregoriocolor}{m} & \textcolor{gregoriocolor}{n} & \textcolor{gregoriocolor}{p} & \textcolor{gregoriocolor}{q} & \textcolor{gregoriocolor}{r} & \textcolor{gregoriocolor}{s} & \textcolor{gregoriocolor}{t} & \textcolor{gregoriocolor}{u} \\
 26 & 27 & 28 & 29 & 1 & 2 & 3 & 4 & 5 & 6 & 7 & 8 & 9 & 10 & 11 & 12 & 13 & 14 & 15 \\
\end{tabularx}
\vspace{0.5\baselineskip}
\noindent\begin{tabularx}{\linewidth}{*{12}{>{\centering\arraybackslash}X}}
 \textcolor{gregoriocolor}{A} & \textcolor{gregoriocolor}{B} & \textcolor{gregoriocolor}{C} & \textcolor{gregoriocolor}{D} & \textcolor{gregoriocolor}{E} & F & \textcolor{gregoriocolor}{F} & \textcolor{gregoriocolor}{G} & \textcolor{gregoriocolor}{H} & \textcolor{gregoriocolor}{M} & \textcolor{gregoriocolor}{N} & \textcolor{gregoriocolor}{P} \\
 16 & 17 & 18 & 19 & 20 & 21 & 20 & 21 & 22 & 23 & 24 & 25 \\
\end{tabularx}

\begin{paracol}{2}
\selectlanguage{latin}
\lettrine[lines=2]{I}{n} Judæa natális sancti Matthíæ Apóstoli, qui, 
 post Ascensiónem Dómini ab Apóstolis in Judæ proditóris locum sorte eléctus, 
 pro Evangélii prædicatióne martyrium passus est.
\switchcolumn
\selectlanguage{english}
\lettrine[lines=2]{I}{n} Judea, the birthday of St. Matthias the Apostle. After the 
 Ascension of our Lord, the Apostles chose him, by lot, to fill the place of 
 Judas the traitor, and he suffered martyrdom for the preaching of the 
 Gospel.
\switchcolumn*
\selectlanguage{latin}
Romæ sanctæ Primitívæ Mártyris.
\switchcolumn
\selectlanguage{english}
At Rome, St. Primitiva, martyr.
\switchcolumn*
\selectlanguage{latin}
Rotómagi pássio sancti Prætextáti, Epíscopi et 
 Mártyris.
\switchcolumn
\selectlanguage{english}
At Rouen, the passion of St. Praetextatus, bishop and martyr.
\switchcolumn*
\selectlanguage{latin}
Cæsaréæ, in Cappadócia, sancti Sérgii Mártyris, 
 cujus gesta præclára habéntur.
\switchcolumn
\selectlanguage{english}
At Caesarea in Cappadocia, St. Sergius, martyr, of whose life a beautiful 
 account still exists.
\switchcolumn*
\selectlanguage{latin}
In Africa sanctórum Mártyrum Montáni, Lúcii, Juliáni, Victórici, 
 Flaviáni et Sociórum, qui discípuli fuérunt sancti Cypriáni, et, sub 
 Valeriáno Imperatóre, martyrium consummárunt.
\switchcolumn
\selectlanguage{english}
In Africa, the holy martyrs Montanus, Lucius, Julian, Victoricus, Flavian, 
 and their companions. They were disciples of St. Cyprian and suffered 
 martyrdom under Emperor Valerian.
\switchcolumn*
\selectlanguage{latin}
Tréviris sancti Modésti, Epíscopi et 
 Confessóris.
\switchcolumn
\selectlanguage{english}
At Treves, St. Modestus, bishop and confessor.
\switchcolumn*
\selectlanguage{latin}
Apud Stylum, in Calábria, sancti Joánnis, cognoménto Therísti, monásticæ 
 vitæ laude, et sanctitáte insígnis.
\switchcolumn
\selectlanguage{english}
At Stylo in Calabria, St. John Therestus, noted for his sanctity, and his 
 high regard for the monastic life.
\switchcolumn*
\selectlanguage{latin}
In Anglia sancti Edilbérti, Regis Cantiórum, quem sanctus Augustínus, 
 Anglórum Epíscopus, ad Christi fidem convértit.
\switchcolumn
\selectlanguage{english}
In England, St. Ethelbert, ruler of Kent, converted to the faith of Christ 
 by the English bishop, St. Augustine.
\switchcolumn*
\selectlanguage{latin}
Hierosólymis prima Invéntio cápitis sancti Joánnis Baptístæ, 
 Præcursóris Domini.
\switchcolumn
\selectlanguage{english}
At Jerusalem, the finding for the first time of the head of St. John the 
 Baptist, Precursor of the Lord.
\switchcolumn*
\selectlanguage{latin}
\end{paracol}


% ---- martyrology/mart02/mart0224leap.htm
\needspace{10\baselineskip}
\begin{paracol}{2}
\selectlanguage{latin}
\begin{center}{\color{gregoriocolor} Sexto Kaléndas Mártii. 
 Luna\dots\ }\end{center}
\switchcolumn
\selectlanguage{english}
\begin{center}{\color{gregoriocolor} The 24\textsuperscript{th} Day of February. The\dots\ Day of the Moon.}\end{center}
\end{paracol}

\noindent\begin{tabularx}{\linewidth}{*{19}{>{\centering\arraybackslash}X}}
 \textcolor{gregoriocolor}{a} & \textcolor{gregoriocolor}{b} & \textcolor{gregoriocolor}{c} & \textcolor{gregoriocolor}{d} & \textcolor{gregoriocolor}{e} & \textcolor{gregoriocolor}{f} & \textcolor{gregoriocolor}{g} & \textcolor{gregoriocolor}{h} & \textcolor{gregoriocolor}{i} & \textcolor{gregoriocolor}{k} & \textcolor{gregoriocolor}{l} & \textcolor{gregoriocolor}{m} & \textcolor{gregoriocolor}{n} & \textcolor{gregoriocolor}{p} & \textcolor{gregoriocolor}{q} & \textcolor{gregoriocolor}{r} & \textcolor{gregoriocolor}{s} & \textcolor{gregoriocolor}{t} & \textcolor{gregoriocolor}{u} \\
 26 & 27 & 28 & 29 & 1 & 2 & 3 & 4 & 5 & 6 & 7 & 8 & 9 & 10 & 11 & 12 & 13 & 14 & 15 \\
\end{tabularx}
\vspace{0.5\baselineskip}
\noindent\begin{tabularx}{\linewidth}{*{12}{>{\centering\arraybackslash}X}}
 \textcolor{gregoriocolor}{A} & \textcolor{gregoriocolor}{B} & \textcolor{gregoriocolor}{C} & \textcolor{gregoriocolor}{D} & \textcolor{gregoriocolor}{E} & F & \textcolor{gregoriocolor}{F} & \textcolor{gregoriocolor}{G} & \textcolor{gregoriocolor}{H} & \textcolor{gregoriocolor}{M} & \textcolor{gregoriocolor}{N} & \textcolor{gregoriocolor}{P} \\
 16 & 17 & 18 & 19 & 20 & 21 & 20 & 21 & 22 & 23 & 24 & 25 \\
\end{tabularx}

% NOTE: this is a special case, so we keep the conclusion and response
\begin{paracol}{2}
\selectlanguage{latin}
\lettrine[lines=2]{V}{igília} sancti Matthíæ 
Apóstoli.
\switchcolumn
\selectlanguage{english}
\lettrine[lines=2]{T}{he} Vigil of St. Matthias the Apostle.
\switchcolumn*
\selectlanguage{latin}
Item commemorátio plurimórum sanctórum Mártyrum et Confessórum, atque sanctárum Vírginum. \textcolor{gregoriocolor}{\Rbar.}~Deo grátias.
\switchcolumn
\selectlanguage{english}
Also the commemoration of many other holy martyrs, confessors, and holy virgins. \textcolor{gregoriocolor}{\Rbar.}~Thanks be to God.
\switchcolumn*
\selectlanguage{latin}
\end{paracol}


% ---- martyrology/mart02/mart0225.htm
\needspace{10\baselineskip}
\begin{paracol}{2}
\selectlanguage{latin}
\begin{center}{\color{gregoriocolor} Quinto Kaléndas Mártii. 
 Luna\dots\ }\end{center}
\switchcolumn
\selectlanguage{english}
\begin{center}{\color{gregoriocolor} The 25\textsuperscript{th} Day of February. The\dots\ Day of the Moon.}\end{center}
\end{paracol}

\noindent\begin{tabularx}{\linewidth}{*{19}{>{\centering\arraybackslash}X}}
 \textcolor{gregoriocolor}{a} & \textcolor{gregoriocolor}{b} & \textcolor{gregoriocolor}{c} & \textcolor{gregoriocolor}{d} & \textcolor{gregoriocolor}{e} & \textcolor{gregoriocolor}{f} & \textcolor{gregoriocolor}{g} & \textcolor{gregoriocolor}{h} & \textcolor{gregoriocolor}{i} & \textcolor{gregoriocolor}{k} & \textcolor{gregoriocolor}{l} & \textcolor{gregoriocolor}{m} & \textcolor{gregoriocolor}{n} & \textcolor{gregoriocolor}{p} & \textcolor{gregoriocolor}{q} & \textcolor{gregoriocolor}{r} & \textcolor{gregoriocolor}{s} & \textcolor{gregoriocolor}{t} & \textcolor{gregoriocolor}{u} \\
 27 & 28 & 29 & 1 & 2 & 3 & 4 & 5 & 6 & 7 & 8 & 9 & 10 & 11 & 12 & 13 & 14 & 15 & 16 \\
\end{tabularx}
\vspace{0.5\baselineskip}
\noindent\begin{tabularx}{\linewidth}{*{12}{>{\centering\arraybackslash}X}}
 \textcolor{gregoriocolor}{A} & \textcolor{gregoriocolor}{B} & \textcolor{gregoriocolor}{C} & \textcolor{gregoriocolor}{D} & \textcolor{gregoriocolor}{E} & F & \textcolor{gregoriocolor}{F} & \textcolor{gregoriocolor}{G} & \textcolor{gregoriocolor}{H} & \textcolor{gregoriocolor}{M} & \textcolor{gregoriocolor}{N} & \textcolor{gregoriocolor}{P} \\
 17 & 18 & 19 & 20 & 21 & 22 & 21 & 22 & 23 & 24 & 25 & 26 \\
\end{tabularx}

\begin{paracol}{2}
\selectlanguage{latin}
\lettrine[lines=2]{I}{n} Ægypto natális sanctórum Mártyrum Victoríni, 
 Victóris, Nicéphori, Claudiáni, Dióscori, Serapiónis et Pápiæ, sub Numeriáno 
 Imperatóre. Horum duo primi, pro confessióne fídei, exquisíta 
 suppliciórum génera constánter passi, cápite plectúntur; Nicéphorus, post 
 cratículas candéntes ignésque superátos, minutátim concísus est; Claudiánus 
 et Dióscorus flammis incénsi; Serápion vero et Pápias gládio cæsi sunt.
\switchcolumn
\selectlanguage{english}
\lettrine[lines=2]{I}{n} Egypt, under Emperor Numerian, the birthday of the holy martyrs 
 Victorinus, Victor, Nicephorus, Claudian, Dioscorus, Serapion, and Papias. 
 After patiently enduring extreme tortures, the first two were beheaded for 
 the confession of the faith, Nicephorus was laid on a heated gridiron, 
 placed over the fire, then thoroughly hacked with a knife; Claudian and 
 Dioscorus were burned at the stake; Serapion and Papias were slain with the 
 sword.
\switchcolumn*
\selectlanguage{latin}
In Africa sanctórum Mártyrum Donáti, Justi, Herénæ 
 et Sociórum.
\switchcolumn
\selectlanguage{english}
In Africa, the holy martyrs Donatus, Justus, Herenas, and their companions.
\switchcolumn*
\selectlanguage{latin}
Constantinópoli sancti Tharásii Epíscopi, eruditióne et pietáte insígnis; ad 
 quem exstat Hadriáni Papæ Primi epístola pro 
 defensióne sanctárum Imáginum.
\switchcolumn
\selectlanguage{english}
At Constantinople, St. Tharasius, bishop, a man of great learning and piety. 
 There exists a letter defending sacred images, written to him by Pope 
 Hadrian I.
\switchcolumn*
\selectlanguage{latin}
Naziánzi, in Cappadócia, sancti Cæsárii, qui 
 beátæ Nonnæ fílius ac beatórum Gregórii Theólogi et Gorgóniæ fuit frater, et 
 quem idem Gregórius inter ágmina beatórum se vidísse testátur.
\switchcolumn
\selectlanguage{english}
At Nazianzus, St. Caesarius, who was the son of blessed Nonna, and whom his 
 brother, blessed Gregory the Theologian, says he saw among the hosts of the 
 blessed.
\switchcolumn*
\selectlanguage{latin}
In monastério Heidenhémii, diœcésis 
 Eistetténsis, in Germánia, sanctæ Walbúrgæ Vírginis, quæ fuit fília sancti 
 Richárdi, Anglórum Regis, et soror sancti Willebáldi, Eistetténsis Epíscopi.
\switchcolumn
\selectlanguage{english}
In the monastery of Heidenheim, in the Eichstadt diocese in Germany, St. 
 Walburga, virgin. She was the daughter of St. Richard, king of 
 England, and sister of St. Willebald, bishop of Eichstadt.
\switchcolumn*
\selectlanguage{latin}
\end{paracol}


% ---- martyrology/mart02/mart0225leap.htm
\needspace{10\baselineskip}
\begin{paracol}{2}
\selectlanguage{latin}
\begin{center}{\color{gregoriocolor} Sexto Kaléndas Mártii. 
 Luna\dots\ }\end{center}
\switchcolumn
\selectlanguage{english}
\begin{center}{\color{gregoriocolor} The 25\textsuperscript{th} Day of February. The\dots\ Day of the Moon.}\end{center}
\end{paracol}

\noindent\begin{tabularx}{\linewidth}{*{19}{>{\centering\arraybackslash}X}}
 \textcolor{gregoriocolor}{a} & \textcolor{gregoriocolor}{b} & \textcolor{gregoriocolor}{c} & \textcolor{gregoriocolor}{d} & \textcolor{gregoriocolor}{e} & \textcolor{gregoriocolor}{f} & \textcolor{gregoriocolor}{g} & \textcolor{gregoriocolor}{h} & \textcolor{gregoriocolor}{i} & \textcolor{gregoriocolor}{k} & \textcolor{gregoriocolor}{l} & \textcolor{gregoriocolor}{m} & \textcolor{gregoriocolor}{n} & \textcolor{gregoriocolor}{p} & \textcolor{gregoriocolor}{q} & \textcolor{gregoriocolor}{r} & \textcolor{gregoriocolor}{s} & \textcolor{gregoriocolor}{t} & \textcolor{gregoriocolor}{u} \\
 26 & 27 & 28 & 29 & 1 & 2 & 3 & 4 & 5 & 6 & 7 & 8 & 9 & 10 & 11 & 12 & 13 & 14 & 15 \\
\end{tabularx}
\vspace{0.5\baselineskip}
\noindent\begin{tabularx}{\linewidth}{*{12}{>{\centering\arraybackslash}X}}
 \textcolor{gregoriocolor}{A} & \textcolor{gregoriocolor}{B} & \textcolor{gregoriocolor}{C} & \textcolor{gregoriocolor}{D} & \textcolor{gregoriocolor}{E} & F & \textcolor{gregoriocolor}{F} & \textcolor{gregoriocolor}{G} & \textcolor{gregoriocolor}{H} & \textcolor{gregoriocolor}{M} & \textcolor{gregoriocolor}{N} & \textcolor{gregoriocolor}{P} \\
 16 & 17 & 18 & 19 & 20 & 21 & 20 & 21 & 22 & 23 & 24 & 25 \\
\end{tabularx}

\begin{paracol}{2}
\selectlanguage{latin}
\lettrine[lines=2]{I}{n} Judæa natális sancti Matthíæ Apóstoli, qui, 
 post Ascensiónem Dómini ab Apóstolis in Judæ proditóris locum sorte eléctus, 
 pro Evangélii prædicatióne martyrium passus est.
\switchcolumn
\selectlanguage{english}
\lettrine[lines=2]{I}{n} Judea, the birthday of St. Matthias the Apostle. After the 
 Ascension of our Lord, the Apostles chose him, by lot, to fill the place of 
 Judas the traitor, and he suffered martyrdom for the preaching of the 
 Gospel.
\switchcolumn*
\selectlanguage{latin}
Romæ sanctæ Primitívæ Mártyris.
\switchcolumn
\selectlanguage{english}
At Rome, St. Primitiva, martyr.
\switchcolumn*
\selectlanguage{latin}
Rotómagi pássio sancti Prætextáti, Epíscopi et 
 Mártyris.
\switchcolumn
\selectlanguage{english}
At Rouen, the passion of St. Praetextatus, bishop and martyr.
\switchcolumn*
\selectlanguage{latin}
Cæsaréæ, in Cappadócia, sancti Sérgii Mártyris, 
 cujus gesta præclára habéntur.
\switchcolumn
\selectlanguage{english}
At Caesarea in Cappadocia, St. Sergius, martyr, of whose life a beautiful 
 account still exists.
\switchcolumn*
\selectlanguage{latin}
In Africa sanctórum Mártyrum Montáni, Lúcii, Juliáni, Victórici, 
 Flaviáni et Sociórum, qui discípuli fuérunt sancti Cypriáni, et, sub 
 Valeriáno Imperatóre, martyrium consummárunt.
\switchcolumn
\selectlanguage{english}
In Africa, the holy martyrs Montanus, Lucius, Julian, Victoricus, Flavian, 
 and their companions. They were disciples of St. Cyprian and suffered 
 martyrdom under Emperor Valerian.
\switchcolumn*
\selectlanguage{latin}
Tréviris sancti Modésti, Epíscopi et 
 Confessóris.
\switchcolumn
\selectlanguage{english}
At Treves, St. Modestus, bishop and confessor.
\switchcolumn*
\selectlanguage{latin}
Apud Stylum, in Calábria, sancti Joánnis, cognoménto Therísti, monásticæ 
 vitæ laude, et sanctitáte insígnis.
\switchcolumn
\selectlanguage{english}
At Stylo in Calabria, St. John Therestus, noted for his sanctity, and his 
 high regard for the monastic life.
\switchcolumn*
\selectlanguage{latin}
In Anglia sancti Edilbérti, Regis Cantiórum, quem sanctus Augustínus, 
 Anglórum Epíscopus, ad Christi fidem convértit.
\switchcolumn
\selectlanguage{english}
In England, St. Ethelbert, ruler of Kent, converted to the faith of Christ 
 by the English bishop, St. Augustine.
\switchcolumn*
\selectlanguage{latin}
Hierosólymis prima Invéntio cápitis sancti Joánnis Baptístæ, 
 Præcursóris Domini.
\switchcolumn
\selectlanguage{english}
At Jerusalem, the finding for the first time of the head of St. John the 
 Baptist, Precursor of the Lord.
\switchcolumn*
\selectlanguage{latin}
\end{paracol}


% ---- martyrology/mart02/mart0226.htm
\needspace{10\baselineskip}
\begin{paracol}{2}
\selectlanguage{latin}
\begin{center}{\color{gregoriocolor} Quarto Kaléndas Mártii. 
 Luna\dots\ }\end{center}
\switchcolumn
\selectlanguage{english}
\begin{center}{\color{gregoriocolor} The 26\textsuperscript{th} Day of February. The\dots\ Day of the Moon.}\end{center}
\end{paracol}

\noindent\begin{tabularx}{\linewidth}{*{19}{>{\centering\arraybackslash}X}}
 \textcolor{gregoriocolor}{a} & \textcolor{gregoriocolor}{b} & \textcolor{gregoriocolor}{c} & \textcolor{gregoriocolor}{d} & \textcolor{gregoriocolor}{e} & \textcolor{gregoriocolor}{f} & \textcolor{gregoriocolor}{g} & \textcolor{gregoriocolor}{h} & \textcolor{gregoriocolor}{i} & \textcolor{gregoriocolor}{k} & \textcolor{gregoriocolor}{l} & \textcolor{gregoriocolor}{m} & \textcolor{gregoriocolor}{n} & \textcolor{gregoriocolor}{p} & \textcolor{gregoriocolor}{q} & \textcolor{gregoriocolor}{r} & \textcolor{gregoriocolor}{s} & \textcolor{gregoriocolor}{t} & \textcolor{gregoriocolor}{u} \\
 28 & 29 & 1 & 2 & 3 & 4 & 5 & 6 & 7 & 8 & 9 & 10 & 11 & 12 & 13 & 14 & 15 & 16 & 17 \\
\end{tabularx}
\vspace{0.5\baselineskip}
\noindent\begin{tabularx}{\linewidth}{*{12}{>{\centering\arraybackslash}X}}
 \textcolor{gregoriocolor}{A} & \textcolor{gregoriocolor}{B} & \textcolor{gregoriocolor}{C} & \textcolor{gregoriocolor}{D} & \textcolor{gregoriocolor}{E} & F & \textcolor{gregoriocolor}{F} & \textcolor{gregoriocolor}{G} & \textcolor{gregoriocolor}{H} & \textcolor{gregoriocolor}{M} & \textcolor{gregoriocolor}{N} & \textcolor{gregoriocolor}{P} \\
 18 & 19 & 20 & 21 & 22 & 23 & 22 & 23 & 24 & 25 & 26 & 27 \\
\end{tabularx}

\begin{paracol}{2}
\selectlanguage{latin}
\lettrine[lines=2]{P}{erge,} in Pamphylia, natális beáti Néstoris 
 Epíscopi, qui, in persecutióne Décii, cum diu noctúque oratióni insísteret 
 póstulans ut grex Christi custodirétur, comprehénsus est, ac, nomen Dómini 
 mira liberáte et alacritáte conféssus, Præsidis Polliónis jussu equúleo 
 sævíssime est cruciátus; ac demum, cum se Christo semper adhæsúrum 
 constánter profiterétur, crucis suspéndio victor in cælum migrávit.
\switchcolumn
\selectlanguage{english}
\lettrine[lines=2]{A}{t} Pergen in Pamphylia, during the persecution of Decius, the birthday of 
 blessed Nestor, bishop, who praying night and day for the safety of the 
 flock of Christ, was put under arrest. Because he confessed the Name 
 of the Lord with great zeal and freedom, he was cruelly tortured on the rack 
 by order of Pollio the governor. When he still courageously proclaimed 
 that he would remain ever faithful to Christ, he was crucified, and thus 
 triumphantly went to heaven.
\switchcolumn*
\selectlanguage{latin}
Ibídem pássio sanctórum Pápiæ, Diodóri, Conónis 
 et Claudiáni, qui sanctum Néstorem martyrio præcessérunt.
\switchcolumn
\selectlanguage{english}
In the same place, the passion of Saints Papias, Diodorus, Conon, and 
 Claudian, who preceded St. Nestor to martyrdom.
\switchcolumn*
\selectlanguage{latin}
Item sanctórum Mártyrum Fortunáti, Felícis et aliórum vigínti septem.
\switchcolumn
\selectlanguage{english}
Also, the holy martyrs Fortunatus, Felix, and twenty-seven others.
\switchcolumn*
\selectlanguage{latin}
Alexandríæ sancti Alexándri Epíscopi, gloriósi 
 senis, qui, post beátum Petrum, ejúsdem civitátis Epíscopum, zelo fídei 
 succénsus, Aríum, Presbyterum suum, hærética impietáte depravátum et divina 
 veritáte convíctum, de Ecclésia ejécit; ac póstea inter trecéntos decem et 
 octo Patres, in Nicǽno Concílio eúndem damnávit.
\switchcolumn
\selectlanguage{english}
At Alexandria, Bishop St. Alexander, an aged man held in great honour, who 
 succeeded blessed Peter as bishop of that city. He expelled Arius, one 
 of his priests, from the Church because he was tainted with heretical 
 ímpiety and convicted in the face of divine truth. Later on he was one 
 of the three hundred and eighteen Fathers who condemned him in the Council 
 of Nicaea.
\switchcolumn*
\selectlanguage{latin}
Bonóniæ sancti Faustiniáni Epíscopi, qui eam 
 Ecclésiam, Diocletiáni persecutióne vexátam, verbo prædicatiónis firmávit et 
 auxit.
\switchcolumn
\selectlanguage{english}
At Bologna, the bishop St. Faustinian. His preaching strengthened and 
 multiplied the faithful of that church when it was so much afflicted during 
 the persecution of Diocletian.
\switchcolumn*
\selectlanguage{latin}
Gazæ, in Palæstína, sancti Porphyrii Epíscopi, 
 qui, témpore Arcádii Imperatóris, Marnam idólum ejúsque templum evértit, ac, 
 multa passus, quiévit in Dómino.
\switchcolumn
\selectlanguage{english}
At Gaza in Palestine, St. Porphyry, bishop, in the time of Emperor Arcadius. 
 He overthrew the idol Marna and its temple, and after many sufferings, went 
 to his rest in the Lord.
\switchcolumn*
\selectlanguage{latin}
Floréntiæ sancti Andréæ, Epíscopi et 
 Confessóris.
\switchcolumn
\selectlanguage{english}
At Florence, St. Andrew, bishop and confessor.
\switchcolumn*
\selectlanguage{latin}
In território Archiacénsi, in Gállia, sancti Victóris Confessóris, cujus 
 laudes sanctus Bernárdus conscrípsit.
\switchcolumn
\selectlanguage{english}
In the province of Champagne in France, St. Victor, confessor, about whom 
 eulogies have been written by St. Bernard.
\switchcolumn*
\selectlanguage{latin}
\end{paracol}


% ---- martyrology/mart02/mart0226leap.htm
\needspace{10\baselineskip}
\begin{paracol}{2}
\selectlanguage{latin}
\begin{center}{\color{gregoriocolor} Quinto Kaléndas Mártii. 
 Luna\dots\ }\end{center}
\switchcolumn
\selectlanguage{english}
\begin{center}{\color{gregoriocolor} The 26\textsuperscript{th} Day of February. The\dots\ Day of the Moon.}\end{center}
\end{paracol}

\noindent\begin{tabularx}{\linewidth}{*{19}{>{\centering\arraybackslash}X}}
 \textcolor{gregoriocolor}{a} & \textcolor{gregoriocolor}{b} & \textcolor{gregoriocolor}{c} & \textcolor{gregoriocolor}{d} & \textcolor{gregoriocolor}{e} & \textcolor{gregoriocolor}{f} & \textcolor{gregoriocolor}{g} & \textcolor{gregoriocolor}{h} & \textcolor{gregoriocolor}{i} & \textcolor{gregoriocolor}{k} & \textcolor{gregoriocolor}{l} & \textcolor{gregoriocolor}{m} & \textcolor{gregoriocolor}{n} & \textcolor{gregoriocolor}{p} & \textcolor{gregoriocolor}{q} & \textcolor{gregoriocolor}{r} & \textcolor{gregoriocolor}{s} & \textcolor{gregoriocolor}{t} & \textcolor{gregoriocolor}{u} \\
 27 & 28 & 29 & 1 & 2 & 3 & 4 & 5 & 6 & 7 & 8 & 9 & 10 & 11 & 12 & 13 & 14 & 15 & 16 \\
\end{tabularx}
\vspace{0.5\baselineskip}
\noindent\begin{tabularx}{\linewidth}{*{12}{>{\centering\arraybackslash}X}}
 \textcolor{gregoriocolor}{A} & \textcolor{gregoriocolor}{B} & \textcolor{gregoriocolor}{C} & \textcolor{gregoriocolor}{D} & \textcolor{gregoriocolor}{E} & F & \textcolor{gregoriocolor}{F} & \textcolor{gregoriocolor}{G} & \textcolor{gregoriocolor}{H} & \textcolor{gregoriocolor}{M} & \textcolor{gregoriocolor}{N} & \textcolor{gregoriocolor}{P} \\
 17 & 18 & 19 & 20 & 21 & 22 & 21 & 22 & 23 & 24 & 25 & 26 \\
\end{tabularx}

\begin{paracol}{2}
\selectlanguage{latin}
\lettrine[lines=2]{I}{n} Ægypto natális sanctórum Mártyrum Victoríni, 
 Victóris, Nicéphori, Claudiáni, Dióscori, Serapiónis et Pápiæ, sub Numeriáno 
 Imperatóre. Horum duo primi, pro confessióne fídei, exquisíta 
 suppliciórum génera constánter passi, cápite plectúntur; Nicéphorus, post 
 cratículas candéntes ignésque superátos, minutátim concísus est; Claudiánus 
 et Dióscorus flammis incénsi; Serápion vero et Pápias gládio cæsi sunt.
\switchcolumn
\selectlanguage{english}
\lettrine[lines=2]{I}{n} Egypt, under Emperor Numerian, the birthday of the holy martyrs 
 Victorinus, Victor, Nicephorus, Claudian, Dioscorus, Serapion, and Papias. 
 After patiently enduring extreme tortures, the first two were beheaded for 
 the confession of the faith, Nicephorus was laid on a heated gridiron, 
 placed over the fire, then thoroughly hacked with a knife; Claudian and 
 Dioscorus were burned at the stake; Serapion and Papias were slain with the 
 sword.
\switchcolumn*
\selectlanguage{latin}
In Africa sanctórum Mártyrum Donáti, Justi, Herénæ 
 et Sociórum.
\switchcolumn
\selectlanguage{english}
In Africa, the holy martyrs Donatus, Justus, Herenas, and their companions.
\switchcolumn*
\selectlanguage{latin}
Constantinópoli sancti Tharásii Epíscopi, eruditióne et pietáte insígnis; ad 
 quem exstat Hadriáni Papæ Primi epístola pro 
 defensióne sanctárum Imáginum.
\switchcolumn
\selectlanguage{english}
At Constantinople, St. Tharasius, bishop, a man of great learning and piety. 
 There exists a letter defending sacred images, written to him by Pope 
 Hadrian I.
\switchcolumn*
\selectlanguage{latin}
Naziánzi, in Cappadócia, sancti Cæsárii, qui 
 beátæ Nonnæ fílius ac beatórum Gregórii Theólogi et Gorgóniæ fuit frater, et 
 quem idem Gregórius inter ágmina beatórum se vidísse testátur.
\switchcolumn
\selectlanguage{english}
At Nazianzus, St. Caesarius, who was the son of blessed Nonna, and whom his 
 brother, blessed Gregory the Theologian, says he saw among the hosts of the 
 blessed.
\switchcolumn*
\selectlanguage{latin}
In monastério Heidenhémii, diœcésis 
 Eistetténsis, in Germánia, sanctæ Walbúrgæ Vírginis, quæ fuit fília sancti 
 Richárdi, Anglórum Regis, et soror sancti Willebáldi, Eistetténsis Epíscopi.
\switchcolumn
\selectlanguage{english}
In the monastery of Heidenheim, in the Eichstadt diocese in Germany, St. 
 Walburga, virgin. She was the daughter of St. Richard, king of 
 England, and sister of St. Willebald, bishop of Eichstadt.
\switchcolumn*
\selectlanguage{latin}
\end{paracol}


% ---- martyrology/mart02/mart0227.htm
\needspace{10\baselineskip}
\begin{paracol}{2}
\selectlanguage{latin}
\begin{center}{\color{gregoriocolor} Tértio Kaléndas Mártii. 
 Luna\dots\ }\end{center}
\switchcolumn
\selectlanguage{english}
\begin{center}{\color{gregoriocolor} The 27\textsuperscript{th} Day of February. The\dots\ Day of the Moon.}\end{center}
\end{paracol}

\noindent\begin{tabularx}{\linewidth}{*{19}{>{\centering\arraybackslash}X}}
 \textcolor{gregoriocolor}{a} & \textcolor{gregoriocolor}{b} & \textcolor{gregoriocolor}{c} & \textcolor{gregoriocolor}{d} & \textcolor{gregoriocolor}{e} & \textcolor{gregoriocolor}{f} & \textcolor{gregoriocolor}{g} & \textcolor{gregoriocolor}{h} & \textcolor{gregoriocolor}{i} & \textcolor{gregoriocolor}{k} & \textcolor{gregoriocolor}{l} & \textcolor{gregoriocolor}{m} & \textcolor{gregoriocolor}{n} & \textcolor{gregoriocolor}{p} & \textcolor{gregoriocolor}{q} & \textcolor{gregoriocolor}{r} & \textcolor{gregoriocolor}{s} & \textcolor{gregoriocolor}{t} & \textcolor{gregoriocolor}{u} \\
 29 & 1 & 2 & 3 & 4 & 5 & 6 & 7 & 8 & 9 & 10 & 11 & 12 & 13 & 14 & 15 & 16 & 17 & 18 \\
\end{tabularx}
\vspace{0.5\baselineskip}
\noindent\begin{tabularx}{\linewidth}{*{12}{>{\centering\arraybackslash}X}}
 \textcolor{gregoriocolor}{A} & \textcolor{gregoriocolor}{B} & \textcolor{gregoriocolor}{C} & \textcolor{gregoriocolor}{D} & \textcolor{gregoriocolor}{E} & F & \textcolor{gregoriocolor}{F} & \textcolor{gregoriocolor}{G} & \textcolor{gregoriocolor}{H} & \textcolor{gregoriocolor}{M} & \textcolor{gregoriocolor}{N} & \textcolor{gregoriocolor}{P} \\
 19 & 20 & 21 & 22 & 23 & 24 & 23 & 24 & 25 & 26 & 27 & 28 \\
\end{tabularx}

\begin{paracol}{2}
\selectlanguage{latin}
\lettrine[lines=2]{I}{nsulæ,} in Aprútio, sancti Gabriélis a Vírgine 
 Perdolénte, Clérici Congregatiónis a Cruce et Passióne Dómini nuncupátæ, et 
 Confessóris; qui, magnis intra breve vitæ spátium méritis et post mortem 
 miráculis clarus, a Benedícto Papa Décimo quinto in Sanctórum cánonem 
 relátus est.
\switchcolumn
\selectlanguage{english}
\lettrine[lines=2]{A}{t} Isola, in the province of Abruzzi, St. Gabriel of the Sorrowful Virgin, 
 confessor and cleric of the Passionist Congregation. Having been known 
 for his merits during his short life, and after death renowned for miracles, 
 Pope Benedict XV enrolled him in the canon of the saints.
\switchcolumn*
\selectlanguage{latin}
Romæ natális sanctórum Mártyrum Alexándri, 
 Abúndii, Antígoni et Fortunáti.
\switchcolumn
\selectlanguage{english}
At Rome, the birthday of the holy martyrs, Alexander, Abundius, Antigonus, 
 and Fortunatus.
\switchcolumn*
\selectlanguage{latin}
Alexandríæ pássio sancti Juliáni Mártyris, qui, 
 cum ita pódagra constríctus esset, ut neque incédere neque stare posset, una 
 cum duóbus fámulis, qui eum in sella gestábant, Júdici offértur; quorum 
 alter fidem negávit, alter, nómine Eunus, cum dómino suo perdurávit in 
 confessióne Christi. Ipse porro Juliánus et Eunus, camélis impósiti, 
 per totam urbem circumdúci jubéntur, et flagris laniári, ac tandem, incénso 
 rogo, hinc inde spectánte pópulo, combúri.
\switchcolumn
\selectlanguage{english}
At Alexandria, the passion of St. Julian, martyr. Although he was so 
 afflicted with gout that he could neither walk nor stand, he was taken 
 before the judge with two servants, who carried him in a chair. One of 
 these denied his faith, but the other, named Eunus, persevered with Julian 
 in confessing Christ. Both were set on camels, led through the whole 
 city, scourged, and then burned alive in the presence of all the people.
\switchcolumn*
\selectlanguage{latin}
Ibídem sancti Besæ mílitis, qui, cum 
 insultántes in prædíctos Mártyres cohibéret, delátus est ad Júdicem, et, pro 
 fide constánter agens, cápite truncátus.
\switchcolumn
\selectlanguage{english}
In the same city, St. Besas, a soldier. He had rebuked those who 
 insulted the martyrs just mentioned, and so was denounced before the judge. 
 Because he continued to proclaim his attachment to the faith he was 
 beheaded.
\switchcolumn*
\selectlanguage{latin}
Híspali, in Hispánia, natális sancti Leándri, 
 ejúsdem civitátis Epíscopi, qui, sanctórum Isidóri Epíscopi ac Florentínæ 
 Vírginis frater, sua prædicatióne et indústria gentem Visigothórum, 
 adjuvánte Reccarédo, eórum Rege, ab Ariána impietáte ad cathólicam fidem 
 convértit.
\switchcolumn
\selectlanguage{english}
At Seville in Spain, the birthday of St. Leander, bishop of that city, and 
 of St. Florentina, virgin. By his preaching and zeal the Visigoths, 
 with the help of King Recared, were converted from the Arian heresy to the 
 Catholic faith.
\switchcolumn*
\selectlanguage{latin}
Constantinópoli sanctórum Confessórum Basilíi 
 et Procópii, qui, témpore Leónis Imperatóris, pro cultu sanctárum Imáginum 
 strénue decertárunt.
\switchcolumn
\selectlanguage{english}
At Constantinople, in the time of Emperor Leo, the holy confessors Basil and 
 Procopius, who fought courageously in behalf of the veneration of sacred 
 images.
\switchcolumn*
\selectlanguage{latin}
Lugdúni, in Gállia, sancti Baldoméri Subdiáconi, viri Deo devóti, cujus 
 sepúlcrum crebris miráculis illustrátur.
\switchcolumn
\selectlanguage{english}
At Lyons, St. Baldomer, subdeacon and man of God, whose tomb is graced by 
 many miracles.
\switchcolumn*
\selectlanguage{latin}
\end{paracol}


% ---- martyrology/mart02/mart0227leap.htm
\needspace{10\baselineskip}
\begin{paracol}{2}
\selectlanguage{latin}
\begin{center}{\color{gregoriocolor} Quarto Kaléndas Mártii. 
 Luna\dots\ }\end{center}
\switchcolumn
\selectlanguage{english}
\begin{center}{\color{gregoriocolor} The 27\textsuperscript{th} Day of February. The\dots\ Day of the Moon.}\end{center}
\end{paracol}

\noindent\begin{tabularx}{\linewidth}{*{19}{>{\centering\arraybackslash}X}}
 \textcolor{gregoriocolor}{a} & \textcolor{gregoriocolor}{b} & \textcolor{gregoriocolor}{c} & \textcolor{gregoriocolor}{d} & \textcolor{gregoriocolor}{e} & \textcolor{gregoriocolor}{f} & \textcolor{gregoriocolor}{g} & \textcolor{gregoriocolor}{h} & \textcolor{gregoriocolor}{i} & \textcolor{gregoriocolor}{k} & \textcolor{gregoriocolor}{l} & \textcolor{gregoriocolor}{m} & \textcolor{gregoriocolor}{n} & \textcolor{gregoriocolor}{p} & \textcolor{gregoriocolor}{q} & \textcolor{gregoriocolor}{r} & \textcolor{gregoriocolor}{s} & \textcolor{gregoriocolor}{t} & \textcolor{gregoriocolor}{u} \\
 28 & 29 & 1 & 2 & 3 & 4 & 5 & 6 & 7 & 8 & 9 & 10 & 11 & 12 & 13 & 14 & 15 & 16 & 17 \\
\end{tabularx}
\vspace{0.5\baselineskip}
\noindent\begin{tabularx}{\linewidth}{*{12}{>{\centering\arraybackslash}X}}
 \textcolor{gregoriocolor}{A} & \textcolor{gregoriocolor}{B} & \textcolor{gregoriocolor}{C} & \textcolor{gregoriocolor}{D} & \textcolor{gregoriocolor}{E} & F & \textcolor{gregoriocolor}{F} & \textcolor{gregoriocolor}{G} & \textcolor{gregoriocolor}{H} & \textcolor{gregoriocolor}{M} & \textcolor{gregoriocolor}{N} & \textcolor{gregoriocolor}{P} \\
 18 & 19 & 20 & 21 & 22 & 23 & 22 & 23 & 24 & 25 & 26 & 27 \\
\end{tabularx}

\begin{paracol}{2}
\selectlanguage{latin}
\lettrine[lines=2]{P}{erge,} in Pamphylia, natális beáti Néstoris 
 Epíscopi, qui, in persecutióne Décii, cum diu noctúque oratióni insísteret 
 póstulans ut grex Christi custodirétur, comprehénsus est, ac, nomen Dómini 
 mira liberáte et alacritáte conféssus, Prǽsidis Polliónis jussu equúleo 
 sævíssime est cruciátus; ac demum, cum se Christo semper adhæsúrum 
 constánter profiterétur, crucis suspéndio victor in cælum migrávit.
\switchcolumn
\selectlanguage{english}
\lettrine[lines=2]{A}{t} Pergen in Pamphylia, during the persecution of Decius, the birthday of 
 blessed Nestor, bishop, who praying night and day for the safety of the 
 flock of Christ, was put under arrest. Because he confessed the Name 
 of the Lord with great zeal and freedom, he was cruelly tortured on the rack 
 by order of Pollio the governor. When he still courageously proclaimed 
 that he would remain ever faithful to Christ, he was crucified, and thus 
 triumphantly went to heaven.
\switchcolumn*
\selectlanguage{latin}
Ibídem pássio sanctórum Pápiæ, Diodóri, Conónis 
 et Claudiáni, qui sanctum Néstorem martyrio præcessérunt.
\switchcolumn
\selectlanguage{english}
In the same place, the passion of Saints Papias, Diodorus, Conon, and 
 Claudian, who preceded St. Nestor to martyrdom.
\switchcolumn*
\selectlanguage{latin}
Item sanctórum Mártyrum Fortunáti, Felícis et aliórum vigínti septem.
\switchcolumn
\selectlanguage{english}
Also, the holy martyrs Fortunatus, Felix, and twenty-seven others.
\switchcolumn*
\selectlanguage{latin}
Alexandríæ sancti Alexándri Epíscopi, gloriósi 
 senis, qui, post beátum Petrum, ejúsdem civitátis Episcopum, zelo fídei 
 succénsus, Aríum, Presbyterum suum, hærética impietáte depravátum et divína 
 veritáte convíctum, de Ecclésia ejécit; ac póstea inter trecéntos decem et 
 octo Patres, in Nicǽno Concílio eúndem damnávit.
\switchcolumn
\selectlanguage{english}
At Alexandria, Bishop St. Alexander, an aged man held in great honour, who 
 succeeded blessed Peter as bishop of that city. He expelled Arius, one 
 of his priests, from the Church because he was tainted with heretical 
 ímpiety and convicted in the face of divine truth. Later on he was one 
 of the three hundred and eighteen Fathers who condemned him in the Council 
 of Nicaea.
\switchcolumn*
\selectlanguage{latin}
Bonóniæ sancti Faustiniáni Epíscopi, qui eam 
 Ecclésiam, Diocletiáni persecutióne vexátam, verbo prædicatiónis firmávit et 
 auxit.
\switchcolumn
\selectlanguage{english}
At Bologna, the bishop St. Faustinian. His preaching strengthened and 
 multiplied the faithful of that church when it was so much afflicted during 
 the persecution of Diocletian.
\switchcolumn*
\selectlanguage{latin}
Gazæ, in Palæstína, sancti Porphyrii Epíscopi, 
 qui, témpore Arcádii Imperatóris, Marnam idólum ejúsque templum evértit, ac, 
 multa passus, quiévit in Dómino.
\switchcolumn
\selectlanguage{english}
At Gaza in Palestine, St. Porphyry, bishop, in the time of Emperor Arcadius. 
 He overthrew the idol Marna and its temple, and after many sufferings, went 
 to his rest in the Lord.
\switchcolumn*
\selectlanguage{latin}
Floréntiæ sancti Andréæ, Epíscopi et 
 Confessóris.
\switchcolumn
\selectlanguage{english}
At Florence, St. Andrew, bishop and confessor.
\switchcolumn*
\selectlanguage{latin}
In território Archiacénsi, in Gállia, sancti Victóris Confessóris, cujus 
 laudes sanctus Bernárdus conscrípsit.
\switchcolumn
\selectlanguage{english}
In the province of Champagne in France, St. Victor, confessor, about whom 
 eulogies have been written by St. Bernard.
\switchcolumn*
\selectlanguage{latin}
\end{paracol}


% ---- martyrology/mart02/mart0228.htm
\needspace{10\baselineskip}
\begin{paracol}{2}
\selectlanguage{latin}
\begin{center}{\color{gregoriocolor} Prídie Kaléndas Mártii. 
 Luna\dots\ }\end{center}
\switchcolumn
\selectlanguage{english}
\begin{center}{\color{gregoriocolor} The 28\textsuperscript{th} Day of February. The\dots\ Day of the Moon.}\end{center}
\end{paracol}

\noindent\begin{tabularx}{\linewidth}{*{19}{>{\centering\arraybackslash}X}}
 \textcolor{gregoriocolor}{a} & \textcolor{gregoriocolor}{b} & \textcolor{gregoriocolor}{c} & \textcolor{gregoriocolor}{d} & \textcolor{gregoriocolor}{e} & \textcolor{gregoriocolor}{f} & \textcolor{gregoriocolor}{g} & \textcolor{gregoriocolor}{h} & \textcolor{gregoriocolor}{i} & \textcolor{gregoriocolor}{k} & \textcolor{gregoriocolor}{l} & \textcolor{gregoriocolor}{m} & \textcolor{gregoriocolor}{n} & \textcolor{gregoriocolor}{p} & \textcolor{gregoriocolor}{q} & \textcolor{gregoriocolor}{r} & \textcolor{gregoriocolor}{s} & \textcolor{gregoriocolor}{t} & \textcolor{gregoriocolor}{u} \\
 1 & 2 & 3 & 4 & 5 & 6 & 7 & 8 & 9 & 10 & 11 & 12 & 13 & 14 & 15 & 16 & 17 & 18 & 19 \\
\end{tabularx}
\vspace{0.5\baselineskip}
\noindent\begin{tabularx}{\linewidth}{*{12}{>{\centering\arraybackslash}X}}
 \textcolor{gregoriocolor}{A} & \textcolor{gregoriocolor}{B} & \textcolor{gregoriocolor}{C} & \textcolor{gregoriocolor}{D} & \textcolor{gregoriocolor}{E} & F & \textcolor{gregoriocolor}{F} & \textcolor{gregoriocolor}{G} & \textcolor{gregoriocolor}{H} & \textcolor{gregoriocolor}{M} & \textcolor{gregoriocolor}{N} & \textcolor{gregoriocolor}{P} \\
 20 & 21 & 22 & 23 & 24 & 25 & 24 & 25 & 26 & 27 & 28 & 29 \\
\end{tabularx}

\begin{paracol}{2}
\selectlanguage{latin}
\lettrine[lines=2]{R}{omæ} natális sanctórum Mártyrum Macárii, Rufíni, 
 Justi et Theóphili.
\switchcolumn
\selectlanguage{english}
\lettrine[lines=2]{A}{t} Rome, the birthday of the holy martyrs Macarius, Rufinus, Justus, and 
 Theophilus.
\switchcolumn*
\selectlanguage{latin}
Alexandríæ pássio sanctórum Cæreális, Púpuli, 
 Caji et Serapiónis.
\switchcolumn
\selectlanguage{english}
At Alexandria, the passion of the Saints Caerealis, Pupulus, Caius, and 
 Serapion.
\switchcolumn*
\selectlanguage{latin}
Ibídem commemorátio sanctórum Presbyterórum, Diaconórum et aliórum 
 plurimórum; qui, témpore Valeriáni Imperatóris, cum pestis sævíssima 
 grassarétur, morbo laborántibus ministrántes, libentíssime mortem oppetiére, 
 et quos velut Mártyres religiósa piórum fides venerári consuévit.
\switchcolumn
\selectlanguage{english}
In the same city, in the reign of Emperor Valerian, the commemoration of the 
 holy priests, deacons, and many others. When a most deadly epidemic 
 was raging, they willingly met their death by ministering to the sick. 
 The religious sentiment of the pious faithful has generally venerated them 
 as martyrs.
\switchcolumn*
\selectlanguage{latin}
Romæ sancti Hílari, Papæ et Confessóris.
\switchcolumn
\selectlanguage{english}
At Rome, St. Hilary, pope and confessor.
\switchcolumn*
\selectlanguage{latin}
In território Lugdunénsi, locis Jurénsibus, deposítio sancti Románi Abbátis, 
 qui primum illic eremíticam vitam duxit, et, multis virtútibus ac miráculis 
 clarus, plurimórum póstea Pater éxstitit Monachórum.
\switchcolumn
\selectlanguage{english}
In the territory of Lyons, in the Jura Mountains, the death of St. Romanus, 
 abbot, who first had led the life of a hermit there. His reputation 
 for virtues and miracles brought under his guidance many monks.
\switchcolumn*
\selectlanguage{latin}
Papíæ Translátio córporis sancti Augustíni 
 Epíscopi, Confessóris et Ecclésiæ Doctóris, ex Sardínia ínsula, ópera 
 Luitprándi, Regis Longobardórum.
\switchcolumn
\selectlanguage{english}
At Papia, the transfer, ordered by the Lombard King Luitprand, of the body 
 of St. Augustine, bishop, away from the island of Sardinia.
\switchcolumn*
\selectlanguage{latin}
\end{paracol}


% ---- martyrology/mart02/mart0228leap.htm
\needspace{10\baselineskip}
\begin{paracol}{2}
\selectlanguage{latin}
\begin{center}{\color{gregoriocolor} Tértio Kaléndas Mártii. 
 Luna\dots\ }\end{center}
\switchcolumn
\selectlanguage{english}
\begin{center}{\color{gregoriocolor} The 28\textsuperscript{th} Day of February. The\dots\ Day of the Moon.}\end{center}
\end{paracol}

\noindent\begin{tabularx}{\linewidth}{*{19}{>{\centering\arraybackslash}X}}
 \textcolor{gregoriocolor}{a} & \textcolor{gregoriocolor}{b} & \textcolor{gregoriocolor}{c} & \textcolor{gregoriocolor}{d} & \textcolor{gregoriocolor}{e} & \textcolor{gregoriocolor}{f} & \textcolor{gregoriocolor}{g} & \textcolor{gregoriocolor}{h} & \textcolor{gregoriocolor}{i} & \textcolor{gregoriocolor}{k} & \textcolor{gregoriocolor}{l} & \textcolor{gregoriocolor}{m} & \textcolor{gregoriocolor}{n} & \textcolor{gregoriocolor}{p} & \textcolor{gregoriocolor}{q} & \textcolor{gregoriocolor}{r} & \textcolor{gregoriocolor}{s} & \textcolor{gregoriocolor}{t} & \textcolor{gregoriocolor}{u} \\
 29 & 1 & 2 & 3 & 4 & 5 & 6 & 7 & 8 & 9 & 10 & 11 & 12 & 13 & 14 & 15 & 16 & 17 & 18 \\
\end{tabularx}
\vspace{0.5\baselineskip}
\noindent\begin{tabularx}{\linewidth}{*{12}{>{\centering\arraybackslash}X}}
 \textcolor{gregoriocolor}{A} & \textcolor{gregoriocolor}{B} & \textcolor{gregoriocolor}{C} & \textcolor{gregoriocolor}{D} & \textcolor{gregoriocolor}{E} & F & \textcolor{gregoriocolor}{F} & \textcolor{gregoriocolor}{G} & \textcolor{gregoriocolor}{H} & \textcolor{gregoriocolor}{M} & \textcolor{gregoriocolor}{N} & \textcolor{gregoriocolor}{P} \\
 19 & 20 & 21 & 22 & 23 & 24 & 23 & 24 & 25 & 26 & 27 & 28 \\
\end{tabularx}

\begin{paracol}{2}
\selectlanguage{latin}
\lettrine[lines=2]{I}{nsulæ,} in Aprútio, sancti Gabriélis a Vírgine 
 Perdolénte, Clérici Congregatiónis a Cruce et Passióne Dómini nuncupátæ, et 
 Confessóris; qui, magnis intra breve vitæ spátium méritis et post mortem 
 miráculis clarus, a Benedícto Papa Décimo quinto in Sanctórum cánonem 
 relátus est.
\switchcolumn
\selectlanguage{english}
\lettrine[lines=2]{A}{t} Isola, in the province of Abruzzi, St. Gabriel of the Sorrowful Virgin, 
 confessor and cleric of the Passionist Congregation. Having been known 
 for his merits during his short life, and after death renowned for miracles, 
 Pope Benedict XV enrolled him in the canon of the saints.
\switchcolumn*
\selectlanguage{latin}
Romæ natális sanctórum Mártyrum Alexándri, 
 Abúndii, Antígoni et Fortunáti.
\switchcolumn
\selectlanguage{english}
At Rome, the birthday of the holy martyrs, Alexander, Abundius, Antigonus, 
 and Fortunatus.
\switchcolumn*
\selectlanguage{latin}
Alexandríæ pássio sancti Juliáni Mártyris, qui, 
 cum ita pódagra constríctus esset, ut neque incédere neque stare posset, una 
 cum duóbus fámulis, qui eum in sella gestábant, Júdici offértur; quorum 
 alter fidem negávit, alter, nómine Eunus, cum dómino suo perdurávit in 
 confessióne Christi. Ipse porro Juliánus et Eunus, camélis impósiti, 
 per totam urbem circumdúci jubéntur, et flagris laniári, ac tandem, incénso 
 rogo, hinc inde spectánte pópulo, combúri.
\switchcolumn
\selectlanguage{english}
At Alexandria, the passion of St. Julian, martyr. Although he was so 
 afflicted with gout that he could neither walk nor stand, he was taken 
 before the judge with two servants, who carried him in a chair. One of 
 these denied his faith, but the other, named Eunus, persevered with Julian 
 in confessing Christ. Both were set on camels, led through the whole 
 city, scourged, and then burned alive in the presence of all the people.
\switchcolumn*
\selectlanguage{latin}
Ibídem sancti Besæ mílitis, qui, cum 
 insultántes in prædictos Mártyres cohibéret, delátus est ad Júdicem, et, pro 
 fide constánter agens, cápite truncátus.
\switchcolumn
\selectlanguage{english}
In the same city, St. Besas, a soldier. He had rebuked those who 
 insulted the martyrs just mentioned, and so was denounced before the judge. 
 Because he continued to proclaim his attachment to the faith he was 
 beheaded.
\switchcolumn*
\selectlanguage{latin}
Híspali, in Hispánia, natális sancti Leándri, 
 ejúsdem civitátis Epíscopi, qui, sanctórum Isidóri Epíscopi ac Florentínæ 
 Vírginis frater, sua prædicatióne et indústria gentem Visigothórum, 
 adjuvánte Reccarédo, eórum Rege, ab Ariána impietáte ad cathólicam fidem 
 convértit.
\switchcolumn
\selectlanguage{english}
At Seville in Spain, the birthday of St. Leander, bishop of that city, and 
 of St. Florentina, virgin. By his preaching and zeal the Visigoths, 
 with the help of King Recared, were converted from the Arian heresy to the 
 Catholic faith.
\switchcolumn*
\selectlanguage{latin}
Constantinópoli sanctórum Confessórum Basilíi 
 et Procópii, qui, témpore Leónis Imperatóris, pro cultu sanctárum Imáginum 
 strénue decertárunt.
\switchcolumn
\selectlanguage{english}
At Constantinople, in the time of Emperor Leo, the holy confessors Basil and 
 Procopius, who fought courageously in behalf of the veneration of sacred 
 images.
\switchcolumn*
\selectlanguage{latin}
Lugdúni, in Gállia, sancti Baldoméri Subdiáconi, viri Deo devóti, cujus 
 sepúlcrum crebris miráculis illustrátur.
\switchcolumn
\selectlanguage{english}
At Lyons, St. Baldomer, subdeacon and man of God, whose tomb is graced by 
 many miracles.
\switchcolumn*
\selectlanguage{latin}
\end{paracol}


% ---- martyrology/mart02/mart0229.leap.htm
\needspace{10\baselineskip}
\begin{paracol}{2}
\selectlanguage{latin}
\begin{center}{\color{gregoriocolor} Prídie Kaléndas Mártii. 
 Luna\dots\ }\end{center}
\switchcolumn
\selectlanguage{english}
\begin{center}{\color{gregoriocolor} The 29\textsuperscript{th} Day of February. The\dots\ Day of the Moon.}\end{center}
\end{paracol}

\noindent\begin{tabularx}{\linewidth}{*{19}{>{\centering\arraybackslash}X}}
 \textcolor{gregoriocolor}{a} & \textcolor{gregoriocolor}{b} & \textcolor{gregoriocolor}{c} & \textcolor{gregoriocolor}{d} & \textcolor{gregoriocolor}{e} & \textcolor{gregoriocolor}{f} & \textcolor{gregoriocolor}{g} & \textcolor{gregoriocolor}{h} & \textcolor{gregoriocolor}{i} & \textcolor{gregoriocolor}{k} & \textcolor{gregoriocolor}{l} & \textcolor{gregoriocolor}{m} & \textcolor{gregoriocolor}{n} & \textcolor{gregoriocolor}{p} & \textcolor{gregoriocolor}{q} & \textcolor{gregoriocolor}{r} & \textcolor{gregoriocolor}{s} & \textcolor{gregoriocolor}{t} & \textcolor{gregoriocolor}{u} \\
 1 & 2 & 3 & 4 & 5 & 6 & 7 & 8 & 9 & 10 & 11 & 12 & 13 & 14 & 15 & 16 & 17 & 18 & 19 \\
\end{tabularx}
\vspace{0.5\baselineskip}
\noindent\begin{tabularx}{\linewidth}{*{12}{>{\centering\arraybackslash}X}}
 \textcolor{gregoriocolor}{A} & \textcolor{gregoriocolor}{B} & \textcolor{gregoriocolor}{C} & \textcolor{gregoriocolor}{D} & \textcolor{gregoriocolor}{E} & F & \textcolor{gregoriocolor}{F} & \textcolor{gregoriocolor}{G} & \textcolor{gregoriocolor}{H} & \textcolor{gregoriocolor}{M} & \textcolor{gregoriocolor}{N} & \textcolor{gregoriocolor}{P} \\
 20 & 21 & 22 & 23 & 24 & 25 & 24 & 25 & 26 & 27 & 28 & 29 \\
\end{tabularx}

\begin{paracol}{2}
\selectlanguage{latin}
\lettrine[lines=2]{R}{omæ} natális sanctórum Mártyrum Macárii, Rufíni, 
 Justi et Theóphili.
\switchcolumn
\selectlanguage{english}
\lettrine[lines=2]{A}{t} Rome, the birthday of the holy martyrs Macarius, Rufinus, Justus, and 
 Theophilus.
\switchcolumn*
\selectlanguage{latin}
Alexandríæ pássio sanctórum Cæreális, Púpuli, 
 Caji et Serapiónis.
\switchcolumn
\selectlanguage{english}
At Alexandria, the passion of the Saints Caerealis, Pupulus, Caius, and 
 Serapion.
\switchcolumn*
\selectlanguage{latin}
Ibídem commemorátio sanctórum Presbyterórum, Diaconórum et aliórum 
 plurimórum; qui, témpore Valeriáni Imperatóris, cum pestis sævíssima 
 grassarétur, morbo laborántibus ministrántes, libentíssime mortem oppetiére, 
 et quos velut Mártyres religiósa piórum fides venerári consuévit.
\switchcolumn
\selectlanguage{english}
In the same city, in the reign of Emperor Valerian, the commemoration of the 
 holy priests, deacons, and many others. When a most deadly epidemic 
 was raging, they willingly met their death by ministering to the sick. 
 The religious sentiment of the pious faithful has generally venerated them 
 as martyrs.
\switchcolumn*
\selectlanguage{latin}
Romæ sancti Hílari, Papæ et Confessóris.
\switchcolumn
\selectlanguage{english}
At Rome, St. Hilary, pope and confessor.
\switchcolumn*
\selectlanguage{latin}
In território Lugdunénsi, locis Jurénsibus, deposítio sancti Románi Abbátis, 
 qui primum illic eremíticam vitam duxit, et, multis virtútibus ac miráculis 
 clarus, plurimórum póstea Pater éxstitit Monachórum.
\switchcolumn
\selectlanguage{english}
In the territory of Lyons, in the Jura Mountains, the death of St. Romanus, 
 abbot, who first had led the life of a hermit there. His reputation 
 for virtues and miracles brought under his guidance many monks.
\switchcolumn*
\selectlanguage{latin}
Papíæ Translátio córporis sancti Augustíni 
 Epíscopi, Confessóris et Ecclésiæ Doctóris, ex Sardínia ínsula, ópera 
 Luitprándi, Regis Longobardórum.
\switchcolumn
\selectlanguage{english}
At Papia, the transfer, ordered by the Lombard King Luitprand, of the body 
 of St. Augustine, bishop, away from the island of Sardinia.
\switchcolumn*
\selectlanguage{latin}
\end{paracol}

